% Santa Clara University PhD Thesis
% Document class: report, 12pt, letter paper, single-sided
\documentclass[12pt, letterpaper, oneside]{report}

% -----------------------------------------------------------------------
% Packages
% -----------------------------------------------------------------------
\usepackage[left=1.5in, right=1in, top=1in, bottom=1in, headheight=15pt]{geometry}
\usepackage{setspace}
\usepackage{amsmath, amsfonts, amssymb}
\usepackage{amsthm}
\usepackage{graphicx}
\usepackage{array}
\usepackage{caption}
\usepackage{algorithm}
\usepackage{algorithmic}
\usepackage{fancyhdr}
\usepackage{tikz}
\usepackage{booktabs}
\usepackage[hidelinks]{hyperref}

% -----------------------------------------------------------------------
% Theorem environments -- ONE running counter for the whole thesis, tied
% to chapter number. This resolves the Chapter 5 same-numbered "Lemma 1"
% collision (estimation-section noise gains vs. stability-section transverse
% contraction) automatically: under this setup they become, e.g., Lemma 5.3
% and Lemma 5.7 rather than both being "Lemma 1". See Thesis_Execution_Plan.md
% Section 3 (Chapter 5 mapping) and Section 6 item 10.
% -----------------------------------------------------------------------
\newtheorem{theorem}{Theorem}[chapter]
\newtheorem{lemma}[theorem]{Lemma}
\newtheorem{proposition}[theorem]{Proposition}
\newtheorem{corollary}[theorem]{Corollary}
\newtheorem{remark}[theorem]{Remark}
\newtheorem{assumption}[theorem]{Assumption}

% -----------------------------------------------------------------------
% Thesis metadata -- edit these fields
% -----------------------------------------------------------------------
\newcommand{\thesistitle}{Distributed Multirobot Estimation and Navigation in Vector Fields: From Hardware Testbed to Second-Order Field Structure}
\newcommand{\thesisauthor}{Christopher J. Waight}
\newcommand{\thesisdepartment}{Mechanical Engineering}
\newcommand{\thesisdegreefield}{Mechanical Engineering}
\newcommand{\thesisadvisor}{Christopher A. Kitts}
\newcommand{\thesisreaderA}{Reader Name A}
\newcommand{\thesisreaderB}{Reader Name B}
\newcommand{\thesisreaderC}{Reader Name C}
\newcommand{\thesisreaderD}{Reader Name D}
\newcommand{\thesisdeptchair}{Department Chair Name}
\newcommand{\thesisyear}{2026}
\newcommand{\thesismonth}{Month}
\newcommand{\thesisday}{Day}

% -----------------------------------------------------------------------
% Page style setup
% -----------------------------------------------------------------------
\pagestyle{fancy}
\fancyhf{}
\fancyhead[R]{\thepage}
\renewcommand{\headrulewidth}{0pt}

\fancypagestyle{plain}{%
  \fancyhf{}
  \fancyhead[R]{\thepage}
  \renewcommand{\headrulewidth}{0pt}
}

% -----------------------------------------------------------------------
% TOC / LOF / LOT formatting
% -----------------------------------------------------------------------
\renewcommand{\contentsname}{Table of Contents}
\renewcommand{\listfigurename}{List of Figures}
\renewcommand{\listtablename}{List of Tables}

% -----------------------------------------------------------------------
% Hyphenation hints
% -----------------------------------------------------------------------
\hyphenation{multi-robot cluster-space Jacobian separ-atrix}

% -----------------------------------------------------------------------
% Document
% -----------------------------------------------------------------------
\begin{document}

% Front matter uses roman numerals
\pagenumbering{roman}
\setcounter{page}{1}

% ===================================================================
% 1. APPROVAL / SIGNATURE PAGE
% ===================================================================
\begin{titlepage}
\begin{center}
  \large
  Santa Clara University \\[6pt]
  Department of \thesisdepartment

  \vspace{1.5in}

  \normalsize
  Date: \thesismonth\ \thesisday, \thesisyear

  \vspace{1in}

  I HEREBY RECOMMEND THAT THE THESIS PREPARED UNDER MY SUPERVISION BY

  \vspace{0.5in}
  \MakeUppercase{\thesisauthor}
  \vspace{0.5in}

  ENTITLED

  \vspace{0.4in}
  \textit{\thesistitle}
  \vspace{0.4in}

  BE ACCEPTED IN PARTIAL FULFILLMENT OF THE REQUIREMENTS FOR THE DEGREE

  \vspace{0.3in}
  \textbf{DOCTOR OF PHILOSOPHY IN \MakeUppercase{\thesisdegreefield}}
\end{center}

\vspace{0.8in}

\noindent
\begin{tabular}{p{3in} p{0.3in} p{3in}}
  \hrulefill & & \hrulefill \\[2pt]
  Thesis Advisor & & Thesis Reader \\[24pt]
  \hrulefill & & \hrulefill \\[2pt]
  Chairman of Department & & Thesis Reader \\[24pt]
  & & \hrulefill \\[2pt]
  & & Thesis Reader \\[24pt]
  & & \hrulefill \\[2pt]
  & & Thesis Reader
\end{tabular}

\end{titlepage}

% ===================================================================
% 2. TITLE PAGE
% ===================================================================
\begin{titlepage}
\begin{center}
  \vspace*{1in}

  {\Large \thesistitle}

  \vspace{0.6in}
  by \\[6pt]
  \thesisauthor

  \vspace{1.2in}
  Dissertation \\[6pt]
  Submitted in Partial Fulfillment of the Requirements \\
  for the Degree of Doctor of Philosophy \\
  in \thesisdegreefield \\[6pt]
  in the School of Engineering at \\
  Santa Clara University, \thesisyear

  \vspace{1in}
  Santa Clara, California \\
  \thesisyear
\end{center}
\end{titlepage}

% ===================================================================
% 3. COPYRIGHT
% ===================================================================
\clearpage
\thispagestyle{plain}
\vspace*{3in}
\begin{center}
  \copyright\ Copyright by \thesisauthor\ \thesisyear \\
  All Rights Reserved
\end{center}
\clearpage

% ===================================================================
% 4. ACKNOWLEDGEMENTS
% ===================================================================
\chapter*{Acknowledgements}
\addcontentsline{toc}{chapter}{Acknowledgements}

% %% TODO: SYNC WITH LATEST DRAFT -- placeholder acknowledgements. Per
% Thesis_Execution_Plan.md Section 5, item 3: fold in the names from the
% Decabot testbed paper's own Acknowledgements (Scot Tomer, Jiayi Wang,
% Michael Waight, Sara Alvarez) rather than keeping a separate per-chapter
% block in Chapter 2.

I would like to thank my advisor, Christopher A. Kitts, and the members of
the Santa Clara University Robotic Systems Laboratory for their support
throughout this work.

A portion of this work has been published or submitted in [1] and [2], and
portions are in preparation for future publication.

\clearpage

% ===================================================================
% 5. ABSTRACT
% ===================================================================
\chapter*{Abstract}
\addcontentsline{toc}{chapter}{Abstract}

\begin{center}
  \thesistitle \\[6pt]
  \thesisauthor \\[6pt]
  Department of \thesisdepartment \\
  Santa Clara University \\
  Santa Clara, California \\
  \thesisyear
\end{center}

\vspace{0.4in}

% %% TODO: SYNC WITH LATEST DRAFT -- this abstract is drafted from the
% Thesis_Execution_Plan.md Section 1 narrative and contributions list; it
% has not been reviewed by the user and should be treated as a first pass.

\noindent
This dissertation develops a progression of tools for multirobot estimation
and navigation in physical fields, moving from the hardware substrate upward
to increasingly higher-order field information. A functional hardware testbed
of omnidirectional rovers, printed HSV-encoded vector field maps, and
neural-network sensor calibration establishes the experimental platform used
throughout. A compositional grammar for cluster-space kinematics then
generalizes the geometric machinery underneath every multirobot formation:
formations are trees built from a small library of base charts, and the
resulting inverse Jacobian is complete for any tree and assembles in a single
pass, replacing one-off, per-formation derivations. Building on this
substrate, a three-robot team is shown to estimate a two-dimensional vector
field's local Jacobian from simultaneous local measurements alone, with no
prior knowledge of the field, enabling closed-form critical point
localization, eigenvalue-based classification, and attraction and orbital
control laws; three robots is proven both necessary and sufficient for this
first-order estimate, and the method is validated in simulation across eight
canonical fields and on hardware with a 100\% convergence success rate over
157 trials. This first-order result is then extended to second order: a
six-robot team recovers a flow's local Jacobian, Okubo-Weiss parameter, and
Hessian from simultaneous measurements, enabling two modified-Newton
controllers that track dynamic boundaries and separatrices no single critical
point can describe, with a proven stability guarantee and validation on both
a canonical double-gyre flow and real oceanographic current data. Together,
these contributions establish a coherent methodology for distributed,
model-free field estimation that scales from isolated critical points to
extended coherent structures using only local, cooperative sensing.

\clearpage

% ===================================================================
% 6. GLOSSARY OF TERMS
% ===================================================================
% Modeled on the Glossary of Terms front-matter chapter in
% ShaeThesisForexample of template.pdf (Shae Hart's 2023 SCU ME PhD thesis),
% which is not present in the blank phd_thesis_template.tex and has been
% added here per Thesis_Execution_Plan.md Section 5, Prompt 1. This is a
% single alphabetized symbol table, distinct from the discursive notation
% discussion inside Chapter 1's "Notation and Definitions" section: this
% table is for quick lookup while reading later chapters, the Chapter 1
% section is for understanding the collisions and the reasoning behind the
% chosen conventions the first time.
%
% %% TODO: SYNC WITH LATEST DRAFT -- populate this table's page-number-style
% quick reference once all chapter numbers/labels are final; for now it
% lists symbol and meaning only, matching Shae Hart's own glossary format.
\chapter*{Glossary of Terms}
\addcontentsline{toc}{chapter}{Glossary of Terms}

\begin{table}[h]
\centering
\begin{tabular}{l p{3.5in}}
$\mathbf{p}$ & Position vector in the plane, $\mathbf{p}=(x,y)$ \\
$\mathbf{p}^*$ & Critical point of a vector field, $\mathbf{v}(\mathbf{p}^*)=\mathbf{0}$ \\
$\mathbf{v}(\mathbf{p})$ & Vector field value at $\mathbf{p}$, components $(u,v)$ \\
$\mathbf{J}$ & Field Jacobian, $\mathbf{J} = \left[\begin{smallmatrix} u_x & u_y \\ v_x & v_y \end{smallmatrix}\right]$ \\
$\mathbf{J}_c$ & Kinematic (cluster-space) Jacobian, robot-space to cluster-space velocity map \\
$\alpha_{\text{eig}}$ & Real part of a complex eigenvalue pair of $\mathbf{J}$ (spiral classification) \\
$\alpha_{\text{mom}}$ & Momentum / discretization coefficient, $\alpha_{\text{mom}} = \exp(-\Delta t/\tau)$ \\
$\phi_1, \phi_3$ & SAS triangle vertex angles at robots 1 and 3 (Appendix, Ch.~3) \\
$\alpha_{\text{spine}}$ & Spine-relative shape-angle offset, composition-tree extension (Ch.~3) \\
$p, q, \beta$ & SAS triangle side lengths and interior angle (inherited Kitts-lab notation, unchanged across chapters) \\
$\mathbf{q}$ (bold) & Cluster-space state vector (pose + shape, dimension $2N$) \\
$q$ (scalar) & SAS triangle's second side length (2--3 distance); never confused with bold $\mathbf{q}$ \\
$r$ & Orbital radial vector / magnitude, $\mathbf{r} = \mathbf{p}^*-\mathbf{p}_c$ \\
$s$ & SAS triangle's third side length (1--3 diagonal) \\
$\rho$ & Polar radius (field definitions, Ch.~4 appendix) and pentagon ring radius (Ch.~5) \\
$D$ & Okubo-Weiss parameter, $D = \det(\mathbf{J})$ \\
$\mathbf{H}_D$ & Hessian of $D$ \\
$\boldsymbol{\phi}(\tilde{\mathbf{p}})$ & Quadratic sensing basis vector (Ch.~5) \\
$\boldsymbol{\Phi}$ & Six-robot formation matrix (capital Phi; Ch.~5) \\
$\tau$ & Actuator time constant, 0.3~s on Decabot hardware \\
$\Delta t$ & Control period, 0.1~s (10~Hz) \\
$v_{\max}$ & Maximum robot speed, 0.3~m/s \\
\end{tabular}
\end{table}

\clearpage

% ===================================================================
% 7. TABLE OF CONTENTS
% ===================================================================
\tableofcontents
\clearpage

% ===================================================================
% 8. LIST OF FIGURES
% ===================================================================
\listoffigures
\addcontentsline{toc}{chapter}{List of Figures}
\clearpage

% ===================================================================
% 9. LIST OF TABLES
% ===================================================================
\listoftables
\addcontentsline{toc}{chapter}{List of Tables}
\clearpage

% ===================================================================
% BODY -- switch to arabic page numbering
% ===================================================================
\pagenumbering{arabic}
\setcounter{page}{1}
\doublespacing

% ===================================================================
% CHAPTER 1: INTRODUCTION
% ===================================================================
\chapter{Introduction}
\label{ch:intro}

\section{Motivation}

Many environments of scientific and practical interest are organized by an
underlying vector field: an ocean current, an atmospheric flow, a diffusing
plume, a magnetic field. Locating and characterizing the special points and
structures of such a field, its sinks, sources, vortices, saddles, and the
separatrices that divide one flow regime from another, is often more useful
than simply following the field's local direction. A search-and-rescue team
drifting with an ocean current benefits from knowing where a nearby eddy will
trap debris; an environmental monitoring team benefits from knowing where two
water masses meet. In many of these settings no map of the field exists in
advance, and no single sensor can see the whole field at once. A team of
mobile robots, each sampling the field locally and sharing measurements with
its neighbors, can estimate field structure that no individual robot could
determine alone, and can then use that estimate to steer.

This dissertation develops such a capability in stages. It begins with the
physical substrate every later result depends on: a working hardware testbed
of small robots that can sense a printed, encoded vector field and be
tracked precisely enough to validate a control law, not just simulate it. It
continues with the geometric machinery every multirobot formation needs
regardless of team size: a systematic way to relate the motion of a formation,
as a rigid or semi-rigid shape, to the motion of its individual robots. It
then asks what a minimal team, three robots, can infer about the field
surrounding it using only first-order (Jacobian) information, and shows this
is enough to locate and classify isolated critical points and navigate to or
around them. Finally it asks what a larger team, six robots, can infer using
second-order (Hessian) information, and shows this is enough to track
extended, moving structures, such as a separatrix or an Okubo-Weiss boundary,
that no isolated critical point can describe.

\section{Problem Statement}

Existing approaches to multirobot navigation in scalar and vector fields
generally fall into two categories. The first requires prior knowledge of
the field, a map, a model, or an assumed functional form, against which a
robot team's local measurements are compared or fit. The second uses only
local information but provides directional guidance alone: climb the
gradient, follow the current, avoid the strain. Neither category directly
answers the question this dissertation poses: given only local, real-time
measurements and no prior model, can a robot team estimate the location and
type of a field's critical points, and, more generally, the field's local
first- and second-order structure, well enough to navigate with respect to
it, not just along it.

This dissertation answers that question constructively. It shows that three
robots are both necessary and sufficient to recover a two-dimensional vector
field's local linearization (its Jacobian) from simultaneous local
measurements alone, and that six robots are both necessary and (almost
always) sufficient to recover the same field's local quadratic structure
(its Jacobian, Okubo-Weiss parameter, and Hessian). In both cases the
estimate is exact for the class of fields the robots' formation size implies
(affine for three robots, quadratic for six), assembled without any
assumption about the field's global shape, and used directly inside a
navigation control law whose stability is established analytically and
validated in simulation and, where hardware permits, on a physical testbed.

\section{Contributions}

The main contributions of this dissertation are:

\begin{enumerate}
\item A functional, validated hardware testbed for multirobot vector-field
  navigation using HSV-encoded printed fields and neural-network sensor
  calibration (R\textsuperscript{2}~=~0.96 direction, 0.91 magnitude),
  removing the RGB channel-interference failure of the prior testbed
  generation (Chapter~\ref{ch:testbed}).
\item A compositional grammar for cluster-space kinematics: any multirobot
  formation is a tree of three base charts (pair, SAS triangle, $k$-body
  star); the resulting inverse Jacobian is complete and square for any tree
  and assembles in a single depth-first pass, $O(N)$ at bounded hierarchy
  depth and $O(N\log N)$ for balanced trees, replacing a prior method's
  $O(\text{depth}^2)$ rotation-product recomputation. A consensus
  orientation gauge additionally removes a single-pointer-robot fragility
  near field critical points (Chapter~\ref{ch:cluster_builder}).
\item A distributed, model-free method for three robots to estimate a
  two-dimensional vector field's local Jacobian from simultaneous local
  measurements alone, with no prior field knowledge, proven minimal at
  exactly three robots, enabling closed-form critical point localization,
  eigenvalue-based type classification, and attraction and orbital control
  laws; validated in simulation (eight fields, 1000 Monte Carlo trials each,
  100\% convergence) and on hardware (157 convergence trials plus 12 orbital
  trials, 100\% success, sub-centimeter bias) (Chapter~\ref{ch:vector_field}).
\item A distributed, model-free method for six robots (pentagon plus center)
  to estimate a two-dimensional flow's local Jacobian, Okubo-Weiss parameter,
  and Hessian from simultaneous local measurements, proven minimal at exactly
  six robots and almost-always sufficient, ruling out the degenerate hexagon
  and conic trap, enabling two modified-Newton controllers that track the
  Okubo-Weiss boundary and the separatrix trench respectively, with a proven
  stability and traversal guarantee and validation on both a canonical
  double-gyre model and real ocean high-frequency-radar current data
  (Chapter~\ref{ch:separatrix}).
\end{enumerate}

\section{Publications}

Portions of this dissertation have appeared, or are in preparation to appear,
in the following publications.

\begin{table}[h]
\centering
\caption{Mapping of dissertation chapters to publications.}
\begin{tabular}{p{0.6in} p{2.6in} p{1in} p{0.8in}}
\toprule
Chapter & Working title & Venue & Status \\
\midrule
\ref{ch:testbed} & A Functional Indoor Testbed for Multirobot Adaptive
  Navigation in Vector Field Environments & IDETC/CIE 2025 & Submitted \\
\ref{ch:cluster_builder} & A Compositional Grammar for Cluster-Space
  Kinematics with Single-Pass Inverse Jacobian Assembly (evolves the
  author's Master's thesis, An Algorithm for Calculating the Inverse
  Jacobian of Multirobot Systems in a Cluster Space Formulation) &
  IDETC/CIE (target venue TBD) & In preparation \\
\ref{ch:vector_field} & Adaptive Navigation of Multirobot Systems to
  Critical Points in 2D Vector Fields & IEEE/ASME Transactions on
  Mechatronics & Near submission \\
\ref{ch:separatrix} & (Separatrix / Okubo-Weiss tracking title) & IEEE
  Transactions on Robotics & In preparation \\
\bottomrule
\end{tabular}
\end{table}

\section{Organization}

Chapter~\ref{ch:testbed} describes the Decabot hardware testbed: the robot
platform, the HSV-encoded printed field representation, and the
neural-network sensor calibration that together make hardware validation
possible for every later chapter. Chapter~\ref{ch:cluster_builder}
generalizes the cluster-space kinematics underlying every formation used in
this dissertation, presenting a compositional grammar that assembles a
formation's kinematic Jacobian from a small library of base charts rather
than deriving each formation by hand. Chapter~\ref{ch:vector_field} presents
the first-order estimation result: three robots recovering a vector field's
local Jacobian, and the resulting critical point localization, classification,
and navigation control laws, validated in simulation and on hardware.
Chapter~\ref{ch:separatrix} presents the second-order extension: six robots
recovering a flow's local Hessian and Okubo-Weiss structure, and the
resulting boundary- and separatrix-tracking controllers, validated in
simulation on a canonical double-gyre model and on real ocean current data.
Chapter~\ref{ch:conclusions} synthesizes the four contributions and
discusses future work. Appendix~\ref{app:four_robot} presents supplementary,
unpublished four-robot hardware trials that generalize
Chapter~\ref{ch:vector_field}'s three-robot result.

\section{Notation and Definitions}
\label{sec:notation}

Throughout this dissertation, bold lowercase letters denote vectors, e.g.,
$\mathbf{p} \in \mathbb{R}^2$, and bold uppercase letters denote matrices,
e.g., $\mathbf{J} \in \mathbb{R}^{m \times n}$. The four studies collected
in this dissertation were originally written as independent papers, and
their notations were never cross-checked against one another. Most symbols
already agree; a small number collide and are resolved here, once, for the
whole document. Chapter-specific notation not covered here is introduced
where it is first used.

\subsection*{The bold-\texorpdfstring{$\mathbf{p}$}{p} versus scalar-\texorpdfstring{$p$}{p} rule}

Bold $\mathbf{p}$ always denotes a two-dimensional position,
$\mathbf{p} = (x,y)$, and is never written unbolded when used as a position.
Scalar $p$ always denotes the side-length formation parameter inherited from
the Kitts-lab cluster-space lineage (Chapter~\ref{ch:cluster_builder} onward):
the distance from robot 1 to robot 2 in a side-angle-side (SAS) triangular
formation. There is no unbolded $p$ for position and no bold $\mathbf{p}$ for
side length. This rule is carried unchanged from the source papers and holds
throughout every chapter of this dissertation.

\subsection*{A new rule: the bold-\texorpdfstring{$\mathbf{q}$}{q} versus scalar-\texorpdfstring{$q$}{q} rule}

Chapter~\ref{ch:cluster_builder} introduces bold $\mathbf{q}$ as the general
cluster-space state vector (pose and shape stacked, dimension $2N$ for $N$
robots), while scalar $q$, following the same inherited SAS convention as
scalar $p$, denotes the triangle's second side length (the 2--3 distance).
These two uses never appear in the same equation and are distinguished
exactly as bold $\mathbf{p}$ and scalar $p$ are: there is no unbolded
$\mathbf{q}$ for the state vector and no bold $q$ for the side length. This
rule did not exist in any of the four source papers individually; it is
introduced here because this dissertation is the first place the general
cluster-space state vector and the SAS side length appear together.

\subsection*{The alpha collision (four-way)}

The symbol $\alpha$ is used for four unrelated quantities across the source
materials. Each is disambiguated with a subscript in this dissertation and
must never be written as bare $\alpha$ after this point.

\begin{itemize}
\item $\alpha_{\text{eig}}$: the real part of a complex eigenvalue pair of
  the field Jacobian $\mathbf{J}$, written $\alpha_{\text{eig}} \pm i\omega$.
  This is a property of the vector field's linearization at a critical point
  and classifies spiral-type critical points (Chapter~\ref{ch:vector_field}).
\item $\alpha_{\text{mom}}$: the discrete first-order momentum filter
  coefficient for robot dynamics, $\alpha_{\text{mom}} = \exp(-\Delta t/\tau)$.
  At 10~Hz with $\tau = 0.3$~s, $\alpha_{\text{mom}} \approx 0.717$. This
  appears identically in the robot dynamics of Chapters~\ref{ch:vector_field}
  and~\ref{ch:separatrix}.
\item $\phi_1, \phi_3$: the two SAS-triangle vertex angles at robots 1 and 3,
  used only in the triangle's law-of-cosines derivation
  (Chapter~\ref{ch:cluster_builder} appendix material). These were originally
  written as $\alpha$ and $\gamma$ in one source paper's appendix; they are
  renamed here since they are otherwise unused and the original choice
  collided with both $\alpha_{\text{eig}}$ and $\alpha_{\text{mom}}$.
\item $\alpha_{\text{spine}}$: a spine-relative shape-angle offset,
  $\alpha_{\text{spine},i} = \theta_{c_i} - \theta_v$, appearing only in the
  optional composition-tree extension of Chapter~\ref{ch:cluster_builder}.
  This is unrelated to all three uses above.
\end{itemize}

\subsection*{The \texorpdfstring{$r$}{r} collision}

Three meanings of $r$ appear across the source materials: the orbital radial
vector/magnitude $\mathbf{r} = \mathbf{p}^* - \mathbf{p}_c$
(Chapter~\ref{ch:vector_field}), the SAS triangle's third side length (the
1--3 diagonal), and the polar radius used in several canonical field
definitions. This dissertation keeps $r$ for the orbital radial quantity,
since it is the most load-bearing use and appears in a named control law.
The triangle's third side is renamed $s$, and the polar radius in field
definitions is renamed $\rho$, matching the pentagon ring radius already used
in Chapter~\ref{ch:separatrix} so that ``radius of a circular construction''
reads consistently throughout.

\subsection*{The \texorpdfstring{$\mathbf{J}$}{J} collision}

Bold $\mathbf{J}$ denotes the field Jacobian, $\mathbf{J} =
\left[\begin{smallmatrix} u_x & u_y \\ v_x & v_y \end{smallmatrix}\right]$,
throughout this dissertation: it is the central mathematical object of both
Chapter~\ref{ch:vector_field} (first order) and Chapter~\ref{ch:separatrix}
(embedded in the second-order estimator). Every kinematic or cluster-space
Jacobian, mapping shape-parameter velocities to individual robot velocities,
is written $\mathbf{J}_c$ instead, whether it is the $6\times 6$ triangle
Jacobian or the $12\times 12$ pentagon Jacobian, both relocated into
Chapter~\ref{ch:cluster_builder}.

\subsection*{Other symbols}

The full symbol reference for field and estimation notation, critical point
classification, control laws, robot dynamics, and formation parameters
follows \texttt{docs/notation.md} unchanged, extended with the
Chapter~\ref{ch:cluster_builder} and Chapter~\ref{ch:separatrix} additions
listed in the Glossary of Terms. Of particular note: capital $\boldsymbol{\Phi}$
(the six-robot formation matrix, Chapter~\ref{ch:separatrix}) is visually
close to lowercase $\phi_1, \phi_3$ (the triangle vertex angles,
Chapter~\ref{ch:cluster_builder}) in some fonts; the two never appear in the
same chapter, but the reader should not assume a relationship between them.
Similarly, $\rho$ appears in three guises across this dissertation, the field
definitions' polar radius, the pentagon formation's ring radius, and the
consensus orientation gauge's alignment magnitude, $\rho_{\text{gauge}} =
\sqrt{C^2+S^2}$ (Chapter~\ref{ch:cluster_builder}); the gauge quantity is
always written with the \texttt{gauge} subscript to keep it distinct from the
other two.

As in Chapter~\ref{ch:vector_field}, we write a critical point of a flow,
including a saddle point of the Okubo-Weiss field in
Chapter~\ref{ch:separatrix}, as $\mathbf{p}^*$.

% ===================================================================
% CHAPTER 2: DECABOT HARDWARE TESTBED
% ===================================================================
\chapter{The Decabot Hardware Testbed}
\label{ch:testbed}

\section{Introduction}

Investigating multirobot adaptive navigation in a vector field requires a
field to navigate: some way to present a robot team with a physically real,
locally sensed vector field rich enough to contain the critical points and
structures the rest of this dissertation studies. This chapter presents the
Decabot testbed, a hardware platform built for exactly this purpose. A prior
generation of this testbed encoded vector field readings directly in RGB
color and printed the result on the floor as a map for the robots to sense.
That approach suffered a persistent limitation: interference between color
channels in the robots' undercarriage-mounted sensors degraded the system's
ability to recover accurate direction and magnitude information, and the
discrete color palette limited the effective sampling rate to roughly
0.15~Hz, too slow to support the continuous, closed-loop control laws
developed in Chapters~\ref{ch:vector_field} and~\ref{ch:separatrix}.

This chapter describes a redesigned encoding, direction in hue and magnitude
in saturation, together with a neural-network calibration procedure that
corrects for the channel interference the prior approach could not. The
result is a testbed capable of continuous robot motion at full control rate,
validated here against two navigation primitives originally verified only in
simulation. The chapter's contributions are: (1) an HSV-based encoding for
representing two-dimensional vector fields on a printed floor map, (2) a
neural-network calibration technique that achieves encoding-decoding accuracy
of R\textsuperscript{2}~=~0.96 in direction and 0.91 in magnitude, (3)
hardware validation of the vector-sum and vector-to-scalar adaptive
navigation control primitives, previously verified only in simulation, and
(4) a general method for building this class of testbed that later chapters'
hardware validation depends on directly.

\section{Control Architecture Overview}
\label{sec:testbed_architecture}

Every hardware and simulation result in this dissertation uses the same
three-layer control architecture, introduced here and reused without
modification in Chapters~\ref{ch:vector_field} and~\ref{ch:separatrix}.

\begin{enumerate}
\item \textbf{Robot controller layer.} Each Decabot is an omnidirectional
  planar rover carrying an RGB color sensor (a tcs34725 undercarriage sensor
  whose channels are not fully isolated per its datasheet, the source of the
  interference problem this chapter addresses), an Arduino-class onboard
  controller, and a WiFi/TCP-IP link to a central computer. Robots accept a
  commanded velocity and execute it subject to a maximum linear velocity of
  0.3~m/s and a stiction floor below which commanded motion does not produce
  actual motion. The discrete-time momentum model governing this layer is
  developed in full in Chapter~\ref{ch:vector_field}, Section~3.1, and reused
  unchanged in Chapter~\ref{ch:separatrix}.
\item \textbf{Cluster-space controller layer.} A formation's shape and pose
  are described by a small set of cluster-space variables (for a three-robot
  triangular formation, two side lengths and an interior angle, plus
  centroid position and heading) rather than by each robot's individual
  pose. A kinematic Jacobian maps between cluster-space velocities and
  individual robot velocities in both directions. Chapter~\ref{ch:cluster_builder}
  generalizes this mapping to arbitrary formation topologies; the specific
  triangular and pentagon instances used in later chapters are presented
  there as well.
\item \textbf{Adaptive navigation layer.} A navigation primitive computes a
  desired cluster-space (typically centroid) velocity from the team's shared
  field measurements, using a feature estimator (the subject of
  Chapters~\ref{ch:vector_field} and~\ref{ch:separatrix}) and a navigation
  controller built on top of it. A formation controller maintains the
  commanded shape independently of navigation, holding the cluster-space
  shape variables at their setpoints, here $p=q=0.35$~m and $\beta=1.05$~rad
  for the equilateral-triangle-like formation used throughout this chapter's
  hardware trials.
\end{enumerate}

The full loop, sensing, estimation, navigation, cluster-space inversion, and
robot-level actuation, runs at 10~Hz over the WiFi link described below.

\section{Field Representation and Sensor Calibration}
\label{sec:testbed_fields}

\subsection{Testbed Overview}

The physical testbed consists of three Decabot rovers operating on a
multicolored printed floor map (2.3~m\textsuperscript{2} per map), an
OptiTrack motion capture system providing centimeter-level position accuracy
via reflective markers on each robot's roof, and a central computer running
the control architecture of Section~\ref{sec:testbed_architecture} in
MATLAB/Simulink. Robots communicate with the central computer over WiFi;
workspace size is bounded above by approximately 3~m~$\times$~6~m and below
by roughly 1~m~$\times$~1~m (a collision-clearance minimum for a three-robot
team); battery life limits a continuous trial to approximately one hour.

\subsection{Modeling Vector Fields in Color Space}

A vector field reading has two components at each point: a direction and a
magnitude. This testbed encodes direction in hue and magnitude in
saturation, allowing a two-dimensional vector field to be represented as a
printed image and sensed by a simple color sensor rather than requiring
active field generation. Two test fields are used throughout this chapter's
hardware trials: a fixed vortex and a sinking vortex, both centered on the
printed map and defined analytically so that ground truth is known
everywhere on the map.

\subsection{Feature Estimation and Sensor Calibration}

The undercarriage color sensor's raw channels, red, green, blue, and a clear
channel, do not map cleanly to hue and saturation because the channels are
not fully isolated from one another; this is precisely the interference
problem that limited the prior testbed generation. This chapter's solution
is a pair of neural networks, trained per robot (or, for a generalization
test, trained on three robots and validated on a fourth held out), that map
the sensor's raw seven-dimensional input (the four raw channels plus the
naively computed hue, saturation, and value) to corrected hue and saturation
estimates. Hue is decomposed into sine and cosine components before training
to respect its circular topology; saturation is regressed directly. Networks
use one or two hidden layers (searched over 2--12 units in the first layer,
4--10 in the second), Tanh activations, and the Levenberg-Marquardt optimizer,
with a 70/15/15 train/validation/test split.

Across eight validation datasets, the cross-robot generalized model achieves
hue R\textsuperscript{2}~$\geq$~0.96 (circular RMSE~$\leq$~0.05) and
saturation R\textsuperscript{2}~$\geq$~0.93 (RMSE~$\leq$~0.079); the held-out
fourth robot alone achieves hue R\textsuperscript{2}~=~0.993 and saturation
R\textsuperscript{2}~=~0.908 (Table~\ref{tab:testbed_calibration}).
Individually tuned per-robot models improve further, with hue
R\textsuperscript{2}~$>$~0.99 and saturation R\textsuperscript{2}~$>$~0.95
across all three robots (Table~\ref{tab:testbed_individual}).

\begin{table}[h]
\centering
\caption{Performance of the cross-robot generalized calibration model on
eight validation datasets, and on a fourth robot held out entirely from
training.}
\label{tab:testbed_calibration}
\begin{tabular}{lcc}
\toprule
 & Hue & Saturation \\
\midrule
Cross-robot model (all validation sets) & $R^2 \geq 0.96$ & $R^2 \geq 0.93$ \\
Held-out robot (generalization test) & $R^2 = 0.993$ & $R^2 = 0.908$ \\
\bottomrule
\end{tabular}
\end{table}

\begin{table}[h]
\centering
\caption{Model inference performance for individually tuned per-robot
calibration models.}
\label{tab:testbed_individual}
\begin{tabular}{lcccc}
\toprule
Robot & Hue $R^2$ & Hue RMSE & Sat.\ $R^2$ & Sat.\ RMSE \\
\midrule
1 & 0.995 & 0.020 & 0.963 & 0.055 \\
2 & 0.991 & 0.027 & 0.952 & 0.063 \\
3 & 0.997 & 0.015 & 0.976 & 0.044 \\
\bottomrule
\end{tabular}
\end{table}

\subsection{Robot Velocity from Saturation}

Because saturation encodes field magnitude, and a robot's commanded velocity
in the navigation primitives below scales with sensed magnitude, this
testbed also validates that robot velocity responds linearly to sensed
saturation above the stiction floor (saturation values below approximately
0.2 were excluded from calibration for this reason).

\section{Control Primitive Validation}
\label{sec:testbed_primitives}

Two navigation primitives, previously verified only in simulation, are
validated on hardware here.

\begin{itemize}
\item \textbf{Vector-sum:} each robot's sensed vector reading is summed
  across the team and normalized to set a commanded direction.
\item \textbf{Vector-to-scalar:} only sensed magnitude is used, converting
  navigation to a scalar-field gradient problem; a minimum-seeking variant
  is used to locate a vortex center.
\end{itemize}

On the fixed vortex, the vector-sum primitive produces a consistent outward
drift across repeated trials; on the sinking vortex, the same primitive
settles the cluster to within approximately 0.1~m of the true center. The
vector-to-scalar (minimum-seeking) primitive on the fixed vortex settles the
cluster within approximately 0.2~m of center, with performance essentially
unchanged when starting positions are randomized rather than fixed. Neither
orbiting nor saddle-point navigation, the subjects of
Chapter~\ref{ch:vector_field}, was hardware-tested at this stage; both are
noted in this testbed's own future work as requiring the Jacobian-based
estimation developed in the next two chapters.

Formation shape was held within approximately 1~cm of its commanded values
throughout all trials (Table~\ref{tab:testbed_formation}), consistent with
the resolution limits of the OptiTrack system itself.

\begin{table}[h]
\centering
\caption{Cluster formation-holding accuracy across control primitive trials,
commanded $p=q=0.35$~m, $\beta=1.05$~rad.}
\label{tab:testbed_formation}
\begin{tabular}{lccc}
\toprule
 & $p$ (m) & $q$ (m) & $\beta$ (rad) \\
\midrule
Fixed vortex, vector-sum & $0.36 \pm 0.028$ & $0.36 \pm 0.032$ & $1.03 \pm 0.122$ \\
Fixed vortex, vector-to-scalar & $0.35 \pm 0.015$ & $0.36 \pm 0.009$ & $1.02 \pm 0.045$ \\
Sinking vortex, vector-sum & $0.35 \pm 0.019$ & $0.35 \pm 0.015$ & $1.02 \pm 0.056$ \\
\bottomrule
\end{tabular}
\end{table}

\section{Limitations and Summary}
\label{sec:testbed_limitations}

This testbed's own limitations are primarily physical rather than
algorithmic, and remain in force for every hardware result in later
chapters: the HSV encoding is inherently two-dimensional and does not extend
to a three-dimensional field; workspace size is bounded between roughly
1~m~$\times$~1~m and 3~m~$\times$~6~m; battery life limits a continuous
trial to about one hour; and the printed colormap itself wears over
repeated use and requires periodic replacement. Algorithmic limitations,
team size, static versus dynamic fields, and the estimation order available
to a given formation size, are the proper subject of
Chapters~\ref{ch:vector_field} and~\ref{ch:separatrix} and are not repeated
here.

This chapter has presented a hardware testbed that resolves the RGB
channel-interference limitation of a prior design through HSV encoding and
neural-network sensor calibration, achieving direction and magnitude
recovery accurate enough to validate closed-loop navigation primitives at
full control rate. The remaining chapters of this dissertation build on this
testbed directly: Chapter~\ref{ch:cluster_builder} generalizes the
cluster-space layer introduced in Section~\ref{sec:testbed_architecture}
beyond the three-robot triangle used here, and
Chapters~\ref{ch:vector_field} and~\ref{ch:separatrix} both validate their
own estimation and control results, in part, on this same physical
platform.

% ===================================================================
% CHAPTER 3: CLUSTER BUILDER
% ===================================================================
\chapter{A Compositional Grammar for Cluster-Space Kinematics}
\label{ch:cluster_builder}

\section{Introduction}

Cluster-space control coordinates a group of mobile robots through a small
set of formation-level pose and shape variables rather than through each
robot's individual position. Chapter~\ref{ch:testbed} already relied on this
idea for a three-robot triangular formation, and
Chapters~\ref{ch:vector_field} and~\ref{ch:separatrix} rely on it again for
formations of three and six robots respectively. Each of those instances,
historically, has been derived by hand: a formation's forward and inverse
kinematic maps, and the Jacobian relating cluster-space velocities to
individual robot velocities, are worked out from scratch for each new
formation topology.

This chapter is a direct evolution of the author's Master's thesis work on
this same problem, summarized in Section~\ref{sec:cb_prior_work} below. That
work automated the derivation itself, but its automation still recomputed a
full rotation product from every sub-cluster to the root at every step, an
$O(\text{depth}^2)$ cost that the thesis identified and left as future work.
The present chapter closes that gap and generalizes the construction well
beyond what the earlier tree-propagation algorithm covered. A formation is
recast as a tree: each node applies one of a small
library of base charts (a two-robot pair, a three-robot side-angle-side
triangle, or a $k$-body star) to an ordered list of children, where a child
may be either an individual robot or another sub-cluster. A direct dimension
count shows the resulting chart is complete and square for any such tree, so
a formation is never sized by hand. The forward and inverse maps at each
node are closed form, and the complete $2N \times 2N$ inverse Jacobian is
assembled in a single depth-first traversal using a per-node sensitivity
matrix, removing the redundant rotation product and lowering assembly cost
to $O(N)$ at bounded hierarchy depth and $O(N \log N)$ for balanced trees.
Finally, a consensus orientation gauge distributes a formation's reference
frame across every spoke of a star-shaped chart in place of a single pointer
robot, trading one rank of decoupling for a frame that survives a wandering
or near-coincident reference robot, precisely the situation a formation is
in when it loiters near a field critical point, exactly the regime
Chapters~\ref{ch:vector_field} and~\ref{ch:separatrix} both study.

This chapter's contributions, each a specific advance over the author's
prior tree-propagation algorithm (Section~\ref{sec:cb_prior_work}), are: a
compositional tree grammar that builds a complete, square cluster-space
chart for any formation from a fixed library of closed-form base charts,
with robots and sub-clusters interchangeable as siblings
(Section~\ref{sec:grammar}); a single depth-first pass that assembles the
exact closed-form inverse Jacobian using per-node sensitivity matrices,
removing the depth-quadratic rotation-product recomputation that the prior
algorithm carried (Section~\ref{sec:assembly}); and a consensus orientation
gauge, with no antecedent in the prior work, whose rank-one structure
decouples the formation frame from any single reference robot, improving
conditioning near formation singularities (Section~\ref{sec:gauge}).

\section{Prior Work: The Cluster Tree Inverse Jacobian}
\label{sec:cb_prior_work}

The starting point for this chapter is the author's Master's thesis [66],
which posed the same underlying problem: computing the inverse Jacobian that maps
cluster-space velocities to individual robot velocities had historically
required deriving and differentiating a formation's inverse kinematics by
hand, a process of $\left(\sum_i p_i\right)\times q$ partial derivatives for
$n$ robots with $p_i$ degrees of freedom each and $q$ cluster variables,
tedious and error-prone at any nontrivial formation size. That thesis
represented a cluster as a tree of coordinate frames connected by
homogeneous transformation matrices, robot frames as leaves, cluster and
sub-cluster frames as internal nodes, the global frame as root, and computed
each row of the inverse Jacobian by propagating linear and angular velocity
recursively along the path from a robot's frame to the root:
\begin{equation}
{}^G\mathbf{v}_n = \sum_{k=1}^m \left[{}^G\mathbf{v}_k + {}^G\boldsymbol{\omega}_k \times {}^G\mathbf{p}_{k,n}\right],
\qquad
{}^G\boldsymbol{\omega}_n = \sum_{k=1}^m {}^G\boldsymbol{\omega}_k,
\end{equation}
for a robot $n$ at hierarchical depth $m$. Implemented as a MATLAB toolbox
and validated across ten test configurations, from two-robot pairs to
five-robot nested hierarchies to a six-degree-of-freedom three-dimensional
cluster, this reduced what had been hours of manual derivation to seconds of
automated computation, while requiring no closed-form inverse kinematics at
all: the propagation worked directly from the transformation tree.

That automation carried one structural cost. Because each robot's velocity
was recomputed by walking its own path to the root from scratch, and every
node along a shared path was revisited once per descendant robot, the
prior algorithm's complexity was $O(n\cdot m\cdot p)$, linear in robot count
at the shallow depths its test cases used, but quadratic in hierarchy depth
in the worst case, since the rotation product from a deep sub-cluster to the
root was rebuilt every time that sub-cluster's robots were visited. The
thesis identified this cost explicitly and left removing it, along with
handling closed kinematic loops and optimizing the three-dimensional case
for real-time control, as future work.

This chapter is that future work carried out, restricted deliberately to the
planar case this dissertation needs and pursued further than closing the
complexity gap alone. Section~\ref{sec:assembly} below replaces the prior
algorithm's per-robot root-directed traversal with a single depth-first pass
that computes each node's sensitivity to its own parent exactly once and
passes it upward, removing the redundant rotation product and lowering
worst-case cost to $O(N\log N)$ on a balanced tree, matching the prior
algorithm's cost only in the same degenerate-chain case where both are
$O(N^2)$. Sections~\ref{sec:atoms} and~\ref{sec:grammar} replace the prior
work's generic homogeneous-transformation tree with a fixed library of four
closed-form charts and a dimension theorem proving the result complete and
square for any composition of them, a guarantee the transformation-tree
formulation, general enough to represent any topology, did not by itself
provide. Section~\ref{sec:gauge}'s consensus orientation gauge has no
antecedent in the prior work at all: it addresses a fragility, a single
pointer robot's bearing setting the entire formation frame, that only
becomes a practical concern once a formation is driven, as
Chapters~\ref{ch:vector_field} and~\ref{ch:separatrix} both drive one,
toward the low-signal neighborhood of a field critical point.

\section{Cluster-Space Charts and the Atom Library}
\label{sec:atoms}

Each robot $i$ sits at planar position $(x_i,y_i)$ in a fixed global frame.
A \emph{chart} maps the stacked robot-position vector to a cluster state
vector $\mathbf{q}$ of the same dimension: the \emph{forward} map runs
robots to cluster, and the \emph{inverse} map runs cluster to robots. Every
angle is measured with $\operatorname{atan2}$ and carried in radians. A
planar formation of $n$ robots has $2n$ degrees of freedom; three of these
are always rigid-body pose, two translation and one rotation, and the
remaining $2n-3$ describe internal shape. Every chart below honors this
split.

Almost every entry in every chart below is either a distance or a bearing
between two points. For a vector $\Delta = P_b - P_a$ with length
$L=\lVert\Delta\rVert$ and bearing $\beta=\operatorname{atan2}(\Delta_y,\Delta_x)$,
the gradients with respect to $P_b$ are
\begin{equation}
\frac{\partial L}{\partial P_b} = (\cos\beta,\sin\beta) \quad (\text{unit radial}),
\qquad
\frac{\partial \beta}{\partial P_b} = \frac{1}{L}(-\sin\beta,\cos\beta) \quad (\text{unit tangential}/L),
\label{eq:grads}
\end{equation}
with the opposite sign for $P_a$, and constant fractions for centroid rows.
Once these two patterns are in hand, assembling any chart's Jacobian is
bookkeeping. Four charts make up the atom library used in this dissertation.

\subsection{Three-robot SAS}
\label{sec:atom_sas}

State vector $\mathbf{q}=[x_c,y_c,\theta_c,p,\beta,q]^{\mathsf T}$ (six
entries for six degrees of freedom). Side $p=\lVert P_1-P_2\rVert$, side
$q=\lVert P_2-P_3\rVert$, and $\beta$ is the interior angle at robot 2.
Orientation is fixed to the $1\to2$ vector. This is the formation used
throughout Chapter~\ref{ch:vector_field}.

\paragraph{Forward (robots to cluster).}
\begin{align}
x_c &= \tfrac{1}{3}(x_1+x_2+x_3), \qquad y_c = \tfrac{1}{3}(y_1+y_2+y_3) \\
p &= \sqrt{(x_2-x_1)^2+(y_2-y_1)^2}, \qquad q = \sqrt{(x_3-x_2)^2+(y_3-y_2)^2} \\
\theta_c &= \operatorname{atan2}(y_2-y_1,\,x_2-x_1) \\
s &= \sqrt{(x_3-x_1)^2+(y_3-y_1)^2}, \qquad \beta = \arccos\!\left(\frac{p^2+q^2-s^2}{2pq}\right)
\end{align}
The diagonal $s=\lVert P_1-P_3\rVert$ is an intermediate quantity (this is
the triangle's third side length, renamed from the source material's $r$ to
avoid collision with the orbital radial magnitude of
Chapter~\ref{ch:vector_field}); $\beta$ then follows from the law of cosines
at vertex 2. A smoother alternative replaces the $\arccos$ with a signed
angle, $\beta = \operatorname{atan2}(a\times b,\,a\cdot b)$ with
$a=P_1-P_2,\,b=P_3-P_2$: the signed form carries the handedness of the
triangle (which side of the $1\to2$ baseline robot 3 lies on), so it covers
both mirror images rather than being two-to-one, and its derivative stays
finite as the triangle flattens, where the $\arccos$ derivative diverges.

\paragraph{Inverse (cluster to robots).}
Build the triangle in a local frame aligned so the $1\to2$ edge lies along
the local $x$-axis, then center, rotate, and translate:
\begin{align}
P_1^{\text{loc}} &= (0,0), \quad P_2^{\text{loc}} = (p,0), \quad
P_3^{\text{loc}} = (p-q\cos\beta,\, q\sin\beta) \\
x_c^{\text{loc}} &= \tfrac{1}{3}(2p-q\cos\beta), \qquad y_c^{\text{loc}} = \tfrac{1}{3}q\sin\beta \\
\begin{bmatrix} x_i \\ y_i \end{bmatrix} &= R(\theta_c)
  \begin{bmatrix} x_i^{\text{loc}} - x_c^{\text{loc}} \\ y_i^{\text{loc}} - y_c^{\text{loc}} \end{bmatrix}
  + \begin{bmatrix} x_c \\ y_c \end{bmatrix}
\end{align}
where $R(\theta)$ is the standard $2\times2$ rotation matrix.

\paragraph{Round-trip check.} Take $P_1=(0,0)$, $P_2=(2,0)$, $P_3=(1,1.5)$.
The forward map gives $x_c=1$, $y_c=0.5$, $\theta_c=0$, $p=2$,
$q=\sqrt{3.25}\approx1.8028$, with $s=\sqrt{3.25}$ so
$\cos\beta = 4/(2pq) = 0.5547$ and $\beta\approx0.9828$~rad. The inverse map
gives local $P_3^{\text{loc}}=(1.0,1.5)$, local centroid $(1,0.5)$, and,
since $\theta_c=0$, the rotation is the identity: the three points return
exactly to $(0,0)$, $(2,0)$, $(1,1.5)$.

\subsection{Four-robot cross-diagonal}
\label{sec:atom_xdiag}

State vector $\mathbf{q}=[x_c,y_c,\theta_c,d_1,d_2,r_1,r_2,\psi]^{\mathsf T}$
(eight entries). Diagonal 1 runs $1\to3$ with length $d_1$ and sets
$\theta_c$; diagonal 2 runs $2\to4$ with length $d_2$; $\psi$ is the angle
from diagonal 1 to diagonal 2; $r_1,r_2\in[0,1]$ are the fractional crossing
positions of the two diagonals. This chart is used in
Appendix~\ref{app:four_robot} as the kinematic layer underneath the
Four-Robot Cluster Space Controller.

\paragraph{Forward (robots to cluster).}
\begin{align}
x_c &= \tfrac{1}{4}\textstyle\sum_i x_i, \qquad y_c = \tfrac{1}{4}\textstyle\sum_i y_i \\
d_1 &= \lVert P_3-P_1\rVert, \qquad d_2 = \lVert P_4-P_2\rVert \\
\theta_c &= \operatorname{atan2}(y_3-y_1,\,x_3-x_1), \qquad \psi = \operatorname{atan2}(y_4-y_2,\,x_4-x_2)-\theta_c
\end{align}
With $A=P_3-P_1$, $B=P_4-P_2$, $C=P_2-P_1$, and the two-dimensional cross
product $u\times v = u_xv_y - u_yv_x$,
\begin{equation}
r_1 = \frac{C\times B}{A\times B}, \qquad r_2 = \frac{C\times A}{A\times B}.
\end{equation}
The denominator $A\times B$ vanishes when the diagonals are parallel
($\psi\to0,\pi$); this is the chart's singularity, and it sits near every
near-collinear formation.

\paragraph{Inverse (cluster to robots).}
Let $\hat u_1=(\cos\theta_c,\sin\theta_c)$ and
$\hat u_2=(\cos(\theta_c+\psi),\sin(\theta_c+\psi))$. First place the
diagonals' crossing point relative to the centroid, then step out along
each diagonal:
\begin{align}
\begin{bmatrix} x_{\text{int}} \\ y_{\text{int}} \end{bmatrix}
  &= \begin{bmatrix} x_c \\ y_c \end{bmatrix}
  + \tfrac{1}{4}\Big[d_1(2r_1-1)\hat u_1 + d_2(2r_2-1)\hat u_2\Big] \\
P_1 &= X - r_1 d_1 \hat u_1, \qquad P_3 = X + (1-r_1)d_1\hat u_1 \\
P_2 &= X - r_2 d_2 \hat u_2, \qquad P_4 = X + (1-r_2)d_2\hat u_2
\end{align}
where $X=(x_{\text{int}},y_{\text{int}})$; the $\tfrac{1}{4}[\cdot]$ term
recenters the crossing point because the centroid is the mean of the four
corners, equivalently the midpoint of the two diagonal midpoints.

\paragraph{Round-trip check.} Take $P_1=(0,0)$, $P_2=(1,-1)$, $P_3=(4,0)$,
$P_4=(1,2)$. The forward map gives centroid $(1.5,0.25)$, $d_1=4$,
$\theta_c=0$, $d_2=3$, $\psi=\pi/2$; the cross products
$A\times B=12$, $C\times B=3$, $C\times A=4$ give $r_1=0.25$, $r_2=\tfrac13$.
The inverse map gives $\hat u_1=(1,0)$, $\hat u_2=(0,1)$,
$X=(1.5,0.25)+\tfrac14[(-2,0)+(0,-1)]=(1,0)$, and stepping out returns
exactly $(0,0),(1,-1),(4,0),(1,2)$.

\subsection{Four-robot H-frame / spine}
\label{sec:atom_hframe}

State vector
$\mathbf{q}=[x_c,y_c,\theta_c,L_s,\theta_{c1},w_1,\theta_{c2},w_2]^{\mathsf T}$
(eight entries). Front pair $\{1,2\}$ forms sub-cluster 1, rear pair
$\{3,4\}$ forms sub-cluster 2. The spine connects the two sub-centroids:
$\theta_c$ is the spine bearing, $L_s$ its length. Sub-angles are kept
absolute, which leaves the two layers cleanly block-separable.

\paragraph{Forward (robots to cluster).} Sub-cluster 1:
\begin{equation}
x_{c1}=\tfrac12(x_1+x_2), \quad y_{c1}=\tfrac12(y_1+y_2), \quad
\theta_{c1}=\operatorname{atan2}(y_2-y_1,x_2-x_1), \quad w_1=\lVert P_2-P_1\rVert.
\end{equation}
Sub-cluster 2 is identical with $\{3,4\}$ in place of $\{1,2\}$. At the meta
level, treating the two sub-centroids $C_1, C_2$ as virtual robots,
\begin{equation}
x_c=\tfrac12(x_{c1}+x_{c2}), \quad y_c=\tfrac12(y_{c1}+y_{c2}), \quad
\theta_c=\operatorname{atan2}(y_{c2}-y_{c1},x_{c2}-x_{c1}), \quad L_s=\lVert C_2-C_1\rVert.
\end{equation}

\paragraph{Inverse (cluster to robots).} Meta level to sub-centroids:
\begin{equation}
C_1 = \left(x_c - \tfrac{L_s}{2}\cos\theta_c,\ y_c-\tfrac{L_s}{2}\sin\theta_c\right), \qquad
C_2 = \left(x_c + \tfrac{L_s}{2}\cos\theta_c,\ y_c+\tfrac{L_s}{2}\sin\theta_c\right).
\end{equation}
Sub-centroid to robots:
\begin{align}
P_1 &= C_1 - \tfrac{w_1}{2}(\cos\theta_{c1},\sin\theta_{c1}), \qquad
P_2 = C_1 + \tfrac{w_1}{2}(\cos\theta_{c1},\sin\theta_{c1}) \\
P_3 &= C_2 - \tfrac{w_2}{2}(\cos\theta_{c2},\sin\theta_{c2}), \qquad
P_4 = C_2 + \tfrac{w_2}{2}(\cos\theta_{c2},\sin\theta_{c2}).
\end{align}
There are three independent collapse singularities, $L_s\to0$, $w_1\to0$,
$w_2\to0$, each a physically obvious degeneration that can be floored
directly, which also keeps the sampling stencil rank-2.

\paragraph{Round-trip check.} Take the rectangle $P_1=(0,0)$, $P_2=(0,2)$,
$P_3=(4,0)$, $P_4=(4,2)$. The forward map gives sub-centroids $(0,1)$ and
$(4,1)$, $\theta_{c1}=\theta_{c2}=\pi/2$, $w_1=w_2=2$, meta centroid
$(2,1)$, $\theta_c=0$, $L_s=4$. The inverse map gives $C_1=(0,1)$,
$C_2=(4,1)$, and stepping out by $w/2$ along
$\theta_{c1}=\theta_{c2}=\pi/2$ returns exactly $(0,0),(0,2),(4,0),(4,2)$.

\subsection{Six-robot pentagon plus center}
\label{sec:atom_pentagon}

State vector
$\mathbf{q}=[x_h,y_h,\theta_c,r_1,\ldots,r_5,\gamma_2,\ldots,\gamma_5]^{\mathsf T}$
(twelve entries). Robots 1 through 5 are the vertices, robot 6 is the
center and serves as the hub. Orientation $\theta_c$ is the hub-to-vertex-1
bearing, so $\gamma_1\equiv0$; the other vertices carry angular offsets
$\gamma_i$ relative to $\theta_c$. This is the formation used throughout
Chapter~\ref{ch:separatrix}. The chart is on the decoupled end of the atom
library: perturbing a non-reference vertex ($i=2,\ldots,5$) touches only its
own $r_i$ and $\gamma_i$. Two robots are shared rather than decoupled: the
hub sets the origin of every spoke, so a move of robot 6 reaches all five,
and vertex 1 sets $\theta_c$, so a move of vertex 1 shifts the frame and
with it every $\gamma_j$. Section~\ref{sec:gauge} removes the second
privilege.

\paragraph{Forward (robots to cluster).}
$x_h=x_6$, $y_h=y_6$. For each vertex $i=1,\ldots,5$, with
$\Delta x_i = x_i - x_h$, $\Delta y_i = y_i - y_h$,
\begin{equation}
r_i = \sqrt{\Delta x_i^2 + \Delta y_i^2}, \qquad \beta_i = \operatorname{atan2}(\Delta y_i, \Delta x_i),
\end{equation}
\begin{equation}
\theta_c = \beta_1, \qquad \gamma_i = \beta_i - \theta_c \quad (i=2,\ldots,5,\ \text{wrapped to } (-\pi,\pi]).
\end{equation}

\paragraph{Inverse (cluster to robots).}
$x_6=x_h$, $y_6=y_h$, and for $i=1,\ldots,5$ (with $\gamma_1\equiv0$),
\begin{equation}
x_i = x_h + r_i\cos(\theta_c+\gamma_i), \qquad y_i = y_h + r_i\sin(\theta_c+\gamma_i).
\end{equation}
The only singularity is $r_i\to0$: a vertex meeting the hub, its bearing
undefined. Robots 6 and 1 are privileged under this gauge, since the hub
sets the origin of every spoke and vertex 1 sets the angular reference.

\paragraph{Target formation.} A regular pentagon of circumradius $R$ about
the center: $r_1=\cdots=r_5=R$, $\gamma_2=2\pi/5$, $\gamma_3=4\pi/5$,
$\gamma_4=6\pi/5$, $\gamma_5=8\pi/5$, with $x_h,y_h,\theta_c$ free. At this
setpoint the five vertices average back to the hub, so robot 6 coincides
with the geometric centroid and $(x_h,y_h)$ doubles as a centroid navigation
handle. This is exactly the \texttt{pentagon\_small.yaml} configuration used
throughout Chapter~\ref{ch:separatrix}.

\paragraph{Round-trip check.} Hub $(5,5)$, $R=1$, $\theta_c=\pi/2$. The
inverse map gives vertex 1 at $(5,6)$, vertex 2 at approximately
$(4.0489,5.3090)$, and so on around the ring. The forward map from those
positions returns exactly the hub $(5,5)$, every $r_i=1$, $\theta_c=90°$,
and $\gamma=\{72°,144°,216°,288°\}$.

\section{A Compositional Grammar for Formations}
\label{sec:grammar}

The four charts above are instances of one construction. A cluster is a
base chart applied to an ordered list of children, where each child is
either a robot or another cluster:
\begin{equation}
\texttt{cluster} ::= \texttt{chart(} \text{child, child,} \ldots \texttt{)}
\qquad \texttt{child} ::= \texttt{robot} \mid \texttt{cluster}.
\end{equation}
Leaves are robots, internal nodes are sub-clusters, and the whole formation
is a composition tree. The base chart at a node is one of the atoms already
derived: the two-robot pair, the three-robot SAS triangle, or a $k$-body
star for higher arity. A node does not care whether a given child is a real
robot or the representative point of a sub-cluster; it only places points.
That single fact is what lets robots and sub-clusters sit side by side as
siblings.

\begin{theorem}[Completeness and squareness]
\label{thm:dim}
For any composition tree over $n$ robots, the chart produced by this
grammar has exactly $2n$ variables.
\end{theorem}
\begin{proof}[Proof sketch]
The root contributes two variables (global translation). Every internal
node with $k\geq2$ children contributes $2k-2$: one orientation plus
$2k-3$ shape (a one-child node contributes nothing, discussed below). With
$m$ internal nodes and $n$ leaves, the child links sum to $\sum k=m+n-1$, so
\begin{equation}
\dim\mathbf{q} = 2 + \sum_{\text{nodes}}(2k-2) = 2 + 2(m+n-1)-2m = 2n
\end{equation}
for any tree at all. The chart is complete and square by construction; a
formation is never sized by hand.
\end{proof}

A one-child node is the identity node: $x_{\text{node}}=x_{\text{child}}$,
with no orientation and no shape. It owns $2(1)-2=0$ variables, so it cannot
supply a missing degree of freedom, but it is not forbidden: it is the hub
mechanism. The pentagon chart's $x_h=x_6$ designates robot 6 as the
formation's position reference, which is a cluster of one in all but name.
Oddness needs no phantom padding: use an odd-arity atom where it fits, the
SAS triangle is the natural one, or hang the spare robot as a leaf child of
a node that also carries sub-clusters.

\paragraph{Two worked trees.} Nine robots can be built as three
sub-clusters of three: a SAS meta-node over three sub-centroids, each
branch itself a SAS triangle over three robots. The dimension count is
$2+4+3\times4=18=2\cdot9$. Because every branch holds the same number of
robots, the meta-centroid equals the true centroid of all nine, so the
navigation handle stays the real centroid. Eight robots can instead be
built as a balanced binary tree of two-robot pairs, depth three: a pair over
two pairs-of-pairs over four pairs over eight robots. The dimension count is
$2+7\times2=16=2\cdot8$: seven internal two-body nodes, each the same
reusable pair block, give the full $16\times16$ Jacobian as a three-layer
cascade. The flat single cluster (Section~\ref{sec:atoms}) and this deep
binary tree are the two extremes of the same grammar.

% %% TODO: SYNC WITH LATEST DRAFT -- the source HTML (cluster_kinematics3.html,
% Section 8) presents these two worked trees as inline SVG figures with
% dimension-count annotations. Redraw both as TikZ figures here (a nine-robot
% tree of three SAS-3 sub-clusters, and an eight-robot balanced binary tree of
% pair blocks) rather than screenshotting the HTML's SVGs, per
% Thesis_Execution_Plan.md Section 6, item 2.

\section{Single-Pass Inverse Jacobian Assembly}
\label{sec:assembly}

Because both the forward and inverse maps of Section~\ref{sec:atoms} are
closed form, both the forward Jacobian $\mathbf{J}_c=\partial\mathbf{q}/\partial\mathbf{r}$
and the inverse Jacobian $\mathbf{J}_c^{-1}=\partial\mathbf{r}/\partial\mathbf{q}$
are obtained by direct symbolic differentiation; no matrix is ever inverted
numerically. (Here $\mathbf{J}_c$ denotes the kinematic Jacobian, renamed
from the source material's bare $J$ to avoid collision with the field
Jacobian $\mathbf{J}$ of Chapters~\ref{ch:vector_field}
and~\ref{ch:separatrix}.) The inverse is literally the matrix inverse of the
forward map, so $\mathbf{J}_c^{-1}\mathbf{J}_c=\mathbf{I}$ is a free
consistency check on the algebra.

\subsection{The reusable two-robot block}

The pair map $(x_a,y_a,x_b,y_b)\to(x_c,y_c,\theta,L)$ is the atom that the
H-frame chart and the cluster-of-clusters construction below are built
from. Its forward Jacobian is
\begin{equation}
\mathbf{J}_{\text{pair}} =
\begin{bmatrix}
\tfrac12 & 0 & \tfrac12 & 0 \\[2pt]
0 & \tfrac12 & 0 & \tfrac12 \\[2pt]
\dfrac{\sin\theta}{L} & -\dfrac{\cos\theta}{L} & -\dfrac{\sin\theta}{L} & \dfrac{\cos\theta}{L} \\[6pt]
-\cos\theta & -\sin\theta & \cos\theta & \sin\theta
\end{bmatrix},
\label{eq:jpair}
\end{equation}
and its inverse, the differential of $x_a=x_c-\tfrac{L}{2}\cos\theta$ and
the three companion equations, is
\begin{equation}
\mathbf{J}_{\text{pair}}^{-1} =
\begin{bmatrix}
1 & 0 & \tfrac{L}{2}\sin\theta & -\tfrac12\cos\theta \\[3pt]
0 & 1 & -\tfrac{L}{2}\cos\theta & -\tfrac12\sin\theta \\[3pt]
1 & 0 & -\tfrac{L}{2}\sin\theta & \tfrac12\cos\theta \\[3pt]
0 & 1 & \tfrac{L}{2}\cos\theta & \tfrac12\sin\theta
\end{bmatrix},
\label{eq:jpair_inv}
\end{equation}
whose rows are $(x_a,y_a,x_b,y_b)$ and whose columns are $(x_c,y_c,\theta,L)$.
Multiplying the two confirms $\mathbf{J}_{\text{pair}}\mathbf{J}_{\text{pair}}^{-1}=\mathbf{I}$;
the $\theta$ and $L$ columns rely on $\sin^2\theta+\cos^2\theta=1$. For a
hierarchical chart the total Jacobian is a product of per-layer Jacobians,
exactly as a serial manipulator stacks per-joint contributions:
$\mathbf{J}_{c,\text{total}}=\mathbf{J}_{c,\text{outer}}\mathbf{J}_{c,\text{inner}}$,
so the only object ever inverted by hand is the $4\times4$ block above.

\subsection{Recursive assembly for an arbitrary tree}

The chain rule becomes concrete with one object per node: the
\emph{sensitivity matrix} $S_v$, a $(2n_v\times3)$ matrix whose rows give
the Cartesian velocity of each robot in the subtree rooted at $v$, and whose
three columns correspond to a unit velocity in $(x_c^v,y_c^v,\theta^v)$, the
centroid translation and orientation of that sub-cluster. The recursion has
two parts.

\paragraph{Leaves.} For a size-1 leaf (identity node, robot $i$ is the
sub-centroid),
\begin{equation}
S_{\text{leaf},1} = \begin{bmatrix} 1 & 0 & 0 \\ 0 & 1 & 0 \end{bmatrix}.
\end{equation}
For a size-2 leaf with state $(x_c,y_c,\theta,L)$, the pair block gives
$\mathbf{J}_{\text{pair}}^{-1}$ and $S$ is its first three columns:
\begin{equation}
S_{\text{leaf},2} = \mathbf{J}_{\text{pair}}^{-1}[:,0{:}3] =
\begin{bmatrix}
1 & 0 & \tfrac{L}{2}\sin\theta \\
0 & 1 & -\tfrac{L}{2}\cos\theta \\
1 & 0 & -\tfrac{L}{2}\sin\theta \\
0 & 1 & \tfrac{L}{2}\cos\theta
\end{bmatrix}.
\end{equation}
For a size-3 leaf (SAS) with state $(x_c,y_c,\theta_c,p,\beta,q)$, $S$ is the
first three columns of the SAS chart's own inverse Jacobian
$\mathbf{J}_{c,\text{SAS}}^{-1}$ (Appendix~\ref{app:cluster_kinematics}
derives this block in full). The leaf also directly fills the columns of
the global $\mathbf{J}_c^{-1}$ corresponding to its own shape variables
($L$, or $p,\beta,q$) from the remaining columns of its block.

\paragraph{Internal nodes.} Let node $v$ have children $c_1,\ldots,c_k$ with
sensitivity matrices $S_{c_1},\ldots,S_{c_k}$ already computed. Form the
meta-block $\mathbf{J}_{c,v}^{-1}$, mapping $v$'s variables
$(x_c^v,y_c^v,\theta^v,\text{shape}^v)$ to the $k$ child centroid positions:
for a binary node this is the $4\times4$ pair block; for a ternary node it
is the $6\times6$ SAS block evaluated at the child centroids. Partition by
child,
\begin{equation}
\mathbf{J}_{c,v}^{-1} = \begin{bmatrix} \mathbf{J}_{c,v,c_1}^{-1} \\ \mathbf{J}_{c,v,c_2}^{-1} \\ \vdots \end{bmatrix},
\qquad \mathbf{J}_{c,v,c_i}^{-1} \text{ is } (2\times\dim\mathbf{q}_v).
\end{equation}
Then the sensitivity of every robot in child $c_i$ to the variables of $v$
is
\begin{equation}
\frac{\partial \mathbf{r}_{c_i}}{\partial \mathbf{q}_v} = S_{c_i} \cdot \mathbf{J}_{c,v,c_i}^{-1}[:,0{:}2]^{\mathsf T}.
\end{equation}
Only the first two columns of $\mathbf{J}_{c,v,c_i}^{-1}$ (the $x_c,y_c$
response) enter this product, since every orientation variable in this
tree is an \emph{absolute} angle: a change in the parent's $\theta^v$ does
not alter any child's own orientation directly, only the child centroid's
position, which those first two columns already capture. (If spine-relative
angles $\alpha_{\text{spine},i}=\theta^{c_i}-\theta^v$ are used instead,
each child's orientation co-rotates with the parent and the full
three-column chain applies.) The columns of the global $\mathbf{J}_c^{-1}$
for $v$'s own shape variables are filled analogously from
$\mathbf{J}_{c,v,c_i}^{-1}[:,3{:}]$ chained through $S_{c_i}$, and $v$'s own
sensitivity matrix, passed up to $v$'s parent, is assembled from all
children's contributions in the same way.

One depth-first pass fills every column of $\mathbf{J}_c^{-1}$ and returns
the root's sensitivity matrix, which is discarded (the root has no parent).
The assembled matrix is the complete $(2N\times2N)$ inverse Jacobian. This
recursion is exactly what \texttt{clusterbuilder.py} implements in
\texttt{\_jac\_node}: each call receives the global $\mathbf{J}_c^{-1}$
matrix and a column-index map, fills its own shape columns in place, and
returns its sensitivity matrix; the root call additionally fills the
$(x_c,y_c,\theta_c)$ columns using the same chain rule applied to the root
meta-block and all child sensitivities.

\begin{algorithm}[t]
\caption{Single-pass inverse Jacobian assembly}
\label{alg:assemble}
\begin{algorithmic}[1]
\Function{Assemble}{node $v$, matrix $\mathbf{J}_c^{-1}$, column map}
  \If{$v$ is a leaf}
    \State build the leaf block (identity, pair, or SAS) at $v$'s robots
    \State write $v$'s shape columns into $\mathbf{J}_c^{-1}$ in place
    \State \Return $S_v \gets$ first three columns of the leaf inverse block
  \EndIf
  \For{each child $c_i$ of $v$}
    \State $S_{c_i} \gets$ \Call{Assemble}{$c_i$, $\mathbf{J}_c^{-1}$, map}
  \EndFor
  \State form meta-block $\mathbf{J}_{c,v}^{-1}$: $v$'s variables $\to$ child centroids
  \For{each child $c_i$}
    \State write $S_{c_i}\, \mathbf{J}_{c,v,c_i}^{-1}[:,0{:}2]$ into the rows of $c_i$
    \State fill $v$'s shape columns from $\mathbf{J}_{c,v,c_i}^{-1}[:,3{:}]$ via $S_{c_i}$
  \EndFor
  \State \Return $S_v$ assembled from all $S_{c_i}$
\EndFunction
\end{algorithmic}
\end{algorithm}

\subsection{Complexity}

A prior hierarchical toolbox built the inverse Jacobian in $O(n)$ operations
per cluster variable and per robot, but at $O(n^2)$ in the depth of the
hierarchy, because the rotation product from each node to the root was
recomputed at every step; that prior work flagged this recomputation
explicitly as future work. The recursive method above eliminates it: each
node computes its sensitivity matrix once and passes it to its parent, so
no ancestor ever re-descends to recompute it. Within a single depth-first
pass, an internal node at depth $d$ chains over every robot in its subtree,
so a robot at depth $d$ is touched once per ancestor, and the total work is
$\sum_{i=1}^N \text{depth}(i)$. For a balanced binary tree this is
$O(N\log N)$; at bounded depth it is $O(N)$; only a degenerate chain (every
robot a direct child of the previous) reaches $O(N^2)$, matching the prior
worst case while beating it on every balanced tree and removing the
redundant rotation product entirely (Table~\ref{tab:complexity}).

\begin{table}[t]
\centering
\caption{Assembly cost by tree shape.}
\label{tab:complexity}
\begin{tabular}{lccc}
\toprule
Tree shape & Depth & This work & Prior penalty \\
\midrule
Flat (single cluster) & 1 & $O(N)$ & none \\
Balanced binary & $\log N$ & $O(N\log N)$ & extra depth factor \\
Degenerate chain & $N$ & $O(N^2)$ & matches (worst case) \\
\bottomrule
\end{tabular}
\end{table}

\section{Orientation Gauge: Pointer versus Consensus}
\label{sec:gauge}

The pentagon chart of Section~\ref{sec:atom_pentagon} fixes $\theta_c$ to a
single pointer robot, $\theta_c=\beta_1$. That keeps the orientation row of
the Jacobian sparse but stakes the entire global frame on one bearing. The
choice is a gauge: any rule that reads the formation's rotation from its
spokes is admissible. This section sets that sparse pointer gauge against a
consensus gauge that spreads orientation across every spoke, and shows what
the swap costs.

\subsection{The pointer gauge}

Baseline from Section~\ref{sec:atom_pentagon}: $\theta_c=\beta_1$ and
$\gamma_1\equiv0$. Vertex 1 sets the frame, the hub anchors position. The
orientation row of the forward Jacobian is nonzero only at those two
robots, so the chart is as sparse as a star can be. The exposure is that
$r_1\to0$ sends $\theta_c$, and through it every reconstructed vertex, to an
undefined frame. A formation loitering on a critical point, where the field
signal is weakest and the pointer is most likely to wander, is exactly
where staking the frame on one spoke hurts most, and is exactly the
situation Chapter~\ref{ch:separatrix}'s separatrix and Okubo-Weiss
controllers create by design.

\subsection{The consensus gauge}

Give each spoke a nominal slot $\gamma_i^{\text{set}}$. Each spoke then
implies an orientation estimate $\beta_i-\gamma_i^{\text{set}}$, and the
gauge aligns the template to the observed spokes in least squares. The
radius-weighted (Procrustes) optimum is
\begin{equation}
\theta_c = \operatorname{atan2}(S,C), \qquad
C = \sum_i\left(\Delta x_i\cos\gamma_i^{\text{set}} + \Delta y_i \sin\gamma_i^{\text{set}}\right), \qquad
S = \sum_i\left(\Delta y_i\cos\gamma_i^{\text{set}} - \Delta x_i \sin\gamma_i^{\text{set}}\right),
\label{eq:consensus_theta}
\end{equation}
with $\Delta x_i=x_i-x_h,\Delta y_i=y_i-y_h$. When every radius is equal
this reduces to the plain circular mean,
$\theta_c=\operatorname{atan2}\big(\sum_i\sin(\beta_i-\gamma_i^{\text{set}}),\,\sum_i\cos(\beta_i-\gamma_i^{\text{set}})\big)$.
The radius weighting is not cosmetic: a bearing-only mean still reads
$\beta_i$, which is undefined as $r_i\to0$, so it would inherit the very
singularity the gauge is meant to remove; the $\Delta$-form instead lets a
vanishing spoke fall out of $C$ and $S$. Offsets are then taken for all five
spokes, $\gamma_i=\beta_i-\theta_c$, and $\gamma_1$ is no longer identically
zero.

\subsection{The companion identity}

Pinning $\gamma_1\equiv0$ used vertex 1 as the gauge; the consensus replaces
it with the stationarity condition that defines $\theta_c$:
\begin{equation}
\sum_{i=1}^5 r_i \sin\!\left(\gamma_i - \gamma_i^{\text{set}}\right) = 0.
\label{eq:companion_identity}
\end{equation}
Any measured formation satisfies this automatically, since it is exactly
what makes $\theta_c$ the optimum, so it never constrains a reading. It
constrains only commanded setpoints: the state carries thirteen coordinates
on a twelve-degree-of-freedom surface, the same dimension as the pointer
chart, with the single redundancy absorbed symmetrically rather than
dropped on one robot. The regular target $\gamma_i=\gamma_i^{\text{set}}$
meets it as $\sin 0=0$, so in practice one commands four offsets and lets
the fifth close the constraint. Forward and inverse maps otherwise follow
Section~\ref{sec:atom_pentagon} unchanged, with the one orientation line
replaced by Equation~\eqref{eq:consensus_theta}.

\subsection{What the gauge costs in the Jacobian}

One row moves. With the bearing derivative
$\partial\beta/\partial P_b=\tfrac1L(-\sin\beta,\cos\beta)$ of
Equation~\eqref{eq:grads}, the pointer row is nonzero only at vertex 1 and
the hub. The consensus row spreads over all six robots,
\begin{equation}
\frac{\partial\theta_c}{\partial x_i} = \frac{-C\sin\gamma_i^{\text{set}} - S\cos\gamma_i^{\text{set}}}{C^2+S^2}, \qquad
\frac{\partial\theta_c}{\partial y_i} = \frac{C\cos\gamma_i^{\text{set}} - S\sin\gamma_i^{\text{set}}}{C^2+S^2},
\end{equation}
with the hub entry the negative sum over the vertex entries. At the regular
setpoint $C^2+S^2=25R^2$, and this collapses to
\begin{equation}
\frac{\partial\theta_c}{\partial P_i} = \frac{1}{5}\cdot\frac{1}{r_i}\left(-\sin\beta_i,\cos\beta_i\right).
\end{equation}
The pointer gauge loads the whole orientation sensitivity onto one spoke at
$1/r_1$; the consensus splits it five ways at $1/(5r_i)$. A noisy bearing at
any one vertex moves $\theta_c$ five times less, and as a single $r_i\to0$
that vertex drops out of $C$ and $S$ rather than diverging. The shape angles
pay for it: since $\gamma_i=\beta_i-\theta_c$, every $\gamma_i$ row picks up
the same $-\partial\theta_c/\partial\mathbf{r}$ term. That shared correction
is a rank-one update to the otherwise per-spoke shape block, so it inverts
by the Sherman-Morrison identity without rebuilding the chart; the
decoupling degrades by exactly one rank, not into a dense block.

The pointer gauge fails when $r_1\to0$, a per-robot collapse that takes the
global frame with it. The consensus gauge has no privileged robot; its only
orientation failure is $\rho_{\text{gauge}}=\sqrt{C^2+S^2}\to0$, a
simultaneous collapse of the whole alignment that sits far outside any
regime where $\rho_{\text{gauge}}\approx\sum_i r_i$. (This
$\rho_{\text{gauge}}$ is distinct from the polar radius $\rho$ used in
Appendix~\ref{app:vector_fields}'s field definitions and from the pentagon
ring radius $R$ above; it is always written with the \texttt{gauge}
subscript to avoid confusion.) The irreducible $r_i\to0$ singularity remains
in both gauges, but under consensus it costs only that spoke's $\gamma_i$,
not the frame. The pointer gauge is preferable when the formation stays
well-conditioned and the sparsest Jacobian is worth more than the
redundancy; the consensus gauge is preferable when the frame must survive a
wandering or near-coincident reference, the operating regime of a cluster
parked on a critical point, exactly the regime studied in
Chapters~\ref{ch:vector_field} and~\ref{ch:separatrix}.

\section{Edge Cases and Choosing a Structure}
\label{sec:edge_cases}

Different formation geometries stress different charts. The cross-diagonal
chart (Section~\ref{sec:atom_xdiag}) is the one that suffers from
non-convexity: when the quadrilateral is non-convex, the two diagonals can
fail to cross inside their segments, so the ratios $r_1,r_2$ leave $[0,1]$.
Stars and the SAS triangle make no convexity assumption and represent these
arrangements without complaint. A star chart represents a collinear
formation with no chart singularity at all, since the spokes simply take
bearings of $0$ or $\pi$; the cross-diagonal chart is worst here too, its
diagonals turn parallel and it blows up, while the SAS triangle passes
through a flat configuration only in the signed-angle form of
Section~\ref{sec:atom_sas}. A surviving chart is a kinematic fact only: a
line has no two-dimensional spread, so it stays useless for recovering a
two-dimensional gradient or Hessian regardless of which chart is used, a
point directly relevant to Chapters~\ref{ch:vector_field}
and~\ref{ch:separatrix}, both of which require a formation with genuine
areal extent. Charts whose shape angles are unsigned, the $\arccos$ form of
the SAS triangle among them, cannot separate a formation from its mirror
image and are two-to-one; signed $\operatorname{atan2}$ angles carry the
handedness instead. Finally, any distance or bearing between two named
robots is undefined when they meet: coincident robots are the one
irreducible singularity shared by every chart in the library.

Composition (Section~\ref{sec:grammar}) earns its place when a formation is
large, when one part should be driven without moving the rest, or when
dependence should be distributed: its inverse Jacobian is block-sparse, so
a command to one sub-cluster's shape or orientation maps only to that
sub-cluster's robots and is zero elsewhere, which is what lets one wing of
a formation reshape for local sampling while the rest holds station. A flat
star hangs every robot off one shared orientation and one shared anchor, so
a single $\dot\theta_c$ rotates the entire formation and any disturbance at
the hub reaches all spokes; the consensus gauge of
Section~\ref{sec:gauge} softens half of this cost, but the shared hub anchor
remains. The flat star is preferable when there is a natural center, when
one variable should rotate the whole group, or when the formation must
morph through a straight line, where the star alone stays
singularity-free; composition is preferable for scale, for reusing one
block, and for localized control. Different trees over the same robots are
different charts, with different shape variables and singular loci, so the
tree is still chosen to match the formation, and conditioning compounds
along a cascade: a single ill-conditioned layer, a short spine against wide
pairs for instance, propagates through the whole chart.

\section{Formations Used in This Dissertation}
\label{sec:formations_used}

Two instances of this atom library carry the rest of this dissertation.
Chapter~\ref{ch:vector_field}'s three-robot estimator uses the SAS triangle
of Section~\ref{sec:atom_sas} directly, a single flat node, no composition
required. Chapter~\ref{ch:separatrix}'s six-robot estimator uses the
pentagon-plus-center chart of Section~\ref{sec:atom_pentagon}, under the
consensus gauge of Section~\ref{sec:gauge} where formation loitering near a
saddle or an Okubo-Weiss boundary makes the pointer gauge's fragility a live
concern. Appendix~\ref{app:four_robot}'s supplementary four-robot material
uses the cross-diagonal chart of Section~\ref{sec:atom_xdiag} as its
kinematic layer, giving that appendix's material a second purpose here: a
third worked instance of this chapter's general grammar, at a formation
size between the triangle and the pentagon.

\section{Validation}
\label{sec:cb_validation}

Four claims are validated: that every chart inverts exactly, that the
assembled Jacobian is consistent, that assembly scales as
Table~\ref{tab:complexity} predicts, and that the consensus gauge improves
conditioning near a critical point. All experiments use the open-source
generator (\texttt{cluster\_builder/clusterbuilder.py}); the recursion of
Algorithm~\ref{alg:assemble} is the implementation under test.

\subsection{Round-trip exactness}

For each chart in the library, the three-robot SAS triangle, the four-robot
cross-diagonal, the four-robot H-frame, and the six-robot
pentagon-plus-center, random admissible formations are sampled, mapped from
robots to cluster variables and back, and the recovery error is measured.
Across the library the round trip returns the original positions to
machine precision; the worked numerical examples in
Sections~\ref{sec:atom_sas}--\ref{sec:atom_pentagon} above are representative
cases. The reconstructions confirm that each closed-form inverse is a true
inverse of its forward map, not a least-squares approximation.

\subsection{Jacobian consistency}

Because both maps are closed form, the forward Jacobian
$\mathbf{J}_c=\partial\mathbf{q}/\partial\mathbf{r}$ and the inverse
Jacobian $\mathbf{J}_c^{-1}=\partial\mathbf{r}/\partial\mathbf{q}$ are
obtained by direct symbolic differentiation, with no numerical matrix
inversion anywhere in the pipeline. The identity
$\mathbf{J}_c^{-1}\mathbf{J}_c=\mathbf{I}$ is verified at every sampled
configuration; the residual $\lVert \mathbf{J}_c^{-1}\mathbf{J}_c-\mathbf{I}\rVert$
stays at the level of floating-point round-off for every chart and for
composed trees. This identity is the assembly's built-in correctness test:
a single mis-placed column or a stale rotation term breaks it immediately,
which is what makes the symbolic path a useful oracle for the numeric one.

\subsection{Assembly scaling}

Sweeping the robot count $N$ over three tree shapes, a flat single cluster,
a balanced binary tree of pairs, and a degenerate chain, and recording
assembly time against a baseline that recomputes the node-to-root rotation
product at every level (reproducing the prior hierarchical method), the
flat and balanced trees follow the $O(N)$ and $O(N\log N)$ slopes of
Table~\ref{tab:complexity}, while the baseline carries the extra depth
factor on the balanced tree and separates from the single-pass curve as
depth grows. The two methods coincide only on the degenerate chain, where
both reach $O(N^2)$, confirming that the single pass removes a depth
penalty rather than trading one cost for another.

% %% TODO: SYNC WITH LATEST DRAFT -- fig:scaling is an unfilled placeholder in
% the source IDETC-II draft. Generate real assembly-time-vs-N curves from the
% cluster_builder/tests/ suite (flat, balanced, chain, vs. rotation-product
% baseline) before this figure is finalized, per Thesis_Execution_Plan.md
% Section 6, item 1.
\begin{figure}[t]
\centering
\framebox[0.9\linewidth]{\rule{0pt}{3.0cm}}
\caption{Measured assembly time versus robot count $N$ for three tree
shapes, single-pass assembly against the rotation-product baseline.}
\label{fig:scaling}
\end{figure}

\subsection{Block sparsity and localized singularities}

On a sixteen-robot balanced binary tree, the sparsity pattern of the
assembled inverse Jacobian shows that a command to one sub-cluster's shape
or orientation maps only to that sub-cluster's robots and is zero
elsewhere, so the inverse Jacobian is block-sparse by construction
(Figure~\ref{fig:sparsity}). This is the property that lets one wing of a
large formation reshape for local sampling while the rest holds station.
The chart's singular set is the union of the per-node singularities, each a
physically readable collapse (a vanishing spine, a vanishing sub-cluster
width, a spoke meeting its hub) rather than the entangled blow-up of a
monolithic crossing-diagonal chart; singularities grow more numerous as
nodes multiply but remain local and individually interpretable.

% %% TODO: SYNC WITH LATEST DRAFT -- fig:sparsity is an unfilled placeholder;
% generate a real spy plot of a 16-robot balanced binary tree's inverse
% Jacobian, per Thesis_Execution_Plan.md Section 6, item 1.
\begin{figure}[t]
\centering
\framebox[0.9\linewidth]{\rule{0pt}{3.0cm}}
\caption{Sparsity pattern of the inverse Jacobian for a sixteen-robot
balanced binary tree.}
\label{fig:sparsity}
\end{figure}

\subsection{Gauge conditioning near a critical point}

Placing a six-robot star at a saddle of a representative vector field and
driving the reference spoke radius $r_1$ toward zero, the regime in which
the pointer is least trustworthy, the condition number of the orientation
block is recorded for both the pointer gauge and the consensus gauge. The
pointer-gauge sensitivity grows as $1/r_1$ while the consensus-gauge
sensitivity stays bounded, because a vanishing spoke drops out of the
consensus sums instead of dominating them (Figure~\ref{fig:gauge}). The
consensus gauge pays for this with a single shared correction term across
the shape rows, a rank-one change that the Sherman-Morrison update absorbs,
so the improvement costs one rank of decoupling and nothing more.

% %% TODO: SYNC WITH LATEST DRAFT -- fig:gauge is an unfilled placeholder;
% generate real orientation-block condition number curves (pointer vs.
% consensus) vs. r_1, per Thesis_Execution_Plan.md Section 6, item 1.
\begin{figure}[t]
\centering
\framebox[0.9\linewidth]{\rule{0pt}{3.0cm}}
\caption{Orientation conditioning as the reference radius $r_1 \to 0$,
pointer gauge versus consensus gauge.}
\label{fig:gauge}
\end{figure}

\section{Discussion and Limitations}
\label{sec:cb_discussion}

The construction of this chapter has clear limits, and naming them
sharpens where it helps. The tree is a modeling choice that must match the
formation: a pairing chosen for an elongated front-rear group degrades if a
maneuver wants the pairing to swap. Conditioning compounds along the
cascade, so a single ill-conditioned layer propagates through the whole
chart, and each node's geometry must be kept healthy. Over-listing
coordinates buys no control authority: a planar formation has $2n$ degrees
of freedom however the variables are named (Theorem~\ref{thm:dim}), and the
extra directions a redundant parameterization appears to offer are
dependent, not free. Genuine task redundancy, when wanted, comes from the
navigation objective being lower-dimensional than the formation, and belongs
in a null-space projection at the control layer rather than in the chart
itself.

Several extensions follow naturally but remain future work. A dynamics
layer would carry cluster-space inertia and wrench mappings alongside the
present velocity-level kinematics. Spine-relative shape angles would let a
controller issue splay-relative commands at the price of a small triangular
coupling between layers. A typed field and controller interface would let
navigation primitives, such as those developed in
Chapters~\ref{ch:vector_field} and~\ref{ch:separatrix}, attach to the
generator as optional modules, keeping this chapter's kinematic core, which
is broadly reusable across multi-robot formation control, separate from
results that depend on a particular field or task.

\section{Conclusion}
\label{sec:cb_conclusion}

This chapter recast cluster-space kinematics as a compositional grammar and
gave a single-pass algorithm for the inverse Jacobian it produces. Three
results follow. A formation is a tree of base charts over children that are
robots or sub-clusters, and a direct dimension count
(Theorem~\ref{thm:dim}) shows the chart is complete and square for any such
tree, so the long-standing practice of deriving and sizing each formation
by hand is replaced by one generator. The full $2N\times2N$ inverse Jacobian
is then assembled in one depth-first pass using a per-node sensitivity
matrix, which removes the node-to-root rotation product that a prior
hierarchical construction recomputed at every level and lowers the cost to
$O(N)$ at bounded depth and $O(N\log N)$ for balanced trees. Finally, a
consensus orientation gauge distributes the formation frame across every
spoke, trading one rank of decoupling, recovered by a Sherman-Morrison
update, for a frame that survives a wandering or near-coincident reference
robot. This last result matters specifically for
Chapters~\ref{ch:vector_field} and~\ref{ch:separatrix}, whose control laws
deliberately drive a formation toward the low-signal neighborhood of a
critical point or boundary, exactly where a wandering reference is most
likely.

With this general kinematic substrate established, the remaining chapters
of this dissertation instantiate two specific formations from its atom
library, the SAS triangle for first-order field estimation
(Chapter~\ref{ch:vector_field}) and the pentagon-plus-center for
second-order field estimation (Chapter~\ref{ch:separatrix}), and ask, in
each case, not how the formation's own shape is maintained, which this
chapter has now answered in general, but what the formation can learn about
the field it occupies.

% ===================================================================
% CHAPTER 4: VECTOR FIELD PAPER
% ===================================================================
\chapter{First-Order Estimation: Critical Point Navigation}
\label{ch:vector_field}

\section{Related Work}
\label{sec:vf_related_work}

Multirobot systems exploit the spatial distribution of their members to
collect local samples of an environment and estimate features in their
vicinity. These estimates drive adaptive navigation policies that allow a
team to monitor and explore regions that are hazardous to humans, that vary
in time, or that are too large for a single platform to survey, and such
systems have been field-deployed for environmental monitoring, pollution
tracking, and search and rescue using ground vehicles, surface vessels, and
aerial vehicles. Many of these environments are governed by a spatially
varying velocity or force, and their transport and accumulation behavior is
organized by the field's critical points, the coordinates where the flow
vanishes. A surface convergence zone acts as a sink that gathers buoyant
debris, making it a high-probability search region after a loss at sea. A
ruptured subsea wellhead acts as a source whose origin must be reached
before responders can cap it, against a flow that drives an approaching
agent back outward. A mesoscale ocean eddy acts as a center that traps a
body of water and carries it across a basin for weeks, so monitoring
requires a sensor to orbit the drifting core at a fixed standoff rather than
hold a ground-fixed point against the rotation. In each case the operational
requirement is the coordinate of the feature itself, either to drive an
agent onto it or to hold a fixed radius around it, not merely a heading
toward it.

Adaptive navigation in scalar field environments is comparatively mature: a
range of strategies locate extrema, saddle points, ridges, trenches, and
fronts using gradient following, Hessian estimation, information consensus
filters, constrained nonsmooth distributed optimization, and bio-inspired
search. By contrast, only a few works address navigation of physically
sensed vector fields; much of the broader vector-field literature instead
constructs artificial fields for path-following and obstacle avoidance, a
separate problem from sensing a real flow. The work that does sense
physical fields falls into two camps. A substantial portion presumes prior
knowledge of the environment, unavailable in the unmapped settings that
motivate this chapter. The few online methods split into minimally actuated
drifters that ride and station-keep along ambient streamlines, and
PIM-triple-inspired teams that straddle a saddle and track a single stable
or unstable manifold from local velocity measurements. These either supply a
heading along the flow or track the manifold of a single assumed saddle;
none recovers the coordinate of a critical point together with its identity
across the full taxonomy of feature types. Estimating the locations of
critical points in vector fields remains an open problem for multirobot
systems, and is the problem this chapter solves.

Recovering those coordinates is harder than following a heading, for reasons
of both information and implementation. First, the location and
classification of a critical point are absent from any single velocity
sample and are carried instead by the spatial derivatives of the field. A
single platform must move, sample again, and difference its readings, which
conflates the spatial gradient with the field's change in time, since the
field evolves during transit, and can return features that are not
physically present. A spatially distributed team removes this temporal
aliasing by sampling at one instant, but recovering the derivatives is not
the same as recovering the feature. Magnitude-based methods reduce the flow
to a scalar and yield a descent direction but cannot identify the feature,
since the magnitude vanishes at every critical point regardless of type, and
estimating the two velocity components as independent scalar fields still
does not estimate the location. The coordinate and the identity are
properties of the assembled Jacobian: the location is the point where the
coupled field vanishes, recovered by inverting the matrix, and the type is
fixed by its eigenvalues. Neither follows from treating the components as
separate gradient-following problems. Second, even with the right
information in hand, the local field must be reconstructed on physical
hardware from noisy distributed measurements. The conditioning of this
reconstruction degrades as the formation approaches a collinear geometry, so
the sampling shape must be held by a formation controller operating under
real actuator dynamics and inter-robot communication latency, precisely the
cluster-space machinery generalized in Chapter~\ref{ch:cluster_builder}. The
accuracy of the recovered features is therefore bounded by the mechatronics
of the cluster as much as by the estimator itself.

This chapter presents a distributed cooperative estimation strategy that
recovers the local linear structure of the field from instantaneous vector
measurements. Using a minimum of three robots, the cluster estimates the
local Jacobian in a single time step. From this local estimate it
extrapolates a global feature of the field, treating critical point
navigation as a local-to-global inference: inverting the Jacobian predicts
the location of the nearby critical point that organizes the surrounding
flow, and its eigenvalues fix the identity, with no global map required. The
same estimate supports two control laws, a proportional law that drives the
cluster centroid onto the critical point, and an orbital law that holds the
centroid at a commanded standoff radius around it. This chapter's
contributions are a distributed algorithm that estimates critical point
locations and types from instantaneous measurements alone; a proof that
three robots are the minimum number required for this estimation; control
strategies that use these estimates for navigation to and around critical
points across six vector field types; and validation in 169 hardware
experiments on vortex and saddle fields that demonstrate repeatable critical
point estimation, attraction, and orbit control.

\section{Vector Field Environments and Critical Points}

In its most general form, a vector field assigns a vector
$\mathbf{v}(\mathbf{p}) \in \mathbb{R}^m$ to each point
$\mathbf{p}\in\mathbb{R}^n$ within a given domain, where $m\geq2$. Each
vector carries a magnitude and direction, making vector fields well suited
to represent directional physical quantities such as fluid velocity,
electromagnetic forces, or gravitational forces. For the planar fields
analyzed in this chapter, this takes the specific form of a two-dimensional
vector $\mathbf{v}(\mathbf{p}) = (u(\mathbf{p}),v(\mathbf{p}))$ assigned to
each coordinate $\mathbf{p}=(x,y)$.

Unlike scalar fields where $m=1$, vector fields can exhibit features such as
critical points $\mathbf{p}^*$ where the field magnitude vanishes:
\begin{equation}
\mathbf{v}(\mathbf{p}^*) = \mathbf{0}.
\end{equation}
Critical points model phenomena such as vortices arising from tidal flows
over submarine caves, or plumes when searching for pollution sources. The
ability for autonomous vehicles to estimate the locations of these critical
points is necessary for many missions: underwater robots monitoring an
oceanic vortex or atmospheric drones tracking a developing cyclone must
estimate the field center to maintain stable orbits at specified radii or
compensate for drift. Without an explicit estimate of the critical point
location, these tasks reduce to reactive strategies that provide directional
guidance but cannot achieve precise standoff control.

Near a critical point, the field is well approximated by its linearization,
and the local topology is determined by the Jacobian matrix
$\mathbf{J}=\nabla\mathbf{v}|_{\mathbf{p}}$. The eigenvalues of $\mathbf{J}$
evaluated at $\mathbf{p}^*$ determine the topological type. For
$2\times2$ systems with nonzero determinant, six standard critical point
types arise from the eigenvalue structure, shown in
Figure~\ref{fig:allfield} and classified in Table~\ref{tab:field_summary}.
This classification holds rigorously for hyperbolic critical points; the
center is the exception, since a nonlinear perturbation of a linear center
can produce a spiral (Hartman-Grobman does not apply at $\alpha_{\text{eig}}=0$),
and under measurement noise the estimated real part of a center's
eigenvalue hovers near zero, making center-versus-weak-spiral discrimination
inherently ill-posed. Eigenvalue analysis can also identify degenerate
scenarios where the determinant is zero, corresponding to non-hyperbolic
critical points or regions of near-uniform flow; these are detectable
through the estimation framework of the next section but fall outside the
scope of this chapter.

\begin{figure}[!t]
\centering
\framebox[0.9\linewidth]{\rule{0pt}{3.0cm}}
\caption{Six linear vector field types with critical points.}
\label{fig:allfield}
\end{figure}

\begin{table}[!t]
\centering
\caption{Classification of vector field critical points based on Jacobian
eigenvalue structure ($\lambda$ = eigenvalue; complex eigenvalues
$\alpha_{\text{eig}}\pm i\omega$ where $\alpha_{\text{eig}}$ is the real
part, $\omega$ the imaginary part).}
\label{tab:field_summary}
\begin{tabular}{lll}
\toprule
\textbf{Field Type} & \textbf{Critical Point} & \textbf{Eigenvalues} \\
\midrule
Vortex & Center & $\lambda=\pm i\omega$ \\
Sinking Vortex & Stable Spiral & $\lambda=\alpha_{\text{eig}}\pm i\omega$, $\alpha_{\text{eig}}<0$ \\
Sink & Stable Node & $\lambda_1,\lambda_2<0$ \\
Source & Unstable Node & $\lambda_1,\lambda_2>0$ \\
Spewing Vortex & Unstable Spiral & $\lambda=\alpha_{\text{eig}}\pm i\omega$, $\alpha_{\text{eig}}>0$ \\
Saddle & Saddle Point & $\lambda_1<0<\lambda_2$ \\
\bottomrule
\end{tabular}
\end{table}

\section{Distributed Critical Point Estimation}
\label{sec:vf_estimation}

Following Chapter~\ref{ch:cluster_builder}, we refer to the group of robots
as a cluster: a formation treated as a single virtual rigid body whose
centroid, orientation, and internal geometry can be independently specified
and controlled, here the SAS-3 chart of
Section~\ref{sec:atom_sas}. Given a cluster of three robots sampling a
planar vector field with robot positions $\mathbf{p}_i=(x_i,y_i)$ and
corresponding vector field measurements $(u_i,v_i)$ for $i\in\{1,2,3\}$, the
vector field components across the cluster's spatial extent are approximated
by linear planar models
\begin{equation}
U(x,y)=ax+by+c, \qquad V(x,y)=dx+ey+f.
\end{equation}
The coefficients are determined by solving the linear systems
$\mathbf{A}\boldsymbol{\theta}_u=\mathbf{u}$, $\mathbf{A}\boldsymbol{\theta}_v=\mathbf{v}$,
where
\begin{equation}
\mathbf{A} = \begin{bmatrix} x_1 & y_1 & 1 \\ x_2 & y_2 & 1 \\ x_3 & y_3 & 1 \end{bmatrix},
\quad
\boldsymbol{\theta}_u = \begin{bmatrix} a\\b\\c \end{bmatrix},
\quad
\boldsymbol{\theta}_v = \begin{bmatrix} d\\e\\f \end{bmatrix},
\end{equation}
and $\mathbf{u}=[u_1,u_2,u_3]^{\mathsf T}$, $\mathbf{v}=[v_1,v_2,v_3]^{\mathsf T}$
are the collected field measurements. The critical point
$\mathbf{p}^*=(x^*,y^*)$ satisfies $U(x^*,y^*)=0$ and $V(x^*,y^*)=0$, which
yields $ax^*+by^*+c=0$, $dx^*+ey^*+f=0$, or in matrix form
\begin{equation}
\mathbf{J}\mathbf{p}^* = -\mathbf{h}, \qquad
\mathbf{J} = \begin{bmatrix} a & b \\ d & e \end{bmatrix}, \qquad
\mathbf{h} = \begin{bmatrix} c \\ f \end{bmatrix}.
\label{eq:vf_jacobian_offset}
\end{equation}
Thus, provided $\det(\mathbf{J})\neq0$,
\begin{equation}
\mathbf{p}^* = -\mathbf{J}^{-1}\mathbf{h}.
\label{eq:vf_pstar}
\end{equation}
The eigenvalues of $\mathbf{J}$ then identify the type of the nearby
critical point via Table~\ref{tab:field_summary}.

\section{Minimum Robot Requirements}
\label{sec:vf_minimality}

The two plane-fit systems each contain three coefficients. Each robot
contributes one equation to each system: robot $i$ provides
$u_i=ax_i+by_i+c$ for the $U$-component and $v_i=dx_i+ey_i+f$ for the
$V$-component. To uniquely determine three unknowns requires at least three
equations, so three robots are the minimum required for critical point
estimation. With fewer than three robots the system is underdetermined,
$\mathbf{A}$ has fewer rows than columns, admitting infinitely many
coefficient solutions and preventing unique identification of the critical
point location. With more than three robots the system is overdetermined
and solved via least-squares minimization. Three non-collinear robots are
also sufficient: $\det(\mathbf{A})$ equals twice the signed area of the
formation triangle, so $\mathbf{A}$ is invertible exactly when the robots
are not collinear, connecting this minimality argument directly to the
conditioning discussion of Section~\ref{sec:vf_error_analysis}.

\section{Adaptive Navigation Control Laws}
\label{sec:control_primitives}

Two control primitives drive the cluster centroid
$\mathbf{p}_c=\tfrac13\sum_{i=1}^3\mathbf{p}_i$ toward or around the
estimated critical point $\mathbf{p}^*$. The mapping from cluster-level
velocity commands to individual robot velocities is handled by the SAS-3
cluster-space controller of Chapter~\ref{ch:cluster_builder}.

\subsection{Critical Point Attraction}

To navigate the cluster center $\mathbf{p}_c$ to the critical point, a
proportional control law is applied:
\begin{equation}
\dot{\mathbf{p}}_c = k(\mathbf{p}^*-\mathbf{p}_c),
\label{eq:vf_attraction}
\end{equation}
where $k>0$ is the proportional gain. This linear error dynamic guarantees
global exponential convergence to the estimated critical point with time
constant $\tau=1/k$, provided the estimate is stationary and in the absence
of actuation limits; on the linear fields of Section~\ref{sec:vf_simulation}
the estimate is exact and hence automatically stationary in the noise-free
case.

\subsection{Orbital Navigation}

To maintain a circular orbit of radius $r_d$ around $\mathbf{p}^*$, define
the radial vector $\mathbf{r}=\mathbf{p}^*-\mathbf{p}_c$ with magnitude
$r=\lVert\mathbf{r}\rVert$. The unit radial vector is
$\hat{\mathbf{r}}=\mathbf{r}/r$, and the tangent vector
$\hat{\mathbf{r}}_\perp$ is obtained via a $-90°$ rotation. The orbital
velocity command is
\begin{equation}
\dot{\mathbf{p}}_c = k_t\hat{\mathbf{r}}_\perp + k_r(r-r_d)\hat{\mathbf{r}}.
\label{eq:vf_orbital}
\end{equation}
Here $k_t$ is the commanded orbital (tangential) speed, whose sign fixes
direction, and $k_r>0$ drives radial convergence. This control law decouples
radial and angular dynamics: the radial error $e_r=r-r_d$ follows
$\dot e_r=-k_re_r$, ensuring exponential convergence to $r_d$, and as
$r\to r_d$ the angular velocity stabilizes to $\omega=k_t/r_d$, producing
steady circular motion. Together, Equations~\eqref{eq:vf_attraction}
and~\eqref{eq:vf_orbital} enable a two-phase mission: approach the critical
point, then transition to orbital station-keeping.

\section{Robot and Cluster Dynamics}
\label{sec:vf_dynamics}

The three-layer control architecture of
Section~\ref{sec:testbed_architecture} applies unchanged: an adaptive
navigation layer implementing Sections~\ref{sec:vf_estimation}
through~\ref{sec:control_primitives} generates desired centroid velocities
$\dot{\mathbf{p}}_c$; the SAS-3 cluster-space controller of
Chapter~\ref{ch:cluster_builder} maintains formation geometry at
$p=q=0.33$~m (this chapter's equilateral triangle, chosen to minimize the
condition number of $\mathbf{A}$; see Section~\ref{sec:vf_error_analysis})
while translating cluster velocity commands into individual robot
velocities via the inverse kinematic Jacobian $\mathbf{J}_c$ of
Appendix~\ref{app:cluster_kinematics}; and each robot's own controller
executes the resulting velocity command. At each 10~Hz control cycle, every
robot reports its position and vector field measurement, the cluster-space
kinematics are computed, and velocity commands are mapped in both
directions.

Each robot is modeled as a single omnidirectional point mass with a
first-order response,
\begin{equation}
\dot{\mathbf{v}}(t) = -\frac{1}{\tau}\mathbf{v}(t) + \frac{1}{\tau}\mathbf{v}_{\text{des}}(t),
\label{eq:vf_continuous}
\end{equation}
where $\tau$ is the actuator time constant, measured at 0.3~s on the Decabot
hardware of Chapter~\ref{ch:testbed}. Robots are limited to 0.3~m/s maximum
speed.
% %% TODO: SYNC WITH LATEST DRAFT -- the source text states a 0.05 m/s
% stiction floor, but the simulation that generated Tables II-III actually
% ran with the Omnibot code default of 0.025 m/s (AUDIT_REPORT_4A.md item
% E2). This is a user decision (correct the text to 0.025, or regenerate
% the tables at 0.05), not resolved here per Thesis_Execution_Plan.md
% Section 6, item 5 / Section 10, question 1.
For digital implementation at control period $\Delta t$, discretization
yields
\begin{equation}
\mathbf{v}[k+1] = \alpha_{\text{mom}}\mathbf{v}[k] + (1-\alpha_{\text{mom}})\mathbf{v}_{\text{des}}[k],
\qquad \alpha_{\text{mom}} = e^{-\Delta t/\tau},
\label{eq:vf_momentum}
\end{equation}
so that $\alpha_{\text{mom}}\approx0.717$ at the 10~Hz control rate used
throughout this dissertation.

\section{Control Performance in Simulation}
\label{sec:vf_simulation}

A Python-based simulation implementing the architecture of
Section~\ref{sec:vf_dynamics} was developed to verify the estimation and
control primitives above. For all simulation trials the cluster consisted
of three robots in an equilateral triangular formation with a side length
of 0.33~m, matching the hardware configuration of
Section~\ref{sec:vf_hardware}. Six noise-free analytical fields, one for
each canonical critical point type of Table~\ref{tab:field_summary}, were
used for trials (equations in Appendix~\ref{app:vector_fields}). After
verifying performance in noise-free environments, simulations were also run
on two fields reconstructed from robot-sensed measurements collected on the
hardware testbed (vortex and saddle), incorporating realistic sensor noise.

\subsection{Convergence to Critical Points}

For each of the eight fields, the cluster was initialized at a random
position drawn uniformly from a 1~m~$\times$~1~m box centered on the true
critical point, and the simulation was run for 100 time steps at 10~Hz. A
trial was counted as convergent if the trajectory remained bounded and
terminated within the stiction-limited neighborhood of the critical point;
final-position statistics quantify that neighborhood. This was repeated
1000 times per field, and the final centroid position of each trial was
used to determine systematic bias (mean error) and precision (standard
deviation).

The cluster converged in 100\% of the 1000 trials for every field tested.
On the six analytical fields, systematic bias remained below 0.001~m and
precision ranged from 0.014 to 0.016~m. On the two reconstructed fields,
systematic bias was 0.028~m for the vortex and 0.005~m for the saddle, with
precision near 0.016~m in both cases (Table~\ref{tab:simulation_estimation_results}).
The reconstructed-field trials isolate systematic reconstruction bias rather
than stochastic sensing noise: the NN-reconstructed fields add a frozen,
smooth field warp with almost no per-trial scatter, so reconstructed
precision matches noise-free precision to the third decimal, while per-trial
variance in both cases is dominated by robot dynamics scatter rather than
sensor stochasticity.

\begin{table}[!t]
\footnotesize
\centering
\caption{Simulation results across eight field types (six analytical plus
two reconstructed). All fields achieved 100\% convergence over 1000 trials.}
\label{tab:simulation_estimation_results}
\begin{tabular}{lccc}
\toprule
\textbf{Environment} & \textbf{Avg.\ Rest Point (m)} & \textbf{Bias (m)} & \textbf{$\sigma$ (m)} \\
\midrule
Sink & (0.0006, $-$0.0001) & 0.0006 & 0.0156 \\
Source & ($-$0.0005, $-$0.0001) & 0.0006 & 0.0140 \\
Vortex & ($-$0.0002, $-$0.0000) & 0.0002 & 0.0163 \\
Sinking Vortex & ($-$0.0002, $-$0.0004) & 0.0004 & 0.0145 \\
Spewing Vortex & (0.0001, 0.0006) & 0.0006 & 0.0157 \\
Saddle & (0.0004, $-$0.0001) & 0.0004 & 0.0157 \\
Vortex (recon.) & (0.0206, $-$0.0187) & 0.0278 & 0.0162 \\
Saddle (recon.) & ($-$0.0048, $-$0.0007) & 0.0048 & 0.0157 \\
\bottomrule
\end{tabular}
\end{table}

\subsection{Orbital Control}

Orbital control was evaluated on the same eight fields at six commanded
radii. For each radius, the cluster was positioned at the desired radius to
initialize the run, then the orbital controller
(Equation~\eqref{eq:vf_orbital}) was used for 600 timesteps; mean radial
error and standard deviation were computed after discarding the first 100
timesteps (10~s) as a settling transient.

\begin{table}[!t]
\centering
\caption{Simulated orbital control across eight field types.}
\label{tab:simulation_orbital_results}
\footnotesize
\begin{tabular}{cccc}
\toprule
\textbf{Radius (m)} & \textbf{Noise-free (6 fields)} & \textbf{Vortex (recon.)} & \textbf{Saddle (recon.)} \\
\midrule
0.01 & 0.146 $\pm$ 3e-06 & 0.150 $\pm$ 0.015 & 0.146 $\pm$ 0.012 \\
0.10 & 0.109 $\pm$ 7e-06 & 0.109 $\pm$ 0.013 & 0.098 $\pm$ 0.015 \\
0.20 & 0.081 $\pm$ 1e-05 & 0.104 $\pm$ 0.014 & 0.072 $\pm$ 0.024 \\
0.30 & 0.063 $\pm$ 1e-05 & 0.094 $\pm$ 0.018 & 0.059 $\pm$ 0.042 \\
0.40 & 0.050 $\pm$ 2e-05 & 0.067 $\pm$ 0.026 & 0.035 $\pm$ 0.054 \\
0.50 & 0.042 $\pm$ 2e-05 & 0.027 $\pm$ 0.038 & $-$0.003 $\pm$ 0.061 \\
\bottomrule
\end{tabular}
\end{table}

On the six noise-free fields, all six produced identical performance at
every radius (standard deviations below $2\times10^{-5}$~m), so only one
column is shown to represent all six. This identity is not coincidental: on
exactly linear noise-free fields the three-point fit recovers
$\mathbf{p}^*$ exactly at every step, so the controller input, and hence
the trajectory, is independent of which linear field is used; the six
noise-free trajectories are numerically identical to within $10^{-9}$~m,
and the residual standard deviation is steady-state radius ripple, not
field variation. Mean radial errors had a slight positive bias, decreasing
with increasing radius from 0.146~m at $r_d=0.01$~m to 0.042~m at
$r_d=0.50$~m. The same pattern held on the reconstructed vortex (0.150 to
0.027~m) and, with the trend continuing to a small negative bias at large
radius, the reconstructed saddle (0.146 to $-0.003$~m). In both
reconstructed fields, error standard deviations increased with commanded
radius, ranging from 0.012 to 0.061~m.

\section{Hardware Validation}
\label{sec:vf_hardware}

\subsection{Hardware Testbed and Field Representation}

Validation used the Decabot testbed of Chapter~\ref{ch:testbed}, three
omnidirectional rovers operating within a 1.6~m~$\times$~1.6~m workspace,
sensing an HSV-encoded printed vector field map and localized by an
eight-camera OptiTrack system. Sensor calibration achieved hue RMSE of
0.17~rad (9.7°, $R^2=0.991$) and saturation RMSE of 4.4\% ($R^2=0.951$)
across the operational range, consistent with the calibration methodology
of Section~\ref{sec:testbed_fields}.
% %% TODO: SYNC WITH LATEST DRAFT -- Fig. 6 in the source paper
% (testbed_with_four.png) appears to show four Decabots while captioned as
% three; per Thesis_Execution_Plan.md Section 6, item 6, this needs a
% correct three-robot photo before the figure is finalized here.
\begin{figure}[!t]
\centering
\framebox[0.9\linewidth]{\rule{0pt}{3.0cm}}
\caption{Complete experimental testbed with 1.6~m $\times$ 1.6~m printed
vector field map, three Decabot rovers, and OptiTrack system.}
\label{fig:fullbed}
\end{figure}

Vector fields are physically realized as printed floor maps using HSV color
encoding, hue (0 to $2\pi$) representing flow direction and saturation
(0 to 1) encoding magnitude, remapped in software so that saturation 1.0
corresponds to the maximum robot velocity of 0.3~m/s. This method of
representing vector fields has imperfections in both printing and sensing;
the reconstructed vortex and saddle fields used in
Section~\ref{sec:vf_simulation} were built from exactly this sensed
reconstruction error, which also fit the interpolation models used for the
realistic noise condition in simulation.

\subsection{Validation Test Plan}

Hardware validation focused on two representative field types spanning the
key challenges of vector field navigation: rotational dynamics (vortex) and
hyperbolic instability (saddle). The saddle is the only canonical field
type that contains separatrices, which act as Lagrangian coherent
structures in steady flows, foreshadowing Chapter~\ref{ch:separatrix}'s
direct study of such structures. For each field, 8 starting locations were
selected around the perimeter of the workspace and the attraction
controller (Equation~\eqref{eq:vf_attraction}) was run 10 times from each
location using the three-robot equilateral triangle formation with a
0.33~m side length, yielding 80 trials per field; 2 saddle trials and 1
vortex trial were lost to file management errors, giving 157 usable
convergence trials. Each field was also used to test the orbiting
controller (Equation~\eqref{eq:vf_orbital}) at 6 starting radii matching
the simulation distances, each run for 3000 cycles at 10~Hz, more than
three complete orbits.

\subsection{Attraction Controller Results}

Excluding lost trials, 79 vortex and 78 saddle trials were analyzed. The
cluster converged to the critical point in 100\% of the 157 hardware
trials. In the vortex field, systematic bias was 0.012~m with precision of
0.009~m; in the saddle field, systematic bias was 0.005~m with precision of
0.014~m (Table~\ref{tab:estimation_results}). The narrow standard deviation
bounds across the 10 repetitions per starting position confirm that the
residual scatter reflects sensor noise rather than uncontrolled
experimental variation.

\begin{table}[!t]
\footnotesize
\centering
\caption{Hardware experimental results. Saddle: 78 runs. Vortex: 79 runs.
Bias measures mean offset from the true critical point; precision
($\sigma$) measures repeatability.}
\label{tab:estimation_results}
\begin{tabular}{lccc}
\toprule
\textbf{Environment} & \textbf{Avg.\ Rest Point (m)} & \textbf{Bias (m)} & \textbf{$\sigma$ (m)} \\
\midrule
Saddle (printed) & (0.004, $-$0.001) & 0.005 & 0.014 \\
Vortex (printed) & ($-$0.004, 0.012) & 0.012 & 0.009 \\
\bottomrule
\end{tabular}
\end{table}

\subsection{Orbital Control Results}

For each radius, the orbital controller was run for 3000 cycles at 10~Hz
and the mean radial error and standard deviation were recorded over the
full run.
% %% TODO: SYNC WITH LATEST DRAFT -- the simulation orbital statistics
% (Section~\ref{sec:vf_simulation}) discard a 100-step settling transient
% before averaging; these hardware statistics do not (AUDIT_REPORT_4A.md
% item E4), so Table~\ref{tab:simulation_orbital_results} and
% Table~\ref{tab:orbital_radius_comparison} are not computed over
% matched windows. Disclosed here per audit recommendation; not
% reconciled, since doing so would change reported numbers.
In the vortex field, mean radial errors decreased from 0.246~m at a
commanded radius of 0.01~m to $-0.035$~m at 0.50~m, with standard
deviations ranging from 0.012 to 0.053~m. In the saddle field, mean radial
errors ranged from 0.114~m to 0.060~m over the same commanded radii, with
standard deviations increasing from 0.011 to 0.095~m
(Table~\ref{tab:orbital_radius_comparison}).

\begin{table}[!t]
\centering
\caption{Hardware orbital control results.}
\label{tab:orbital_radius_comparison}
\begin{tabular}{ccc}
\toprule
\textbf{Radius (m)} & \textbf{Vortex} & \textbf{Saddle} \\
\midrule
0.01 & 0.246 $\pm$ 0.053 & 0.114 $\pm$ 0.011 \\
0.10 & 0.190 $\pm$ 0.012 & 0.072 $\pm$ 0.018 \\
0.20 & 0.143 $\pm$ 0.019 & 0.058 $\pm$ 0.037 \\
0.30 & 0.087 $\pm$ 0.015 & 0.089 $\pm$ 0.059 \\
0.40 & 0.032 $\pm$ 0.020 & 0.090 $\pm$ 0.079 \\
0.50 & $-$0.035 $\pm$ 0.023 & 0.060 $\pm$ 0.095 \\
\bottomrule
\end{tabular}
\end{table}

\section{Discussion}

\subsection{Error Analysis}
\label{sec:vf_error_analysis}

Estimation and control errors arise from four sources: robot dynamics,
sensor noise, cluster geometry, and field conditioning. Perfect convergence
does not occur in simulation; however, repeating experiments with no
stiction, no max velocity, and $\alpha_{\text{mom}}=0$ (instantaneous
velocity response) confirms convergence to the controller's $10^{-6}$~m
command deadband, with the critical point estimate itself exact to machine
precision. This confirms the residual variance is due to robot dynamics,
not the estimation method.

Positive radial bias at small orbital radii is also a result of robot
dynamics. At small radii, the commanded angular rate $k_t/r_d$ is high
relative to the radial correction bandwidth $k_r$, and the centripetal
demand ($k_t^2/r$) pushes the cluster outward faster than the radial
controller can correct; as the commanded radius increases, this effect
diminishes. More aggressive radial gains would reduce this error, at the
cost of increased sensitivity to noise in the critical point estimate.

To avoid errors arising from robots becoming collinear (condition number of
$\mathbf{A}$ approaching infinity), an equilateral triangle was chosen as
the cluster shape: among 200{,}000 random formations of the same size
centered on the sampled region, none beat the equilateral triangle's
$\kappa(\mathbf{A})=7.42$. This optimality is not translation-invariant
(the same formation centered elsewhere has a different condition number),
so it is stated for a formation centered on the sampled region, as here.

Further errors arise from poor conditioning of the estimated field
Jacobian $\kappa(\mathbf{J})$, directly affecting the critical point
estimate. The condition number grows when the Jacobian eigenvalues have
very different magnitudes, as occurs near highly asymmetric nodes or
saddle points with disparate stable and unstable manifold strengths, and
diverges as $\det(\mathbf{J})\to0$, corresponding to non-hyperbolic
critical points or near-uniform flow. The negative bias at the largest
commanded radii traces to degradation of the reconstructed field near the
edge of the mapped region: there the estimated critical point distance is
systematically inflated, so the radial controller pulls the cluster inside
the commanded radius, the same range-dependent error amplification
captured by $\kappa(\mathbf{J})$, since far from the critical point the
sampled patch approaches uniform flow and small field errors produce large
estimate shifts. Across all four error sources, estimation errors shift the
equilibrium location of both controllers but preserve the stability
predicted by the linear error dynamics.

\subsection{Comparison with Other Primitives}

The modular architecture of Chapters~\ref{ch:testbed}
and~\ref{ch:cluster_builder} allows navigation controllers to be
interchanged while cluster and robot controller layers remain fixed, so
performance differences reflect algorithmic rather than implementation
variations. Comparing the orbital controller above against two
alternatives in the vortex field, the vector-sum approach (which commands
motion in the direction of the average perceived flow) and the
vector-to-scalar primitive (modified to move tangentially to the direction
of steepest descent), both alternatives exhibit outward spiral drift: both
derive velocity commands from direction-only guidance, with no mechanism to
detect or correct radial error with respect to a specific point, so
centrifugal drift goes uncorrected. In the saddle field the comparison is
more decisive, since the vectors do not all converge on the critical point:
the vector-sum controller cannot produce an orbital trajectory at all,
while the vector-to-scalar technique still produces drift.

\subsection{Limitations}

Several limitations constrain this chapter's implementation and validation.
The approach is restricted to planar vector fields, though many real-world
applications involve three-dimensional flows; extension to 3D would require
at minimum four non-coplanar robots for Jacobian estimation. The testbed
validates performance only on static fields where critical points remain
fixed during operation. Robot team sizes were limited to three; while the
theoretical framework supports arbitrary team sizes through the
overdetermined system approach, larger formations present practical
implementation challenges including communication bandwidth, computational
complexity, and formation control stability. (A supplementary, unpublished
four-robot hardware campaign appears in Appendix~\ref{app:four_robot}.)
Finally, this chapter assumes a sampling rate sufficiently high that the
system approximates continuous control; robustness to lower update rates
and communication dropouts was not tested.

\subsection{Future Work}
\label{sec:vf_future_work}

The estimation framework uses only instantaneous measurements and carries
no state between control cycles. This memoryless property means the
estimator does not assume field stationarity, suggesting applicability to
time-varying environments where field structure evolves during operation.
Beyond critical point detection and orbital control, vector fields contain
other topological features of practical interest. Separatrices connected to
saddle points define boundaries between distinct flow regimes and act as
Lagrangian coherent structures in steady flows; heteroclinic orbits
connecting different critical points through their manifolds offer transit
routes through complex fields. The Jacobian estimation framework developed
in this chapter provides the foundation for detecting such features through
analysis of eigenvector alignment between neighboring critical points,
foundation that Chapter~\ref{ch:separatrix} builds on directly by
extending this chapter's linear, first-order estimate to a quadratic,
second-order one capable of tracking such a boundary rather than only a
single isolated point.

\section{Conclusion}

This chapter presented a distributed framework for multirobot navigation in
vector fields containing critical points. A Jacobian-based approach uses
instantaneous measurements from three robots to identify critical point
locations and types without prior field knowledge. Simulation across six
canonical field types confirmed convergence of these control primitives. In
157 hardware trials, the system achieved a 100\% success rate with position
errors of 0.012~m or less. Orbital control experiments maintained commanded
radii within 0.25~m error at small radii and 0.1~m at larger radii.
Chapter~\ref{ch:separatrix} takes up this chapter's stated limitation to
isolated critical points on static fields directly, extending the
estimation order from first (Jacobian) to second (Hessian) and the
formation from three robots to six, so that a team can track an extended,
potentially moving boundary rather than only a single point.

% ===================================================================
% CHAPTER 5: SEPARATRIX / OKUBO-WEISS PAPER
% ===================================================================
\chapter{Second-Order Estimation: Separatrix and Okubo-Weiss Tracking}
\label{ch:separatrix}

\section{Introduction}
\label{sec:sep_intro}

Objective Eulerian coherent structures (OECSs) are short-term,
frame-invariant transport barriers that organize instantaneous material
transport in general dynamical systems. They can be viewed as the
instantaneous counterparts of stable and unstable manifolds, extracted from
the rate-of-strain field at a single instant rather than from trajectories
integrated over a finite horizon. Identifying OECSs is useful for
applications that demand short-term or real-time forecasting, including
oceanography and weather prediction. This chapter presents a collaborative
robotic control strategy that tracks the short-term attracting and
repelling structures of a flow, including ocean flows, without
prepositioning the robots and without requiring global information about
the dynamics; it relies instead on local sensing, prediction, and
correction. For the two-dimensional incompressible flows considered, the
sign of the determinant of the flow Jacobian, equivalently the Okubo-Weiss
field, partitions the flow into vorticity-dominated and strain-dominated
regions: its zero level set forms the boundary between them, and its
trench marks the strongest hyperbolicity, where the hyperbolic OECS
resides. Two complementary tracking modes are developed. In the first, the
team follows the Okubo-Weiss boundary, the zero level set that divides the
two regimes. In the second, the team navigates the Jacobian-determinant
field itself. Rather than steepest descent or ascent, a modified Newton
step, which reweights the search directions so the team snaps onto the
trench more rapidly, drives the team along the structure, through a saddle
point of this scalar field, and down the opposing branch. The strategy is
implemented with a team of six robots and validated on a time-dependent
model of a wind-driven double-gyre flow and on actual ocean-flow data.

Multirobot systems exploit spatial distribution across a team to sample and
act on an environment more effectively than a single platform can, a
capability with broad application to formation assembly, coverage, and
distributed sensing. Deployed in geophysical flows such as ocean currents
and atmospheric jets, these teams often need to locate and track
dynamically meaningful boundaries rather than isolated critical points, the
subject of Chapter~\ref{ch:vector_field}. Coordinated multi-vehicle
sampling of dynamic ocean features is itself an active area of marine
robotics, from drifter-referenced AUV surveys of advecting patches to
heterogeneous ASV/AUV teams tracking salinity fronts and glider fleets
adaptively sampling upwelling events. Two boundaries appear prominently in
steady two-dimensional flows: the separatrix, which contains both the
stable and unstable manifold connecting to saddle points and partitions the
domain into two distinct basins; and the Okubo-Weiss boundary, the locus of
points where the local flow transitions from rotation-dominated to
strain-dominated behavior, closely related to the physical parameter used
to detect and classify oceanic mesoscale eddies. Both structures are
relevant to dispersion modeling, pollutant containment, and the deployment
of sensor networks in coastal and atmospheric environments, and both can,
in principle, be inferred from instantaneous field measurements alone.

A line of prior work has equipped robot teams to track a coherent structure
directly in the flow rather than reconstruct it offline. Michini et al.
straddle a presumed stable or unstable manifold with a three-robot,
PIM-triple-inspired formation, using only local velocity measurements to
keep one robot near the manifold and the other two on opposite sides of it
[34]; the same strategy was extended to an $N$-robot team for finer spatial
sampling of the same manifold [13] and validated on physical micro-ASVs in a
flow tank [43]. Knizhnik et al. couple a related manifold-tracking law to
vehicle-level flow control in gyre-like marine environments [42]. Kularatne
and Hsieh take a different approach, computing the local finite-time
Lyapunov exponent field on board from the team's own accumulated velocity
measurements, which recovers the structure's type explicitly rather than
assuming it [12]. Every member of this family, however, fixes the
structure's identity before tracking begins: [34], [13], and [43] require a
saddle-straddling line segment seeded around an already-assumed manifold,
and their controller commits in advance to tracking either the stable or
the unstable branch, since the team's estimate of the manifold's location
is read off as the point of maximum measured speed along the straddle
segment if the unstable manifold is being tracked, or of minimum speed if
the stable manifold is being tracked. Neither the controller nor its
stability proof addresses what happens when the tracked branch reaches a
saddle, where this choice becomes ambiguous, since the straddle segment's
validity proof assumes it crosses the manifold exactly once and stops
there. Kularatne and Hsieh's onboard FTLE estimate avoids fixing the
structure's type in advance, but only by accumulating a velocity time
history over a finite integration horizon before an estimate is well
formed, an approach discussed further in
Section~\ref{sec:separatrix}. This chapter's estimator and controllers
require neither a seeded manifold nor an integration horizon: from an
arbitrary six-robot configuration, they infer a structure's location, its
type, and, at a crossing, which branch to continue onto, from the current
instantaneous measurement alone.

Chapter~\ref{ch:vector_field} estimated a single isolated critical point and
classified it from the eigenvalues of a linear Jacobian fit, using a
three-robot formation and a proportional or orbital control law. That is a
point-estimation problem. The work presented here is a field-structure
problem: it fits a quadratic model of the vector field, derives the
curvature of $\det(\mathbf{J})$ from that fit, and uses the curvature to
track one-dimensional manifolds and boundaries rather than converge to a
point.

This chapter presents the following contributions:
\begin{enumerate}
\item A distributed quadratic estimator that uses simultaneous measurements
from six omnidirectional robots to recover the Jacobian $\mathbf{J}$, the
determinant $\det(\mathbf{J})$, its gradient $\nabla\det(\mathbf{J})$, and
its Hessian $\mathbf{H}_D$ at the cluster centroid, with no prior knowledge
of the field. We prove six robots are the exact minimum for the coupled
two-component case, one more than the five-plus-center formation sufficient
for a single scalar field [3], and derive $\det(\mathbf{J})$, its gradient,
and its Hessian analytically by differentiating through the jointly fitted
vector coefficients, rather than by fitting a scalar surrogate directly.
\item Two adaptive controllers built on a modified Newton step: the
ordinary step $-\mathbf{H}_D^{-1}\nabla D$ is decomposed in the eigenbasis
of $\mathbf{H}_D$, each eigen-direction is independently saturated, and, in
the separatrix tracker, each signed eigenvalue is replaced by its magnitude
so the team slides along a trench even where the local curvature is
negative. (a) A hybrid three-mode separatrix tracker (Logic C) that uses
this step to slide along the stable manifold of the double-gyre saddle; (b)
an Okubo-Weiss boundary tracker that applies a Newton-step
contour-following law to drive $\det(\mathbf{J})\to0$ while drifting
tangentially along the boundary.
\item A stability and network-traversal proof for the separatrix tracker:
the trench a controller reaches continues through every saddle it crosses,
with single-saddle capture recovered only as a tightened-threshold special
case. Unlike prior manifold trackers, which fix the tracked branch's
identity before the run begins [34, 13], this controller commits to no
manifold identity in advance, so the branch to continue onto at a crossing
is read from the same instantaneous fit that located the crossing, rather
than assumed ahead of time.
\item Monte Carlo simulation validation of both controllers under a fully
specified noise model, with success-rate maps over a two-dimensional sweep
of measurement and position noise standard deviations, and demonstration
of effectiveness on real oceanic current data. Hardware validation is
discussed as future work (Section~\ref{sec:sep_future}).
\end{enumerate}

\subsection{Lagrangian Coherent Structures and the Objectivity Question}

Haller and Yuan formalized the modern notion of a Lagrangian coherent
structure (LCS) as the extension of a stable or unstable manifold to flows
with arbitrary time dependence, a distinguished material surface that
organizes tracer transport over a finite time horizon [55]. Haller's later
review frames the resulting body of theory as the search for the
Lagrangian skeleton of a flow, and states the property such a skeleton must
have to be physically meaningful: objectivity, invariance under the
time-dependent Euclidean frame changes that a material response cannot
depend on [52]. Shadden, Lekien, and Marsden gave the standard computable
definition, an LCS as a ridge of the finite-time Lyapunov exponent (FTLE)
field, and verified it on the same double-gyre model used throughout this
chapter, where for $\varepsilon>0$ the ridge tracks the underlying
heteroclinic connection [64]. Both the definition and its objectivity come
at a fixed cost: an FTLE ridge is a property of particle trajectories
integrated over a finite horizon, so extracting one requires the velocity
field's future evolution over that horizon, in advance, everywhere the
tracers travel. Haller's review is explicit that this is why the classical
instantaneous criteria, including Okubo-Weiss, the Q-criterion, and the
Hua-Klein criterion, are excluded from the Lagrangian framework entirely:
they can frame a coherent feature of the instantaneous velocity field, but
they are not objective [52]. Serra and Haller later placed a version of
this instantaneous idea on objective footing with objective Eulerian
coherent structures (OECSs), constructed from the spin and rate-of-strain
tensors rather than the raw velocity gradient, and showed OECSs are the
limit LCSs approach as the integration horizon shrinks to zero, a limit
Nolan, Serra, and Ross later made precise for the FTLE field itself [96,
98]. The velocity gradient that this instantaneous cost is paid to avoid is
itself hard to obtain from trajectory data alone: Harms, Dawson, and McKeon
show that standard gradient-based LCS diagnostics degrade badly on the
sparse tracer trajectories typical of ocean drifters, and propose a
regression-based alternative that linearizes in time rather than in space
to recover usable gradients from sparse data [54]. The rate-of-strain and
spin decomposition underlying all of these instantaneous diagnostics is
itself the subject of ongoing refinement; Haller's dynamic polar
decomposition, for instance, identifies a mean material rotation rate that
agrees with the classical vorticity-based rate used throughout this chapter
while correcting a memory effect present in the standard polar
decomposition of the deformation gradient [65]. A robot team sampling the
flow once per control cycle, with no forecast, sits on the instantaneous
side of the resulting tradeoff, objectivity and material accuracy against
zero integration latency, by construction; this chapter adopts that side
deliberately and returns to what is gained and given up by the choice in
Section~\ref{sec:separatrix}.

\subsection{Eulerian Surrogates: Okubo-Weiss, Q-Criterion, Hua-Klein}

Okubo and, independently, Weiss proposed the pointwise criterion
$Q=\tfrac14(s_n^2+s_s^2-\omega^2)$ to partition a two-dimensional flow into
rotation-dominated and strain-dominated regions using only the local
velocity gradient, with their sign convention marking vortex cores by
$Q<0$ [57, 58]. Hunt, Wray, and Moin generalized the same second-invariant
idea to three dimensions as the Q-criterion for identifying eddies,
streams, and convergence zones in turbulent flow fields [63]. Hua and Klein
revisited the two-dimensional criterion for oceanographic use, sharpening
the conditions under which it correctly separates elliptic and hyperbolic
regions of a nearly two-dimensional turbulent flow [59]. This chapter works
with $D=\det(\mathbf{J})$ rather than $Q$; for the incompressible flows
considered here the two carry identical information with opposite sign,
$Q=-D$ (Section~\ref{sec:detj_field}), so $D>0$ marks vortex cores and
$D<0$ marks strain-dominated regions throughout. What every criterion in
this family shares, and what distinguishes it from the FTLE ridges of
Section~\ref{sec:separatrix}, is that each is a scalar function of the
pointwise velocity-gradient tensor, requiring only the instantaneous field
at a point, not an integration over time. Approaches that reduce a sensed
vector field to a scalar this way lose more than time dependence, however:
Cooper et al.'s initial multirobot study of vector field navigation reduced
each measurement to its magnitude, which yields a descent direction but
cannot distinguish a sink from a saddle from a vortex, since the magnitude
vanishes at every critical point regardless of type [50]. Rather than
sensing a scalar reduction of the field directly, whether a speed magnitude
or an Okubo-Weiss-type invariant, we estimate the full local quadratic
vector field and obtain $D$, $\nabla D$, and $\mathbf{H}_D$ as derived
quantities from the jointly fitted components of $u$ and $v$
(Section~\ref{sec:estimation}), which is what makes the surrogate's
gradient and curvature, not merely its sign, available to the controller.

\subsection{Robots Tracking Manifolds and Coherent Structures}

The nearest prior work to this chapter tracks a coherent structure directly
with a robot formation rather than computing it offline. Michini, Hsieh,
Forgoston, and Schwartz proposed a PIM-triple-inspired saddle straddle:
three robots bracket a presumed stable or unstable manifold, using only
local velocity measurements to keep one robot on the manifold and the other
two on either side of it, with no FTLE field or global map required [34].
Michini, Rastgoftar, Hsieh, and Jayasuriya extended the same principle from
three robots to a general $N$-robot team for finer spatiotemporal sampling
of the boundary [13], and Michini, Hsieh, Forgoston, and Schwartz
separately validated the three-robot strategy experimentally on a team of
micro autonomous surface vehicles in a flow tank [43]. Knizhnik, Li, Yu,
and Hsieh applied a related flow-based control strategy to marine robots
operating in gyre-like environments, coupling manifold tracking to
vehicle-level flow control [42]. In each of these, the team's controller
commits to a manifold identity before tracking begins: robot C's estimate
of the manifold's location is read off the straddle segment as the point of
maximum measured speed if the unstable manifold is being tracked, or of
minimum speed if the stable manifold is being tracked, a choice fixed at
the start and never revisited [34]. The accompanying stability proof holds
only under the standing assumption that the straddle segment crosses the
manifold exactly once; neither the proof nor the controller addresses the
saddle point itself, where the linearized Jacobian has one positive
eigenvalue and the direction the straddle line must stay orthogonal to is
exactly the manifold the team is not tracking. This is structurally why the
family offers no guarantee about continuing past a crossing: the
manifold's identity is committed to in advance, and a crossing is precisely
where that identity becomes locally ambiguous. Kularatne and Hsieh took the
complementary approach of computing the local FTLE field on board from a
robot team's own accumulated velocity measurements, which recovers the LCS
type explicitly rather than assuming it, and proved tracking and
formation-keeping guarantees later validated on real ocean flow data [12];
this buys type discrimination with a nonzero integration horizon, in
contrast to the instantaneous estimate used here. None of these strategies
estimates the curvature of the field itself, which is what fixes not just
where a structure lies but what kind of structure it is and how many
branches meet there. The estimator of Section~\ref{sec:estimation} supplies
exactly that, and the separatrix tracker's traversal proof
(Section~\ref{sec:stability}) supplies the continuation guarantee this
family lacks: because the along-trench eigenvalue sign of $\mathbf{H}_D$ is
read fresh from the current instantaneous fit every control cycle, the
branch to continue onto at a crossing is measured rather than assumed.

\subsection{Distributed Gradient and Hessian Estimation}

Ögren, Fiorelli, and Leonard established the template for cooperative
gradient estimation, using a small team of mobile sensors in a fixed
formation to climb the gradient of a scalar field it samples collectively
[2]. Briñón-Arranz, Renzaglia, and Schenato extended this to the Hessian,
showing that a symmetric ring of five robots plus one at the center
recovers both the gradient and the Hessian of a scalar signal in closed
form, and that the resulting curvature information accelerates
Newton-Raphson source seeking over gradient ascent alone once the formation
radius is properly scaled to the noise floor [3]. McDonald, Kitts, and
Neumann developed the complementary line of adaptive navigation primitives,
using three- and four-robot clusters to reach and follow scalar-field
features such as ridges, trenches, and saddle points from local differences
alone, without an explicit Hessian estimate [25, 4], later verified
experimentally on physical hardware [5]. Hessian estimation of a single
scalar field is therefore established, but every prior formation samples
one signal channel. The problem addressed in this chapter, recovering the
Jacobian and its curvature for a coupled two-component vector field,
doubles the estimation problem: Section~\ref{sec:estimation} proves six
robots, not five plus a center, is the exact minimum for the coupled case,
a new estimation-theory result rather than an application of the scalar
one. It also derives $D$, $\nabla D$, and $\mathbf{H}_D$ by differentiating
through $\det(\cdot)$ on the jointly fitted vector coefficients, rather
than fitting $D$ or an Okubo-Weiss-type surrogate as an independent scalar
surface with its own Hessian estimator, which is what preserves the
coupling between $u$ and $v$ that a scalar-surrogate route would discard.
Section~\ref{sec:controllers} shows what this curvature is used for.

\subsection{Cluster Space Control and SAS Formations}

Chapter~\ref{ch:cluster_builder} already presented the cluster space
control lineage this chapter's formation belongs to in full, from Kitts
and Mas's original parameterization [46] through Adamek, Kitts, and Mas's
gradient-based navigation on physical hardware [7] to the general
compositional grammar this dissertation contributes. The side-angle-side
(SAS) triangle parameterization used for three-robot vector field
navigation in Chapter~\ref{ch:vector_field} is one instance of this family;
the six-robot formation of this chapter is built from three such SAS pairs
sharing a common centroid (Appendix~\ref{app:cluster_kinematics}), an
instance of Chapter~\ref{ch:cluster_builder}'s pentagon-plus-center atom
under its consensus orientation gauge. The formation parameterization the
controller regulates is decoupled from the estimator, which reads only the
six measured positions, so the SAS grouping used for actuation is invisible
to the estimation problem of Section~\ref{sec:estimation}. Coordinated
finite-time control results for multi-robot Euler-Lagrange systems, such as
the distributed finite-time tracking and containment control of Sun et al.
[53], point to an alternative convergence-rate analysis that the Lyapunov
arguments of Section~\ref{sec:stability} do not pursue, since this
chapter's controllers prioritize the geometric traversal guarantee over a
finite-time convergence rate. Where Chapter~\ref{ch:vector_field} used a
proportional or orbital law with no curvature, [3]'s Newton-Raphson source
seeking uses the raw Hessian inverse to converge to a scalar extremum, and
the saddle-free Newton method of Dauphin et al. flips signed eigenvalues to
escape saddles in high-dimensional optimization but was never applied to a
robot team tracking a spatial field [61]; the modified Newton step
introduced here is the first to combine per-eigendirection saturation with
the saddle-free sign flip in a distributed field-navigation setting.

\section{Problem Formulation}
\label{sec:detj_field}

\subsection{Flow Fields and Local Polynomial Models}

A planar flow assigns a velocity vector
$\mathbf{v}(\mathbf{p})=(u(\mathbf{p}),v(\mathbf{p}))$ to each position
$\mathbf{p}=(x,y)$ in its domain. Its Jacobian at $\mathbf{p}$ is
\begin{equation}
\mathbf{J}(\mathbf{p}) = \begin{bmatrix} \partial u/\partial x & \partial u/\partial y \\ \partial v/\partial x & \partial v/\partial y \end{bmatrix},
\label{eq:jacobian}
\end{equation}
and we write $D(\mathbf{p})=\det(\mathbf{J}(\mathbf{p}))$ throughout.

In Chapter~\ref{ch:vector_field}, a three-robot cluster fit a linear model
to each component of $\mathbf{v}$, assembled the estimated Jacobian
$\mathbf{J}$ and offset $\mathbf{h}$, and recovered the critical point
$\mathbf{p}^*=-\mathbf{J}^{-1}\mathbf{h}$ where $\mathbf{v}(\mathbf{p}^*)=\mathbf{0}$,
together with its type from the eigenvalues of $\mathbf{J}$. A linear
model, however, carries no information about how $\mathbf{J}$ itself
varies in space. The structures addressed in this chapter, boundaries
between rotation-dominated and strain-dominated flow and the separatrices
attached to saddle points, are level sets and valley lines of $D$, and
navigating them requires the first and second spatial derivatives of $D$.
These derivatives are polynomial functions of the first and second
derivatives of $u$ and $v$. The natural extension is therefore a local
quadratic model of each component,
\begin{align}
u(\mathbf{p}_c+\tilde{\mathbf{p}}) &\approx u_0 + u_x\tilde x + u_y\tilde y + \tfrac12 u_{xx}\tilde x^2 + u_{xy}\tilde x\tilde y + \tfrac12 u_{yy}\tilde y^2, \label{eq:quad_model_u} \\
v(\mathbf{p}_c+\tilde{\mathbf{p}}) &\approx v_0 + v_x\tilde x + v_y\tilde y + \tfrac12 v_{xx}\tilde x^2 + v_{xy}\tilde x\tilde y + \tfrac12 v_{yy}\tilde y^2, \label{eq:quad_model_v}
\end{align}
where $\tilde{\mathbf{p}}=(\tilde x,\tilde y)$ is position relative to the
cluster centroid $\mathbf{p}_c$ and the twelve coefficients are the Taylor
data of the field at $\mathbf{p}_c$: the value, gradient, and Hessian of
each component. Section~\ref{sec:estimation} shows that six robots are the
minimum team that can recover this model from one synchronous set of
measurements.

\subsection{Separatrices, Coherent Structures, and Eulerian Surrogates}
\label{sec:separatrix}

At a saddle point of a steady flow, the stable and unstable manifolds are
the curves of initial conditions that converge to the saddle in forward
and backward time, respectively. Together they form the separatrix, which
partitions the domain into regions whose trajectories do not mix. In
time-dependent flows the analogous objects are Lagrangian coherent
structures (LCSs), commonly identified as ridges of the finite-time
Lyapunov exponent (FTLE) field [55, 56, 64]. Tracking an LCS directly is
not compatible with the sensing model of this chapter: computing FTLE
requires integrating dense grids of particle trajectories over a finite
horizon, which presumes global knowledge of the field over both space and
time. A robot team that measures the flow only at its own footprint, at
the current instant, cannot evaluate it.

The alternative is an Eulerian surrogate: a quantity computed pointwise
from the instantaneous velocity field and its spatial derivatives. Serra
and Haller formalized this idea as objective Eulerian coherent structures
(OECSs), the instantaneous limits of LCSs [60], and Nolan et al. showed
that FTLE fields converge to Eulerian rate-of-strain quantities as the
integration horizon shrinks [62]. This chapter adopts the classical
Okubo-Weiss field as its surrogate. It is not identical to the FTLE ridge
set, and unlike the structures of [60] it is not objective under
time-varying rotations of the observer frame, but it is Galilean
invariant, it requires only the local Jacobian and its derivatives, and
for the flows considered here it coincides with or closely tracks the
structures of interest. The correspondence is exact for the steady double
gyre (Section~\ref{sec:dg_example}) and is evaluated against FTLE fields
computed offline from ocean radar data in Section~\ref{sec:results}.

\subsection{The Okubo-Weiss Field and \texorpdfstring{$\det(\mathbf{J})$}{det(J)}}

Write the entries of $\mathbf{J}$ as in Equation~\eqref{eq:jacobian}. The
normal strain rate, shear strain rate, and vorticity are
$s_n=\partial u/\partial x-\partial v/\partial y$,
$s_s=\partial v/\partial x+\partial u/\partial y$, and
$\omega=\partial v/\partial x-\partial u/\partial y$. The Okubo-Weiss
parameter [57, 58] is $Q=\tfrac14(s_n^2+s_s^2-\omega^2)$, and direct
expansion gives the algebraic identity
\begin{equation}
Q = \tfrac14(\nabla\cdot\mathbf{v})^2 - D.
\label{eq:okubo_weiss}
\end{equation}
For the incompressible flows considered in this chapter,
$\nabla\cdot\mathbf{v}=0$ and $Q=-D$ exactly, so the Okubo-Weiss field and
the Jacobian determinant carry the same information with opposite sign. We
work with $D$.

The sign of $D$ fixes the local eigenstructure. Incompressibility gives
$\operatorname{tr}(\mathbf{J})=0$, so the eigenvalues of $\mathbf{J}$
satisfy $\lambda^2=-D$:
\begin{equation}
\lambda = \pm\sqrt{-D} \ \ (D<0), \qquad \lambda = \pm i\sqrt{D} \ \ (D>0).
\label{eq:eig_detj}
\end{equation}
Where $D<0$ the flow is locally hyperbolic (strain-dominated) and nearby
particles separate at the instantaneous exponential rate $\sqrt{-D}$; where
$D>0$ it is locally elliptic (rotation-dominated). The convention used
throughout this chapter is therefore: $D>0$ marks vortex cores, $D<0$
marks strain regions, and the zero level set
\begin{equation}
\mathcal{C} = \{\mathbf{p} : D(\mathbf{p})=0\}
\label{eq:ow_boundary}
\end{equation}
is the Okubo-Weiss boundary between the two regimes. By
Equation~\eqref{eq:eig_detj}, the more negative $D$ is, the stronger the
local hyperbolicity, so valley lines (trenches) of $D$ inside strain
regions mark the loci of strongest instantaneous attraction and repulsion.
These are the two Eulerian structures tracked in this chapter: the
boundary $\mathcal{C}$ and the trench of $D$.

\subsection{Running Example: the Double Gyre}
\label{sec:dg_example}

The double-gyre flow used throughout is the steady ($\varepsilon=0$)
member of the Shadden family [64], shifted from the canonical domain
$[0,2]\times[0,1]$ to $[-1,1]\times[-0.5,0.5]$ by $x_f=x+1$, $y_f=y+0.5$,
where $(x_f,y_f)$ are field-frame and $(x,y)$ world-frame coordinates. The
field is
\begin{align}
u(\mathbf{p}) &= -\pi A\sin(\pi x_f)\cos(\pi y_f), \label{eq:dg_u} \\
v(\mathbf{p}) &= \pi A\cos(\pi x_f)\sin(\pi y_f), \label{eq:dg_v}
\end{align}
with $A=0.1$. The flow is incompressible, the separatrix is the line
$x=0$, the saddle points are $\mathbf{p}^*_1=(0,-0.5)$ and
$\mathbf{p}^*_2=(0,0.5)$, and the gyre centers are $(\pm0.5,0)$.

All quantities of interest have closed forms (derivations in
Appendix~\ref{app:double_gyre}):
\begin{equation}
D(\mathbf{p}) = -\frac{\pi^4A^2}{2}\bigl(\cos2\pi x_f + \cos2\pi y_f\bigr),
\label{eq:det_analytic_main}
\end{equation}
\begin{equation}
\nabla D = \pi^5A^2\begin{bmatrix}\sin2\pi x_f\\\sin2\pi y_f\end{bmatrix},
\qquad
\mathbf{H}_D = 2\pi^6A^2\begin{bmatrix}\cos2\pi x_f & 0\\0 & \cos2\pi y_f\end{bmatrix}.
\label{eq:gradhess_analytic_main}
\end{equation}
Three facts follow directly and organize everything that comes after.

\paragraph{The Okubo-Weiss boundary is a diamond.} Since
$\cos2\pi x_f+\cos2\pi y_f=2\cos\pi(x_f+y_f)\cos\pi(x_f-y_f)$, the zero set
$\mathcal{C}$ consists of the straight lines
$x_f\pm y_f\in\tfrac12+\mathbb{Z}$: a diamond inscribed around each gyre
center, with $D>0$ inside (vortex core, where $D=\pi^4A^2$ at the center)
and $D<0$ outside. The four corners of each diamond are the only points of
$\mathcal{C}$ where $\nabla D=\mathbf{0}$.

\paragraph{The separatrix is a trench of $D$.} On the line $x=0$ (that is,
$x_f=1$), the cross-track derivative
$\partial D/\partial x=\pi^5A^2\sin2\pi x_f$ vanishes identically and the
cross-track curvature
$\partial^2D/\partial x^2=2\pi^6A^2\cos2\pi x_f=2\pi^6A^2>0$ is positive.
The separatrix is therefore an exact valley line of the signed field $D$.
Moreover the vorticity $\omega=-2\pi^2A\sin(\pi x_f)\sin(\pi y_f)$ vanishes
on this line, so there $\mathbf{J}$ equals the rate-of-strain tensor and
$\sqrt{-D}$ is exactly the instantaneous attraction rate toward the
separatrix. The Eulerian trench and the Lagrangian structure coincide.

\paragraph{The trench is one line of a periodic grid.} Restricted to the
separatrix, $D(0,y)=-\tfrac{\pi^4A^2}{2}(1+\cos2\pi y_f)$: it is deepest
($D=-\pi^4A^2$) at the two saddle points, where the local minima of $D$
sit, and rises to a crest ($D=0$) at the domain midpoint $(0,0)$, where the
two gyres' diamonds touch. A controller that only descends $D$ along the
trench cannot pass the crest: the along-trench gradient vanishes there,
and from starts on the upstream side descent leads away from the flow
direction, toward the upstream well. This motivates the flow-assisted mode
in Section~\ref{sec:sep_controller}. The ambient flow on the separatrix is
tangent to it, directed from the upper saddle toward the lower saddle, and
vanishes only at the saddles.

The trench does not end at a saddle. Since
$\partial^2D/\partial x^2=2\pi^6A^2\cos2\pi x_f$ and
$\partial^2D/\partial y^2=2\pi^6A^2\cos2\pi y_f$ are both periodic, every
line $x_f\in\mathbb{Z}$ is a valley of $D$ in $x$ and every line
$y_f\in\mathbb{Z}$ is a valley of $D$ in $y$; the world-frame separatrix
$x=0$ (that is, $x_f=1$) is one member of the first family. These valley
lines form an infinite grid, and the saddle points are exactly its
crossings: at a saddle, the world-frame separatrix meets the horizontal
trench line through that saddle at a right angle, both branches being
valleys of $D$ in their own cross-track direction. A tracker that reaches
the saddle and continues, rather than stopping, is therefore still on a
trench, now running perpendicular to the one it arrived on.
Section~\ref{sec:stab_sep} states the traversal result for this network;
Section~\ref{sec:results} reports which behavior, capture or continuation,
the estimator in Section~\ref{sec:estimation} actually produces.

% %% TODO: SYNC WITH LATEST DRAFT -- figure not yet generated for thesis;
% see fig:dg_streamlines in Paper_Draft_2A.tex (comment there: generate
% with double_gyre_static from src/fields/environments/Double_Gyre.py).
\begin{figure}[!t]
\centering
\framebox[0.9\linewidth]{\rule{0pt}{3.0cm}}
\caption{Double-gyre vector field (steady case, $A=0.1$). Streamlines are
shown in gray. The separatrix at $x=0$ (dashed) is the stable manifold of
the lower saddle and the unstable manifold of the upper saddle ($\times$).
The $D=0$ Okubo-Weiss boundary (solid) is a diamond around each gyre
center ($\bullet$).}
\label{fig:dg_streamlines}
\end{figure}

% %% TODO: SYNC WITH LATEST DRAFT -- figure not yet generated for thesis;
% see fig:detJ_contour in Paper_Draft_2A.tex.
\begin{figure}[!t]
\centering
\framebox[0.9\linewidth]{\rule{0pt}{3.0cm}}
\caption{$D=\det(\mathbf{J})$ over the double-gyre domain. The $D=0$
boundary (bold) separates rotation-dominated cores ($D>0$) from the
strain-dominated exterior ($D<0$). The separatrix at $x=0$ is a valley
line of $D$ connecting the crest at the origin to the two wells at the
saddle points.}
\label{fig:detJ_contour}
\end{figure}

\begin{figure}[!t]
\centering
\framebox[0.9\linewidth]{\rule{0pt}{3.0cm}}
\caption{Signed $D=\det(\mathbf{J})$ as an elevation surface over the
double-gyre domain. The two gyre cores rise as ridges ($D>0$,
rotation-dominated) and the exterior falls below zero ($D<0$,
strain-dominated); the dashed curve on the floor is the $D=0$ boundary.
The separatrix at $x=0$ (bold) is a valley of $D$ running between the two
ridges, its floor rising to $D=0$ at the origin and deepening toward the
saddle points.}
\label{fig:trench_cross}
\end{figure}

\subsection{Problem Statement}
\label{sec:problem_statement}

Consider $N$ robots at positions $\mathbf{p}_i$, each measuring the local
flow $(u_i,v_i)=\mathbf{v}(\mathbf{p}_i)$ at the common sampling instant.
No map, forecast, or prior model of the field is available, and no robot
can integrate trajectories of the field. Using only the current
synchronous samples $\{(\mathbf{p}_i,u_i,v_i)\}_{i=1}^N$ at each control
cycle, the team must generate a centroid velocity command
$\dot{\mathbf{p}}_c$ that solves:

\emph{Problem 1 (separatrix network traversal).} Drive the cluster
centroid onto the trench of $D$ and travel it, in the direction of the
ambient flow, through the saddle points where trench segments cross.
Terminal capture at a single saddle (Section~\ref{sec:sep_controller}) is
the special case in which the selector's stationary-point test is
tightened to stop the team there instead.

\emph{Problem 2 (Okubo-Weiss boundary tracking).} Drive
$D(\mathbf{p}_c)\to0$ and travel along the boundary $\mathcal{C}$, a
network of the same kind (Section~\ref{sec:dg_example}).

Both problems require, at minimum, pointwise estimates of $D$, its
gradient $\mathbf{g}=\nabla D$, and its Hessian $\mathbf{H}_D$ at the
centroid. The next section shows how a six-robot formation produces them.

\section{Distributed Second-Order Estimation}
\label{sec:estimation}

\subsection{Quadratic Fit from Six Simultaneous Measurements}

Let $\tilde{\mathbf{p}}_i=\mathbf{p}_i-\mathbf{p}_c$ denote robot positions
relative to the centroid and define the quadratic basis vector
\begin{equation}
\boldsymbol{\phi}(\tilde{\mathbf{p}}) = \bigl[\,1,\ \tilde x,\ \tilde y,\ \tfrac12\tilde x^2,\ \tilde x\tilde y,\ \tfrac12\tilde y^2\,\bigr]^{\mathsf T},
\label{eq:basis}
\end{equation}
whose one-half factors make the fitted coefficients equal to the Taylor
derivatives in Equations~\eqref{eq:quad_model_u}--\eqref{eq:quad_model_v}.
Stacking the basis vectors of the six robots row-wise gives the
$6\times6$ formation matrix $\boldsymbol{\Phi}$, the second-order
generalization of the $3\times3$ matrix $\mathbf{A}$ of
Chapter~\ref{ch:vector_field}. The coefficient vectors
$\boldsymbol{\theta}_u=[u_0,u_x,u_y,u_{xx},u_{xy},u_{yy}]^{\mathsf T}$ and
$\boldsymbol{\theta}_v=[v_0,v_x,v_y,v_{xx},v_{xy},v_{yy}]^{\mathsf T}$
solve the pair of linear systems
\begin{equation}
\boldsymbol{\Phi}\boldsymbol{\theta}_u = \mathbf{u}, \qquad \boldsymbol{\Phi}\boldsymbol{\theta}_v = \mathbf{v},
\label{eq:solve}
\end{equation}
where $\mathbf{u}=[u_1,\ldots,u_6]^{\mathsf T}$ and
$\mathbf{v}=[v_1,\ldots,v_6]^{\mathsf T}$ collect the measurements. One
matrix factorization serves both components. When the field is exactly
quadratic over the formation footprint, the solution of
Equation~\eqref{eq:solve} is exact for any nonsingular $\boldsymbol{\Phi}$;
for general smooth fields the model error is quantified in
Section~\ref{sec:sensitivity}.

\subsection{Six Robots Are Necessary, and Almost Always Sufficient}

\begin{proposition}[Minimality]
\label{prop:minimality}
Each robot contributes exactly one linear equation to each system in
Equation~\eqref{eq:solve}, and each system has six unknowns. Hence no
algorithm, linear or otherwise, can uniquely recover the quadratic model
Equations~\eqref{eq:quad_model_u}--\eqref{eq:quad_model_v} from one
synchronous sample of fewer than six robots, and six robots suffice
exactly when $\boldsymbol{\Phi}$ is nonsingular.
\end{proposition}

\begin{proof}
With $N<6$ robots, $\boldsymbol{\Phi}\in\mathbb{R}^{N\times6}$ has a
nontrivial null space, so for any candidate coefficient vector
$\boldsymbol{\theta}_u$ the vector
$\boldsymbol{\theta}_u+\boldsymbol{\theta}_0$ with
$\boldsymbol{\Phi}\boldsymbol{\theta}_0=\mathbf{0}$,
$\boldsymbol{\theta}_0\neq\mathbf{0}$, generates identical measurements at
every robot. The two quadratic fields are indistinguishable from the data,
so no estimator can select between them. With $N=6$,
Equation~\eqref{eq:solve} has a unique solution if and only if
$\det(\boldsymbol{\Phi})\neq0$.
\end{proof}

Nonsingularity has a clean geometric characterization.

\begin{proposition}[Conic criterion]
\label{prop:conic}
$\boldsymbol{\Phi}$ is singular if and only if the six sampling points lie
on a common conic, that is, on the zero set of some nonzero quadratic
polynomial $q(\tilde x,\tilde y)$ (including degenerate conics such as
line pairs).
\end{proposition}

\begin{proof}
$\boldsymbol{\Phi}$ is singular iff its columns are linearly dependent,
iff there is a nonzero coefficient vector $\mathbf{c}$ with
$\boldsymbol{\phi}(\tilde{\mathbf{p}}_i)^{\mathsf T}\mathbf{c}=0$ for all
$i$. The function $q=\boldsymbol{\phi}^{\mathsf T}\mathbf{c}$ is precisely
a nonzero polynomial of degree at most two vanishing at all six points.
\end{proof}

\begin{corollary}[Pentagon plus center]
\label{cor:pentagon}
Place five robots at the vertices of a regular pentagon of radius $\rho$
and the sixth at its center. Then $\boldsymbol{\Phi}$ is nonsingular.
\end{corollary}

\begin{proof}
No three vertices of a regular pentagon are collinear, and five points in
general position determine a unique conic; for the pentagon vertices that
conic is the circumscribed circle. The center does not lie on the circle,
so no conic passes through all six points, and Proposition~\ref{prop:conic}
applies.
\end{proof}

\begin{remark}[The center robot is essential]
Any formation whose six robots lie on one circle, for example a regular
hexagon, is singular regardless of its size or orientation, since the
circle itself is a conic through all six points. Rings are natural
formations for mobile robots, which makes this failure mode easy to adopt
by accident; moving one robot to the center removes it entirely.
\end{remark}

\subsection{The Formation and Its Noise Gains}
\label{sec:formation}

The pentagon-plus-center shape is the same minimal configuration that
Briñón-Arranz et al. derived for Hessian estimation of a scalar field:
$N\geq5$ robots uniformly spaced on a circle plus one at the center [3].
Their estimator combines the ring measurements with fixed symmetric
weights, which is computationally trivial and, for exactly six robots on
the ideal formation, algebraically equivalent to solving
Equation~\eqref{eq:solve}. The two approaches separate when the formation
deforms, as it always does under real formation control: fixed weights
assume the ideal geometry and acquire a bias proportional to the
deformation, whereas Equation~\eqref{eq:solve} uses the measured robot
positions and remains exact for quadratic fields under any deformation
short of the conic degeneracy of Proposition~\ref{prop:conic}. The
estimator therefore degrades gracefully: formation error affects accuracy
only through the conditioning of $\boldsymbol{\Phi}$, which is monitored
online through $\det(\boldsymbol{\Phi})$.

For the ideal formation the noise amplification has closed form. Let each
measurement carry independent zero-mean noise of standard deviation
$\sigma_{uv}$. Since
$\hat{\boldsymbol{\theta}}_u-\boldsymbol{\theta}_u=\boldsymbol{\Phi}^{-1}\boldsymbol{\eta}$,
the standard deviation of each coefficient is $\sigma_{uv}$ times the
Euclidean norm of the corresponding row of $\boldsymbol{\Phi}^{-1}$.

\begin{lemma}[Exact noise gains of the pentagon-plus-center formation]
\label{lem:gains}
For five robots on a ring of radius $\rho$ and one at the centroid, the
coefficient standard deviations are
\begin{equation}
\sigma_{u_0} = \sigma_{uv}, \qquad
\sigma_{u_x} = \sigma_{u_y} = \frac{2}{\sqrt{10}}\frac{\sigma_{uv}}{\rho},
\label{eq:gain_grad}
\end{equation}
\begin{equation}
\sigma_{u_{xx}} = \sigma_{u_{yy}} = \frac{8}{\sqrt{10}}\frac{\sigma_{uv}}{\rho^2},
\qquad
\sigma_{u_{xy}} = \frac{4}{\sqrt{10}}\frac{\sigma_{uv}}{\rho^2},
\label{eq:gain_hess}
\end{equation}
independent of the formation heading, and identically for
$\boldsymbol{\theta}_v$.
\end{lemma}

\begin{proof}
The value coefficient is exact: every basis function except the constant
vanishes at the origin, so interpolation forces $\hat u_0$ to equal the
center robot's reading, giving unit gain. The remaining rows follow from
direct inversion of $\boldsymbol{\Phi}$ using the fifth roots of unity;
the row norms reduce to the stated constants, and repeating the
computation with an arbitrary ring phase leaves them unchanged. The error
covariance of the estimated gradient is in fact isotropic,
$\tfrac{4}{10}\sigma_{uv}^2\rho^{-2}\mathbf{I}$, for every heading.
\end{proof}

Two design consequences. First, the gains are isotropic: the error
covariance of the estimated gradient (and of the Hessian, entrywise as
stated) does not depend on the direction of travel, so estimator quality
does not degrade as the cluster turns. Second, the $1/\rho$ and
$1/\rho^2$ factors exhibit the fundamental scaling: second derivatives are
harder to sense than first derivatives by one power of the formation
radius. A numerical search over all six-point formations of the same
footprint found configurations at best a few percent better in worst-case
Hessian gain, so the symmetric shape is used without further tuning.

% %% TODO: SYNC WITH LATEST DRAFT -- figure not yet generated for thesis;
% see fig:pentagon_formation in Paper_Draft_2A.tex.
\begin{figure}[!t]
\centering
\framebox[0.9\linewidth]{\rule{0pt}{3.0cm}}
\caption{Six-robot formation: five robots on a ring of radius $\rho$ at
$72°$ spacing and one robot at the center, the minimal symmetric
formation for second-order estimation (Corollary~\ref{cor:pentagon}, [3]).
For formation control the six robots are grouped into three pairs whose
midpoints define an SAS triangle (Appendix~\ref{app:cluster_kinematics});
the grouping is invisible to the estimator, which uses only the six
measured positions.}
\label{fig:pentagon_formation}
\end{figure}

\subsection{Closed-Form Recovery of \texorpdfstring{$D$}{D}, \texorpdfstring{$\nabla D$}{grad D}, and \texorpdfstring{$\mathbf{H}_D$}{H\_D}}

The estimated Jacobian at the centroid is assembled from the linear
coefficients,
\begin{equation}
\hat{\mathbf{J}} = \begin{bmatrix} u_x & u_y \\ v_x & v_y \end{bmatrix}, \qquad
\hat D = u_xv_y - u_yv_x.
\label{eq:J_hat}
\end{equation}
Under the quadratic model, each entry of $\mathbf{J}$ is an affine function
of position, so $D(\mathbf{p})$ is itself an exactly quadratic polynomial
over the footprint. Its gradient and Hessian at the centroid are
polynomial functions of the fitted coefficients:
\begin{equation}
\hat{\mathbf{g}} = \nabla\hat D = \begin{bmatrix} u_{xx}v_y + u_xv_{xy} - u_{xy}v_x - u_yv_{xx} \\ u_{xy}v_y + u_xv_{yy} - u_{yy}v_x - u_yv_{xy} \end{bmatrix},
\label{eq:grad_det}
\end{equation}
\begin{equation}
\hat{\mathbf{H}}_D = \begin{bmatrix}
2(u_{xx}v_{xy}-u_{xy}v_{xx}) & u_{xx}v_{yy}-u_{yy}v_{xx} \\
u_{xx}v_{yy}-u_{yy}v_{xx} & 2(u_{xy}v_{yy}-u_{yy}v_{xy})
\end{bmatrix}.
\label{eq:hess_det}
\end{equation}
These follow by differentiating $D=u_xv_y-u_yv_x$ with the third-order
terms absent, and require no field evaluations beyond the six
measurements. The omission is harmless for the gradient but not for the
Hessian; its systematic consequences are quantified in
Section~\ref{sec:sensitivity}. The estimated flow at the centroid,
$\mathbf{f}=(u_0,v_0)$, equals the center robot's reading by
Lemma~\ref{lem:gains}. Together
Equations~\eqref{eq:J_hat}--\eqref{eq:hess_det} supply every quantity the
controllers of Section~\ref{sec:controllers} consume: $\hat D$,
$\hat{\mathbf{g}}$, $\hat{\mathbf{H}}_D$, and $\mathbf{f}$.

\subsection{Sensitivity: Noise, Truncation, and Formation Radius}
\label{sec:sensitivity}

Three error sources separate cleanly. Write
$\boldsymbol{\Phi}=\bar{\boldsymbol{\Phi}}\,\mathbf{S}_\rho$ with
$\mathbf{S}_\rho=\operatorname{diag}(1,\rho,\rho,\rho^2,\rho^2,\rho^2)$,
where $\bar{\boldsymbol{\Phi}}$ depends only on the formation shape. All
scale dependence enters through $\mathbf{S}_\rho^{-1}$.

\paragraph{Measurement noise.} By Lemma~\ref{lem:gains}, gradient
coefficients degrade as $\sigma_{uv}/\rho$ and Hessian coefficients as
$\sigma_{uv}/\rho^2$: larger formations average noise better.

\paragraph{Truncation bias.} If the field is three times differentiable
with $|\partial^3u|,|\partial^3v|\leq M_3$ over the footprint, the Taylor
remainder at robot $i$ is bounded by $\tfrac{M_3}{6}\rho^3$, and passing it
through $\mathbf{S}_\rho^{-1}\bar{\boldsymbol{\Phi}}^{-1}$ bounds the
systematic error by
\begin{equation}
|\hat u_x - u_x| \leq c_1M_3\rho^2, \qquad |\hat u_{xx}-u_{xx}| \leq c_2M_3\rho,
\label{eq:truncation}
\end{equation}
and likewise for the remaining coefficients, where $c_1,c_2$ are absolute
constants of the formation shape (row sums of $|\bar{\boldsymbol{\Phi}}^{-1}|$
divided by six). The Hessian bias decays only linearly in $\rho$, matching
the $O(\rho)$ bound proved for the symmetric estimator in [3]: smaller
formations track curvature better.

\paragraph{The radius trade-off.} Combining the two effects, the total
error of each fitted second-order coefficient is bounded by
$c_2M_3\rho + c_2'\sigma_{uv}/\rho^2$, which is minimized at
\begin{equation}
\rho^* = \left(\frac{2c_2'\sigma_{uv}}{c_2M_3}\right)^{1/3} \propto \left(\frac{\sigma_{uv}}{M_3}\right)^{1/3}.
\label{eq:rho_star}
\end{equation}
The formation radius is thus a sensing parameter, not merely a spacing
choice: it should grow with the noise floor and shrink with the field's
third-derivative scale. This rule is confirmed for $\hat D$ and
$\nabla\hat D$ by the formation-radius sweeps reported in
Section~\ref{sec:results}; the derived Hessian $\hat{\mathbf{H}}_D$ is
instead limited by the model bias discussed next.

\paragraph{Propagation to \texorpdfstring{$D$}{D} and its derivatives.}
Errors in the fitted coefficients enter the navigation quantities through
products of coefficient pairs. For two-by-two matrices the expansion of
the determinant is exact at second order:
$\det(\mathbf{J}+\mathbf{E}) = \det(\mathbf{J}) + \operatorname{tr}(\operatorname{adj}(\mathbf{J})\mathbf{E}) + \det(\mathbf{E})$,
and $\lVert\operatorname{adj}(\mathbf{J})\rVert_F = \lVert\mathbf{J}\rVert_F$,
so the estimated determinant satisfies
\begin{equation}
|\hat D - D| \;\leq\; \lVert\mathbf{J}\rVert_F\lVert\mathbf{E}\rVert_F + \tfrac12\lVert\mathbf{E}\rVert_F^2,
\label{eq:det_perturbation}
\end{equation}
where $\mathbf{E}=\hat{\mathbf{J}}-\mathbf{J}$ collects the four
first-order coefficient errors. The gradient formula
Equation~\eqref{eq:grad_det} is complete: differentiating $D=u_xv_y-u_yv_x$
once produces only products of first- and second-order derivatives, so
$\hat{\mathbf{g}}$ inherits errors solely from the fitted coefficients,
linearly, with slopes set by the true derivative magnitudes.

The Hessian is different in kind. Differentiating twice produces terms
such as $u_{xxx}v_y$ and $u_yv_{xxx}$, third derivatives of the field
weighted by the field itself, and a quadratic model has no third
derivatives: Equation~\eqref{eq:hess_det} is the Hessian of the fitted
quadratic surface, not of $D$. The omitted terms constitute a model bias
that is independent of the formation radius and cannot be reduced by
averaging; the situation is that of fitting a line to a curved signal and
using the line's slope where the curvature is unmodeled, acceptable
exactly where the model is locally adequate. For the double gyre at the
nominal scale this bias is large in norm, on the order of
$\lVert\mathbf{H}_D\rVert_F$ itself (Section~\ref{sec:results}). Two
conclusions nevertheless carry into the stability analysis: the estimate
errors entering the controllers are bounded whenever the coefficient
errors are bounded, with explicit formation-dependent constants; and all
such constants degrade continuously, not abruptly, as the formation
deforms, with outright failure only on the conic configurations of
Proposition~\ref{prop:conic}.

\begin{remark}[What survives the Hessian bias]
\label{rem:eigvec}
The controllers of Section~\ref{sec:controllers} consume
$\hat{\mathbf{H}}_D$ only through its eigenvectors and the signs and rough
magnitudes of its eigenvalues, and these degrade far less than the
Frobenius error suggests. The double gyre is symmetric under reflection
across the separatrix, so $D$ is even in the transverse coordinate and
every odd transverse derivative, including the mixed second derivative,
vanishes on the trench: both the retained and the omitted terms are
diagonal in the trench frame there. The estimated eigen-directions on the
separatrix are exact to numerical precision, remain within a few tenths of
a degree across the tube, and within a few degrees at generic points
(Section~\ref{sec:results}), while the bias loads almost entirely onto the
eigenvalues. The transverse Newton snap divides the transverse gradient by
the transverse eigenvalue, so an eigenvalue bias rescales the step, which
the saturation (Equation~\eqref{eq:sat}) bounds and the closed loop
absorbs; it does not rotate the commanded direction. This is why a
formation too small to sense curvature accurately can still track the
structure that curvature defines. The exception is the immediate
neighborhood of the wells, where the biased eigenvalues nearly vanish and
their signs become unreliable; the consequences for terminal capture are
examined in Section~\ref{sec:results}.
\end{remark}

\section{Adaptive Navigation Controllers}
\label{sec:controllers}

\subsection{Control Architecture}
\label{sec:arch}

Both controllers run inside the three-layer architecture of
Chapters~\ref{ch:testbed} and~\ref{ch:cluster_builder}:
\begin{enumerate}
\item \textbf{Navigation primitive.} From the six positions and
measurements, compute the estimates
Equations~\eqref{eq:J_hat}--\eqref{eq:hess_det} and output a centroid
velocity command $\dot{\mathbf{p}}_c^{\,\text{des}}$. The two primitives
of this section differ only in this block.
\item \textbf{Cluster space controller.} Using the twelve-state cluster
parameterization and its analytic inverse kinematic Jacobian
(Appendix~\ref{app:cluster_kinematics}), distribute the centroid command
to six robot velocity commands while regulating the formation shape to
its nominal geometry.
\item \textbf{Robot dynamics.} Each robot obeys the first-order momentum
model identified on the Decabot hardware (Chapter~\ref{ch:vector_field}),
\begin{equation}
\mathbf{v}[k+1] = \alpha_{\text{mom}}\mathbf{v}[k] + (1-\alpha_{\text{mom}})\mathbf{v}^{\text{des}}[k],
\label{eq:momentum}
\end{equation}
with $\alpha_{\text{mom}}=\exp(-\Delta t/\tau)$, actuator time constant
$\tau=0.3$~s, control period $\Delta t=0.1$~s (so
$\alpha_{\text{mom}}\approx0.717$), speed limit
$v_{\text{max,robot}}=0.3$~m/s, and stiction floor
$v_{\text{stiction}}=0.025$~m/s.
\end{enumerate}
Note on notation: $\alpha_{\text{mom}}$ in
Equation~\eqref{eq:momentum} is the momentum coefficient of the robot
model; it is unrelated to $\alpha_{\text{eig}}$, the real part of a
complex eigenvalue pair of $\mathbf{J}$ used in field classification in
Chapter~\ref{ch:vector_field}.

\subsection{From the Newton Step to the Modified Newton Step}
\label{sec:mod_newton}

Under the quadratic model, $\hat D(\mathbf{p})$ is an exact quadratic
surface with gradient $\hat{\mathbf{g}}$ and constant Hessian
$\hat{\mathbf{H}}_D$ at the centroid, so the natural navigation move is
the Newton step
$\Delta\mathbf{p}=-\hat{\mathbf{H}}_D^{-1}\hat{\mathbf{g}}$, which jumps to
the stationary point of the local model in one move. Written in the
eigenbasis $\hat{\mathbf{H}}_D\mathbf{w}_k=\lambda_k\mathbf{w}_k$, with
$\lambda_1\leq\lambda_2$ throughout, the step decomposes into
per-direction components $-(\mathbf{w}_k^{\mathsf T}\hat{\mathbf{g}})/\lambda_k$.
Two defects make the raw step unsuitable as a velocity command, and each
is repaired by one modification.

First, the magnitude is unbounded and explodes where $\hat{\mathbf{H}}_D$
is near singular. The repair is componentwise saturation through
\begin{equation}
\operatorname{sat}(c) = v_{\max}\tanh\!\left(\frac{c}{v_{\max}}\right),
\label{eq:sat}
\end{equation}
which gives a bounded constant-speed reach far from the structure and
recovers the exact Newton component, with its superlinear local
convergence, once $|c|\ll v_{\max}$. A single parameter $v_{\max}$ sets
both the cruise speed and the width of the terminal boundary layer. The
saturated Newton step is the \emph{attract step},
\begin{equation}
\dot{\mathbf{p}}_c^{\,\text{des}} = \sum_{k=1}^2 \operatorname{sat}\!\left(-\frac{\mathbf{w}_k^{\mathsf T}\hat{\mathbf{g}}}{\lambda_k}\right)\mathbf{w}_k.
\label{eq:attract_step}
\end{equation}

Second, and more fundamentally, Equation~\eqref{eq:attract_step} converges
to whatever stationary point the local model exhibits, with no preference
among minima, maxima, and saddle points of $D$. On a trench the
along-trench curvature can be negative (the crest side), and there the
Newton component deliberately climbs. The modification, adopted from
saddle-free Newton methods in optimization [61], replaces the signed
eigenvalue by its magnitude, giving the \emph{slide step},
\begin{equation}
\dot{\mathbf{p}}_c^{\,\text{des}} = \sum_{k=1}^2 \operatorname{sat}\!\left(-\frac{\mathbf{w}_k^{\mathsf T}\hat{\mathbf{g}}}{|\lambda_k|}\right)\mathbf{w}_k.
\label{eq:slide_step}
\end{equation}
In directions of positive curvature the two laws agree and snap the
centroid onto the trench floor; in the direction of negative curvature
Equation~\eqref{eq:slide_step} reverses the Newton component into a
curvature-normalized descent that slides along the trench, away from
crests. Both laws are well defined despite the sign ambiguity of
eigenvectors: every term is quadratic in $\mathbf{w}_k$, so replacing
$\mathbf{w}_k$ by $-\mathbf{w}_k$ leaves the command unchanged.

\subsection{Controller 1: Separatrix Tracker}
\label{sec:sep_controller}

The separatrix tracker composes the two step bodies with a third,
flow-assisted mode. The need for the third mode is visible in the trench
profile of Section~\ref{sec:dg_example}: the slide step descends $D$
along the trench, but at the crest the along-trench gradient vanishes,
and on the upstream side of the crest descent points away from the flow
direction, while the separatrix continues through the crest. On the
separatrix itself, however, the ambient flow is tangent to the structure
and nonzero away from the saddles, so the flow can carry the team across
the crest. The cluster is declared on the separatrix (FLOW band) when
either
\begin{equation}
|\hat D| < \varepsilon_{\text{raw}} \quad\text{or}\quad \frac{|\hat D|}{\lVert\hat{\mathbf{H}}_D\rVert_F} < \varepsilon_{\text{dim}},
\label{eq:flow_band}
\end{equation}
where $\varepsilon_{\text{raw}}$ is a raw threshold and
$\varepsilon_{\text{dim}}$ a curvature-normalized threshold (units of
length squared) that makes the band scale with the local landscape. In
the FLOW band the command follows the measured flow along the trench
tangent $\mathbf{w}_1$ and retains the Newton snap across it: the
\emph{flow step} is
\begin{equation}
\dot{\mathbf{p}}_c^{\,\text{des}} = \operatorname{sat}\!\left(\mathbf{f}^{\mathsf T}\mathbf{w}_1\right)\mathbf{w}_1 + \operatorname{sat}\!\left(-\frac{\hat{\mathbf{g}}^{\mathsf T}\mathbf{w}_2}{\lambda_2}\right)\mathbf{w}_2.
\label{eq:flow_step}
\end{equation}
Off the band, the local flow direction selects between the attract and
slide laws: if the flow agrees with the along-trench descent direction
the team slides (Equation~\eqref{eq:slide_step}); if it opposes descent
the team attracts toward the local stationary point
(Equation~\eqref{eq:attract_step}), which lies at the crest, inside the
FLOW band, so the attract law hands the team over to the flow. In all
three modes the tangential motion is never directed against the ambient
flow, a fact used in the stability analysis. The complete selector,
executed once per control cycle after the estimation step, is:

\begin{quote}
\textbf{Separatrix tracker: one control cycle}

\textbf{Parameters:} $v_{\max}$, $\varepsilon_{\text{raw}}$, $\varepsilon_{\text{dim}}$

\textbf{Given:} $\hat D$, $\hat{\mathbf{g}}$, $\hat{\mathbf{H}}_D$,
$\mathbf{f}$ from Equations~\eqref{eq:J_hat}--\eqref{eq:hess_det}, and
eigenpairs $(\lambda_1,\mathbf{w}_1)$, $(\lambda_2,\mathbf{w}_2)$ of
$\hat{\mathbf{H}}_D$, $\lambda_1\leq\lambda_2$

If the band test (Equation~\eqref{eq:flow_band}) passes (on the
separatrix): apply the flow step (Equation~\eqref{eq:flow_step}).

Else, if $\lambda_1\lambda_2\geq0$ ($\hat D$ has no saddle; near a well of
$D$): apply the attract step (Equation~\eqref{eq:attract_step}).

Else, orient $\mathbf{w}_1$ so that $\hat{\mathbf{g}}^{\mathsf T}\mathbf{w}_1\leq0$
(descent along the trench): if $\mathbf{f}^{\mathsf T}\mathbf{w}_1>0$
(flow agrees with descent), apply the slide step
(Equation~\eqref{eq:slide_step}); otherwise (flow opposes descent), apply
the attract step (Equation~\eqref{eq:attract_step}).
\end{quote}

The implementation is \textit{separatrix\_logic\_c\_step} in
\textit{pentagon\_primitives.py}. Note that the $\mathbf{w}_2$ component
is one and the same in all three laws
(Equations~\eqref{eq:attract_step}, \eqref{eq:slide_step},
and~\eqref{eq:flow_step}): on a trench, $\lambda_2>0$, so the signed and
magnitude denominators agree, and every mode applies the identical
transverse Newton snap. The modes differ only in their motion along the
structure. This shared component is what makes the switched system
tractable in Section~\ref{sec:stability}.

\subsection{Controller 2: Okubo-Weiss Boundary Tracker}
\label{sec:ow_controller}

The boundary tracker solves Problem 2: drive the scalar $\hat D$ to zero
and travel along the level set $\mathcal{C}$. Its transverse action is a
Newton step for the scalar root-finding problem $D(\mathbf{p})=0$: the
displacement $-(\hat D/\lVert\hat{\mathbf{g}}\rVert^2)\hat{\mathbf{g}}$ is
the smallest step that zeroes the linearized $D$, and its length
$|\hat D|/\lVert\hat{\mathbf{g}}\rVert$ is the estimated distance to the
boundary. Its tangential action drifts along the tangent of the level set,
which is perpendicular to the gradient by definition. With the unit
vectors
\begin{equation}
\hat{\mathbf{u}} = -\operatorname{sign}(\hat D)\frac{\hat{\mathbf{g}}}{\lVert\hat{\mathbf{g}}\rVert},
\qquad
\hat{\mathbf{t}} = \pm\,\mathbf{R}_{90°}\frac{\hat{\mathbf{g}}}{\lVert\hat{\mathbf{g}}\rVert},
\label{eq:ow_directions}
\end{equation}
where $\mathbf{R}_{90°}$ is the quarter-turn rotation and the sign of
$\hat{\mathbf{t}}$ is chosen so that $\mathbf{f}^{\mathsf T}\hat{\mathbf{t}}\geq0$
(the team circulates with the flow), the control law is
\begin{equation}
\dot{\mathbf{p}}_c^{\,\text{des}} = \operatorname{sat}\!\left(\frac{|\hat D|}{\lVert\hat{\mathbf{g}}\rVert}\right)\hat{\mathbf{u}} + \operatorname{sat}\!\left(\mathbf{f}^{\mathsf T}\hat{\mathbf{t}}\right)\hat{\mathbf{t}}.
\label{eq:ow_law}
\end{equation}
This has the same radial-plus-tangential structure as the orbital
controller of Chapter~\ref{ch:vector_field}, with the estimated boundary
distance playing the role of the radial error. When
$\lVert\hat{\mathbf{g}}\rVert < \varepsilon_{\text{grad}}$ the directions
Equation~\eqref{eq:ow_directions} are undefined and the law falls back to
drifting with the flow,
$\dot{\mathbf{p}}_c^{\,\text{des}} = \operatorname{sat}(\lVert\mathbf{f}\rVert)\,\mathbf{f}/\lVert\mathbf{f}\rVert$;
Section~\ref{sec:stab_ow} shows this fallback is needed only at isolated
points of $\mathcal{C}$. The implementation is
\textit{logic\_g\_newton\_contour\_pentagon} in
\textit{pentagon\_primitives.py}.

\subsection{Gains and Velocity Limits}

Both controllers saturate every component through
$v_{\max}\tanh(c/v_{\max})$, so the commanded centroid speed never
exceeds $\sqrt2\,v_{\max}$. The primitive-level saturation speed is
$v_{\max}=0.04$~m/s, amplified by a fixed navigation gain $k=3$ before the
cluster layer so that steady commands clear the robots' stiction floor;
the robot-level limits of Equation~\eqref{eq:momentum} are unchanged from
Chapter~\ref{ch:vector_field}. The band thresholds
$\varepsilon_{\text{raw}}=10^{-3}$ and $\varepsilon_{\text{dim}}=0.025$ and
the gradient floor $\varepsilon_{\text{grad}}=10^{-6}$ are fixed in all
experiments unless noted. Table~\ref{tab:params} collects the values used
in all experiments.

\begin{table}[!t]
\centering
\caption{Simulation and controller parameters.}
\label{tab:params}
\footnotesize
\begin{tabular}{llll}
\toprule
\textbf{Symbol} & \textbf{Value} & \textbf{Units} & \textbf{Description} \\
\midrule
$v_{\max}$ & 0.04 & m/s & primitive saturation speed \\
$k$ & 3.0 & -- & navigation gain \\
$v_{\text{max,robot}}$ & 0.3 & m/s & robot speed limit \\
$v_{\text{stiction}}$ & 0.025 & m/s & robot stiction floor \\
$\tau$ & 0.3 & s & actuator time constant \\
$\Delta t$ & 0.1 & s & control period \\
$\alpha_{\text{mom}}$ & 0.717 & -- & momentum coefficient $e^{-\Delta t/\tau}$ \\
$A$ & 0.1 & -- & double-gyre amplitude \\
$\rho$ & 0.075 & m & formation ring radius \\
$\varepsilon_{\text{raw}}$ & $10^{-3}$ & s$^{-2}$ & FLOW-band raw threshold \\
$\varepsilon_{\text{dim}}$ & 0.025 & m$^2$ & FLOW-band dimensionless threshold \\
$\varepsilon_{\text{grad}}$ & $10^{-6}$ & s$^{-2}$m$^{-1}$ & gradient floor (OW law) \\
\bottomrule
\end{tabular}
\end{table}

\section{Stability Analysis}
\label{sec:stability}

\subsection{Closed-Loop Model and Standing Assumptions}

The analysis is carried out at the navigation level: the cluster centroid
is modeled as a kinematic integrator
\begin{equation}
\dot{\mathbf{p}}_c = k\,\boldsymbol{\delta}(\mathbf{p}_c),
\label{eq:kinematic_model}
\end{equation}
where $\boldsymbol{\delta}$ is the output of the separatrix tracker or of
the boundary-tracking law (Equation~\eqref{eq:ow_law}), and $k>0$ is the
navigation gain. This abstracts the cluster and robot layers, whose
tracking dynamics (Equation~\eqref{eq:momentum}) are faster than the
navigation timescale ($\tau=0.3$~s against transit times of tens of
seconds); the same separation was validated in hardware in
Chapter~\ref{ch:vector_field}. Time is rescaled so that $k=1$. Unless
stated otherwise the estimates are taken exact ($\hat D=D$,
$\hat{\mathbf{g}}=\mathbf{g}$, $\hat{\mathbf{H}}_D=\mathbf{H}_D$,
$\mathbf{f}=\mathbf{v}(\mathbf{p}_c)$); Section~\ref{sec:stab_iss} restores
bounded estimation error. Throughout, $\mathbf{w}_1,\mathbf{w}_2$ denote
unit eigenvectors of $\mathbf{H}_D$ with $\lambda_1\leq\lambda_2$, and
$z_k=(\mathbf{w}_k^{\mathsf T}\mathbf{g})/\lambda_k$ the Newton components,
so that the attract step (Equation~\eqref{eq:attract_step}) reads
$\boldsymbol{\delta}=-\sum_k\operatorname{sat}(z_k)\mathbf{w}_k$.

\subsection{The Attract Mode Is a Gradient-Norm Descent}

The first result needs no trench structure at all: it certifies the
attract step globally, wherever $\mathbf{H}_D$ is nonsingular.

\begin{theorem}[Attract-mode certificate]
\label{thm:attract}
Let $V_A=\tfrac12\lVert\mathbf{g}(\mathbf{p}_c)\rVert^2$. Along
Equation~\eqref{eq:kinematic_model} with the attract step
(Equation~\eqref{eq:attract_step}),
\begin{equation}
\dot V_A = -v_{\max}\sum_{k=1}^2 \lambda_k^2\, z_k\tanh\!\left(\frac{z_k}{v_{\max}}\right) \;\leq\; 0,
\label{eq:VA_dot}
\end{equation}
with equality exactly where $\mathbf{g}=\mathbf{0}$. Consequently every
bounded trajectory on which $\mathbf{H}_D$ remains nonsingular converges
to the critical set $\{\mathbf{g}=\mathbf{0}\}$ of $D$, and isolated
critical points with $\mathbf{H}_D$ nonsingular are locally exponentially
stable.
\end{theorem}

\begin{proof}
Since $\nabla V_A=\mathbf{H}_D\mathbf{g}$,
$\dot V_A = \mathbf{g}^{\mathsf T}\mathbf{H}_D\boldsymbol{\delta}
= -\sum_k v_{\max}\tanh(z_k/v_{\max})\,\mathbf{g}^{\mathsf T}\mathbf{H}_D\mathbf{w}_k
= -\sum_k v_{\max}\tanh(z_k/v_{\max})\,\lambda_k(\mathbf{w}_k^{\mathsf T}\mathbf{g})$,
and $\lambda_k(\mathbf{w}_k^{\mathsf T}\mathbf{g})=\lambda_k^2z_k$, giving
Equation~\eqref{eq:VA_dot}; each summand is nonnegative because
$z\tanh z\geq0$. Equality requires $z_1=z_2=0$, that is
$\mathbf{g}=\mathbf{0}$. Convergence to the critical set follows from the
invariance principle. Near an isolated critical point $|z_k|\ll v_{\max}$,
the saturation is inactive at first order,
$\boldsymbol{\delta}\approx-\mathbf{H}_D^{-1}\mathbf{g}\approx-(\mathbf{p}_c-\mathbf{p}^\star_D)$,
a unit-rate linear contraction.
\end{proof}

Theorem~\ref{thm:attract} explains the two roles the attract step plays in
the selector. Far from the trench it provides the approach: it descends
$\lVert\nabla D\rVert^2$ toward the stationary set of $D$, which for the
double gyre consists of the trench line's crest and wells and the gyre
centers. Near a well of $D$, where $\mathbf{H}_D\succ0$, it is the
terminal capture law: the well is the saddle point of the flow, and
convergence is exponential.

\subsection{Separatrix Tracker}
\label{sec:stab_sep}

The trench is formalized as a valley line. The double gyre satisfies all
of the following exactly, with the constants given below.

\begin{assumption}[Nondegenerate straight trench]
\label{ass:trench}
There is a segment $\Gamma$ (take it along the unit vector
$\hat{\mathbf{e}}_\parallel$, with normal $\hat{\mathbf{e}}_\perp$, and
coordinates $(s,n)$) and a tube $T_w=\{|n|\leq w\}$ on which
\begin{equation}
D(s,n) = \varphi(s) + \tfrac12\lambda_\perp(s)n^2, \qquad 0 < \underline{\lambda} \leq \lambda_\perp(s),
\label{eq:normal_form}
\end{equation}
the cross terms vanish
($\hat{\mathbf{e}}_\perp^{\mathsf T}\mathbf{H}_D\hat{\mathbf{e}}_\parallel=0$
on $T_w$), and on $T_w$ the eigenvalue gap satisfies
$\lambda_2-\lambda_1\geq\gamma>0$ wherever the selector is outside its
attract fallback. For the double gyre separatrix,
Equation~\eqref{eq:normal_form} holds with $n=x$, $s=y$,
$\varphi(s)=-\tfrac{\pi^4A^2}{2}(1+\cos2\pi s_f)$ exactly (the field
Equation~\eqref{eq:det_analytic_main} is additively separable), and
$\lambda_\perp=2\pi^6A^2\cos2\pi x_f\geq\underline\lambda$ on any tube with
$w<1/4$.
\end{assumption}

\begin{assumption}[Flow structure on the trench]
\label{ass:flow}
$\Gamma$ is an invariant curve of the flow: $\mathbf{v}$ restricted to
$\Gamma$ is tangent to it, with $\lVert\mathbf{v}\rVert\geq f_{\min}>0$
outside balls of radius $r_0$ around the flow saddles, where the flow
vanishes. Along $\Gamma$, the stationary points of $\varphi$ are of two
kinds only: wells located at the flow saddles, and crests located inside
the FLOW band (Equation~\eqref{eq:flow_band}); quantitatively,
$|\varphi'(s)|\geq m_\varphi>0$ wherever the band condition fails on
$\Gamma$.
\end{assumption}

\begin{lemma}[Common transverse contraction]
\label{lem:transverse}
Under Assumption~\ref{ass:trench}, in every mode of the separatrix tracker
the transverse coordinate obeys
\begin{equation}
\dot n = -v_{\max}\tanh\!\left(\frac{n}{v_{\max}}\right),
\label{eq:n_dynamics}
\end{equation}
so $|n|$ decreases strictly for $n\neq0$, the tube $T_w$ is forward
invariant, and $n\to0$ exponentially once $|n|\leq v_{\max}$. In
particular Equation~\eqref{eq:n_dynamics} is unaffected by arbitrary
switching among the three modes.
\end{lemma}

\begin{proof}
By Equation~\eqref{eq:normal_form} the Hessian is diagonal in
$(\hat{\mathbf{e}}_\parallel,\hat{\mathbf{e}}_\perp)$ with transverse
eigenvalue $\lambda_\perp>0$ and eigenvector $\hat{\mathbf{e}}_\perp$, and
the transverse gradient component is
$\hat{\mathbf{e}}_\perp^{\mathsf T}\mathbf{g}=\lambda_\perp n$. The
$\hat{\mathbf{e}}_\perp$ component of the attract step
(Equation~\eqref{eq:attract_step}), of the slide step
(Equation~\eqref{eq:slide_step}) (the transverse eigenvalue is positive,
so the signed and magnitude denominators agree), and of the flow step
(Equation~\eqref{eq:flow_step}) are all built from the same argument
$-\lambda_\perp n/\lambda_\perp=-n$. Saturation gives
Equation~\eqref{eq:n_dynamics}. The right side is negative for $n>0$ and
positive for $n<0$, so $V_\perp=\tfrac12n^2$ is a common Lyapunov function
for the switched system; for $|n|\leq v_{\max}$,
$\tanh(n/v_{\max})\geq n\tanh(1)/v_{\max}$ yields
$\dot V_\perp\leq-2\tanh(1)V_\perp$.
\end{proof}

\begin{lemma}[Flow-consistent progress]
\label{lem:progress}
Under Assumptions~\ref{ass:trench} and~\ref{ass:flow}, on $\Gamma$ (after
the transient of Lemma~\ref{lem:transverse}) the tangential velocity
component $\dot s$ satisfies, in every mode,
$\dot s\,(\mathbf{v}^{\mathsf T}\hat{\mathbf{e}}_\parallel)\geq0$: the
team never moves against the ambient flow along the trench. Moreover
$|\dot s|\geq m>0$ outside the $r_0$ balls, with
\begin{equation}
m = \min\left\{v_{\max}\tanh\!\left(\tfrac{f_{\min}}{v_{\max}}\right),\; v_{\max}\tanh\!\left(\tfrac{m_\varphi}{\bar\lambda_\parallel v_{\max}}\right)\right\},
\label{eq:progress_rate}
\end{equation}
where $\bar\lambda_\parallel$ bounds $|\lambda_1|$ on the tube.
\end{lemma}

\begin{proof}[Proof of Lemma~\ref{lem:progress}]
On $\Gamma$ the tangential gradient component is $\varphi'(s)$ and the
along-trench eigenvector is $\hat{\mathbf{e}}_\parallel$. In the FLOW
band the tangential command is
$v_{\max}\tanh((\mathbf{f}^{\mathsf T}\mathbf{w}_1)/v_{\max})$ applied
along $\mathbf{w}_1=\pm\hat{\mathbf{e}}_\parallel$; the product form
makes it a nonnegative multiple of the flow's tangential component, and
by Assumption~\ref{ass:flow} the crest region is covered by this mode,
with magnitude at least the first argument of
Equation~\eqref{eq:progress_rate}. Off the band, the slide step is
selected precisely when the flow agrees with the descent direction
$-\operatorname{sign}(\varphi')\hat{\mathbf{e}}_\parallel$, and its
tangential command moves along descent with magnitude
$v_{\max}\tanh(|\varphi'|/(|\lambda_1|v_{\max}))$; the attract step is
selected when the flow opposes descent, and its tangential command,
$-\varphi'/\lambda_1$ with $\lambda_1<0$, moves along ascent, which by
the selector's own condition is again the flow direction, with the same
magnitude. In the attract fallback ($\mathbf{H}_D\succeq0$ near the
wells) the tangential command descends toward the well, which is the
flow direction near a stable-manifold endpoint. In all cases the sign
agrees with the flow and the magnitude is bounded below by
Equation~\eqref{eq:progress_rate} via Assumption~\ref{ass:flow}.
\end{proof}

\begin{theorem}[Separatrix network traversal]
\label{thm:separatrix}
Under Assumptions~\ref{ass:trench} and~\ref{ass:flow}, with exact
estimates, every trajectory of the separatrix tracker starting in the
tube $T_w$ (i) remains in $T_w$ and converges exponentially to the
trench, and (ii) traverses the trench in the direction of the ambient
flow, reaching the $r_0$ neighborhood of a flow saddle in finite time
bounded by $L_\Gamma/m$ with $L_\Gamma$ the trench length. At the
saddle, $D$ attains a local minimum with $\mathbf{H}_D\succ0$
($\mathbf{H}_D=2\pi^6A^2\mathbf{I}$ at the double-gyre saddles), so the
selector's stationary-point test $\lambda_1\lambda_2\geq0$
(Section~\ref{sec:sep_controller}) is satisfied there and (iii)
Theorem~\ref{thm:attract} gives local exponential convergence to the
saddle, terminal capture, exactly as stated. The trench, however, is
not a single arc: Section~\ref{sec:dg_example} shows the saddle is a
crossing of two trench segments running perpendicular to one another.
A selector that reaches the ball and does not satisfy the
stationary-point test, because its estimate of $\mathbf{H}_D$ fails to
certify positive definiteness there, remains in the FLOW or SLIDE modes
of Assumption~\ref{ass:flow} and continues onto the crossing segment
instead of capturing, repeating (i)-(ii) with $\Gamma$ now the new
segment. Capture and continuation are the same theorem applied to
different premises at the crossing; which one the physical controller
exhibits is an estimation question, addressed next.
\end{theorem}

\begin{proof}
(i) is Lemma~\ref{lem:transverse}. (ii): by Lemma~\ref{lem:progress},
$s(t)$ is monotone in the flow direction with rate at least $m$ outside
the $r_0$ balls, and mode switches cannot reverse it, so the ball is
reached in time at most $L_\Gamma/m$. Capture (iii): inside the ball,
$\mathbf{H}_D\succ0$, the selector is in its attract fallback, and
Theorem~\ref{thm:attract} gives local exponential convergence to the
well, which coincides with the flow saddle $\mathbf{p}^*$.
Continuation: if the fallback test fails to trigger despite
$\mathbf{H}_D\succ0$ holding for the true field, the selector instead
falls through to Lemma~\ref{lem:progress}'s flow or descent case on
whichever of the two crossing segments the flow direction and sign
convention of $\hat{\mathbf{g}}$ select, and (i)-(ii) apply unchanged
with $\Gamma$ relabeled.
\end{proof}

\begin{remark}[Role of each mode]
The proof assigns each mode a precise job: the transverse Newton snap
(common to all modes) makes the trench attractive; the slide step
produces the descent along the trench; the flow step carries the team
across the crest, where descent is undefined; the attract step turns
upstream starts around by climbing to the crest, and, when its
stationary-point test fires at a well, doubles as the terminal capture
law. The hybrid has no chattering issue in the transverse channel
because all modes share the dynamics (Equation~\eqref{eq:n_dynamics}).
\end{remark}

\begin{remark}[Capture depends on an estimated, not exact, Hessian]
\label{rem:capture_gap}
Theorem~\ref{thm:separatrix} is stated for exact estimates, under which
$\mathbf{H}_D\succ0$ at the saddle is certified and capture occurs.
Remark~\ref{rem:eigvec} identifies exactly where this fails for the
closed-form estimator: near the wells the biased eigenvalues nearly
vanish and their signs become unreliable, so the stationary-point test
does not fire on $\hat{\mathbf{H}}_D$ even without noise, and the
selector takes the continuation branch of the proof rather than the
capture branch. Section~\ref{sec:results} confirms this is what the
simulated controller does at every saddle it reaches. Terminal capture
is recovered by tightening the test the selector uses, demonstrated as
a one-parameter alternative in Section~\ref{sec:discussion}; both
behaviors are consequences of the same theorem, differing only in
which premise the deployed estimator satisfies at the crossing.
\end{remark}

\subsection{Okubo-Weiss Boundary Tracker}
\label{sec:stab_ow}

\begin{theorem}[Boundary tracking]
\label{thm:ow}
Let $K$ be a compact set on which $0 < g_{\min} \leq \|\mathbf{g}\| \leq
g_{\max}$. Consider (\ref{eq:kinematic_model}) with the law
(\ref{eq:ow_law}), and let $\chi = (\mathbf{g}/\|\mathbf{g}\|)^\top
\hat{\mathbf{t}}$ measure the misalignment between the commanded drift
direction and the true level-set tangent, with $\bar{\chi} = \sup_K
|\chi| < 1$. Then $V_1 = \tfrac{1}{2}D^2$ satisfies
\begin{equation}
    \dot{V}_1 \leq -\|\mathbf{g}\|\,|D|\,v_{\max}\left[
    \tanh\!\left(\frac{|D|}{\|\mathbf{g}\|\,v_{\max}}\right)
    - \bar{\chi}\right].
    \label{eq:V1_dot}
\end{equation}
Consequently:
(a) with exact estimates, $\hat{\mathbf{t}} \perp \mathbf{g}$ by
construction, so $\chi \equiv 0$ and $V_1$ decreases strictly for
$D \neq 0$; $|D|$ falls at rate at least $\tanh(1)\,g_{\min} v_{\max}$
while $|D| \geq \|\mathbf{g}\| v_{\max}$, and thereafter $\dot{V}_1
\leq -2\tanh(1)V_1$, so $D \to 0$ exponentially and trajectories
converge to the boundary $\mathcal{C}$;
(b) if estimate errors misalign the commanded drift ($\bar{\chi} > 0$),
trajectories reach and remain in the residual band
\begin{equation}
    |D| \;\leq\; g_{\max}\, v_{\max}\,
    \operatorname{artanh}(\bar{\chi}).
    \label{eq:ow_residual}
\end{equation}
\end{theorem}

\begin{proof}
Write (\ref{eq:ow_law}) as $\boldsymbol{\delta} = s_\perp
\hat{\mathbf{u}} + s_{\text{tan}} \hat{\mathbf{t}}$ with $s_\perp =
\mathrm{sat}(|D|/\|\mathbf{g}\|) \geq 0$ and $|s_{\text{tan}}| \leq
v_{\max}$. Then $\dot{V}_1 = D\,\mathbf{g}^\top \boldsymbol{\delta}$.
The perpendicular term: $\mathbf{g}^\top \hat{\mathbf{u}} =
-\mathrm{sign}(D)\|\mathbf{g}\|$, so $D\,\mathbf{g}^\top
\hat{\mathbf{u}}\, s_\perp = -|D|\,\|\mathbf{g}\|\, s_\perp$. The
tangential term is bounded by $|D|\,\|\mathbf{g}\|\,\bar{\chi}\,
v_{\max}$. Summing gives (\ref{eq:V1_dot}). For (a), $\tanh z \geq
\tanh(1)\min(z, 1)$ turns (\ref{eq:V1_dot}) into $\dot{V}_1 \leq
-\tanh(1)\min\{D^2,\, |D|\,\|\mathbf{g}\|v_{\max}\}$, giving the two
regimes stated. For (b), the bracket in (\ref{eq:V1_dot}) is positive
whenever $\tanh(|D|/(\|\mathbf{g}\|v_{\max})) > \bar{\chi}$, that is
whenever $|D| > \|\mathbf{g}\|v_{\max}\operatorname{artanh}
(\bar{\chi})$, so $V_1$ decreases outside the band
(\ref{eq:ow_residual}) and the band is attractive and invariant.
\end{proof}

\begin{remark}[Degenerate points of the boundary]
\label{rem:ow_corners}
The gradient floor $\varepsilon_{\text{grad}}$ in the law
(\ref{eq:ow_law}) exists because $\mathcal{C}$ is not globally regular:
the four corners of each double-gyre diamond are critical points of $D$
lying on $\mathcal{C}$, where $\mathbf{g} = \mathbf{0}$ and
Theorem~\ref{thm:ow}'s premise fails. These are isolated points; the
fallback drifts the team with the flow across the corner region, after
which the Newton action resumes. The same fallback covers
initialization at a gyre core's exact center.
\end{remark}

\subsection{Robustness to Estimation Error}
\label{sec:stab_iss}

Finally the exact-estimate assumption is dropped. By
Section~\ref{sec:sensitivity}, noise of level $\sigma_{uv}$ and
truncation over a footprint of radius $\rho$ perturb every quantity the
controllers use; the perturbations are bounded, with constants
determined by the formation. Their effect enters each control channel
inside a $\tanh$, which is 1-Lipschitz, so a perturbation of a
commanded component is no larger than the perturbation of its argument.

\begin{proposition}[Ultimate bounds under bounded estimate error]
\label{prop:iss}
Suppose the estimate errors induce a bounded disturbance $d(t)$, $|d|
\leq \bar{d} < v_{\max}$, in the transverse channel, so that
(\ref{eq:n_dynamics}) becomes $\dot{n} = -v_{\max}\tanh(n/v_{\max}) +
d$. Then the trench remains practically stable with ultimate bound
\begin{equation}
    \limsup_{t\to\infty} |n(t)| \;\leq\;
    v_{\max}\operatorname{artanh}\!\left(\frac{\bar{d}}{v_{\max}}\right),
    \label{eq:iss_bound}
\end{equation}
which is $\approx \bar{d}$ for $\bar{d} \ll v_{\max}$. Likewise, an
error $|\hat{D} - D| \leq \bar{e}_D$ adds $\bar{e}_D/(\|\mathbf{g}\|
v_{\max})$ inside the bracket of (\ref{eq:V1_dot}), inflating the
Okubo-Weiss residual band (\ref{eq:ow_residual}) to
\begin{equation}
    |D| \;\leq\; g_{\max}\,v_{\max}
    \operatorname{artanh}\!\left(\bar{\chi} +
    \frac{\bar{e}_D}{g_{\min} v_{\max}}\right),
    \label{eq:ow_residual_noisy}
\end{equation}
provided the argument is less than one.
\end{proposition}

\begin{proof}
For (\ref{eq:iss_bound}): $\tfrac{d}{dt}\tfrac{1}{2}n^2 = n(-v_{\max}
\tanh(n/v_{\max}) + d) < 0$ whenever $v_{\max}\tanh(|n|/v_{\max}) >
\bar{d}$, that is outside the stated interval, which is therefore
attractive and invariant. For (\ref{eq:ow_residual_noisy}): the
perpendicular channel of (\ref{eq:ow_law}) computed from $\hat{D},
\hat{\mathbf{g}}$ differs from the exact one by terms bounded by
$\bar{e}_D/(\|\mathbf{g}\|v_{\max})$ inside the saturation (Lipschitz
bound), and the argument of Theorem~\ref{thm:ow}(b) repeats with the
bracket shifted by that amount.
\end{proof}

The disturbance levels $\bar{d}$ and $\bar{e}_D$ inherit the structure
of Section~\ref{sec:sensitivity}: a noise part scaling as
$\sigma_{uv}/\rho$ to $\sigma_{uv}/\rho^2$ through Lemma~\ref{lem:gains}
and (\ref{eq:det_perturbation}), and a truncation part scaling as
$M_3\rho$ to $M_3\rho^2$ through (\ref{eq:truncation}). The ultimate
bounds (\ref{eq:iss_bound}) and (\ref{eq:ow_residual_noisy}) therefore
degrade linearly in the measurement noise level, with the formation
radius trading the two contributions as in (\ref{eq:rho_star}). The
Monte Carlo sweeps of Section~\ref{sec:test_plan} quantify these
predictions empirically.

\section{Simulation and Analysis}
\label{sec:methods}

\subsection{Cluster Space Controller and Pentagon Formation}

Both controllers use the same underlying cluster space controller and
pentagon-of-pairs formation. The forward kinematics maps the 12-state
cluster configuration to global robot positions; the inverse kinematics
and analytic $12 \times 12$ inverse kinematic Jacobian (see
Appendix~\ref{app:cluster_kinematics}) map a desired centroid velocity to
individual robot velocity commands while maintaining nominal pair lengths
and the SAS shape of the midpoint triangle.

Formation maintenance is proportional: each of the nine shape-error
terms in $(p_1, \beta_1, q_1, L_2, \psi_2, L_3, \psi_3, L_4, \psi_4)$ is
driven toward its nominal value at a rate proportional to its current
error, with gain $1.0$ on the length terms and $0.1$ on the angle terms
(\textit{position\_gain}, \textit{angle\_gain} in
\textit{pentagon\_cluster.py}), fixed across all experiments.

\subsection{Six-Robot Configuration}

All simulation trials use the same nominal formation
(\textit{pentagon\_small.yaml}), the pentagon-of-pairs of
Appendix~\ref{app:cluster_kinematics} at $0.5\times$ its base scale: midpoint
triangle $p_1 = 0.0806$, $\beta_1 = 0.7726$~rad, $q_1 = 0.1154$; pair
lengths $L_2 = 0.0750$, $L_3 = L_4 = 0.0882$; base pair orientations
$\psi_2 = -\pi/2$, $\psi_3 = -1.885$, $\psi_4 = 1.885$~rad and base
heading $\theta_c = -0.7726$~rad. A trial's initial heading is applied
as a rigid rotation of this base configuration by a single offset
angle, so the formation's shape is fixed and only its orientation
varies across trials. This gives the ring radius $\rho = 0.075$ used
throughout Section~\ref{sec:sensitivity} and Table~\ref{tab:params}.
The centroid $(x_c, y_c)$ is placed at the trial's start point
(Section~\ref{sec:test_plan}); the heading offset is drawn uniformly
from $[0, 2\pi)$ per trial in the Monte Carlo sweeps and left at zero
(the base orientation above) in the clean, noise-free demonstration
runs of Section~\ref{sec:results}.

\subsection{Simulation Program}

Simulations are run using the Python framework in the
\textit{trunk/Python\_Simulations/Vector\_Fields/VF\_Robot/} directory.
For the double gyre, the field is the static case
(\textit{double\_gyre\_static}, amplitude $A = 0.1$) on the domain
$[-1, 1] \times [-0.5, 0.5]$. The integrator is explicit Euler at
10~Hz ($\Delta t = 0.1$~s); robot dynamics follow the first-order
model~(\ref{eq:momentum}) with $\tau = 0.3$~s. Clean, noise-free
demonstration trials run for a budget of several hundred to 1500 steps
and stop early at domain exit or a fixed number of steps past first
saddle or boundary contact (Section~\ref{sec:results}); the noise-sweep trials of
Section~\ref{sec:test_plan} run up to 500 steps and stop at task
completion or domain exit. The estimator-accuracy sweep of
Section~\ref{sec:results} holds the formation static and draws
independent noise per trial rather than stepping the simulator. The
ocean experiment of Section~\ref{sec:ocean_sim} uses the same
integrator and dynamics under the field-time relabeling of
Table~\ref{tab:ocean_params}, running 168 steps. Every experiment
script under \textit{experiments/} states in a header comment which
paper figure or table it produces and where its CSV or PNG output is
written, so the numbers and figures in this chapter trace back to a
single reproducible script each.

\subsection{Ocean HFR Simulation}
\label{sec:ocean_sim}

The ocean experiment uses hourly surface-current frames from the
Santa Barbara Channel HF radar archive (May 16--17, 2012, 29 hourly
frames covering 28~h of real time). The data are the U.S.~West~Coast
Radial-to-Total Vector (RTV) product produced by the Scripps Coastal
Ocean Observing Laboratory (CORDC) HFRNet network and archived at the
NOAA National Centers for Environmental Information under accession
\texttt{gov.noaa.nodc:IOOS-HFRadarRTVector}, retrieved from the NOAA
THREDDS-Ocean server in June 2026. The $2$~km resolution files were
used; each file contains one hour of surface-current velocity
components ($u$, $v$) on a regular geographic grid. Frames are
linearly interpolated in time inside the field evaluator
(\textit{ocean\_hfr\_socal\_timevarying} in
\textit{src/fields/environments/Ocean\_HFR.py}). The cluster runs in
non-dimensional world coordinates on $[-0.65, 0.65]^2$, and an affine
map carries those coordinates to a geographic region of radius
$0.3^{\circ}$ centered on $(34.2^{\circ}\mathrm{N},
-120.4^{\circ}\mathrm{W})$; one world unit corresponds to
approximately $51$~km of latitude at this center. The pentagon
formation is scaled to $0.7\times$ its nominal radius
(\textit{pentagon\_small\_2km.yaml}): the $2$~km product's lower
valid-data coverage is noisier at the nominal scale, and
$0.5\times$ formations were observed to pick the wrong branch at the
bifurcation near $(34.1^{\circ}\mathrm{N}, -120.4^{\circ}\mathrm{W})$;
averaging over a larger footprint recovers the branch a coarser
product tracks.

The simulator's discrete-time step is treated as one sample per ten
minutes of real time (\textit{TIME\_WARP} $= 6000$), rather than the
testbed's native $0.1$~s. Table~\ref{tab:ocean_params} contrasts the
resulting operating point with the laboratory Decabot values of
Table~\ref{tab:params}. Two changes follow directly from the time
dilation. First, the momentum coefficient is set to
$\alpha_{\text{mom}} = 0$ (\textit{MOMENTUM\_ALPHA}): the Decabot value
$\alpha_{\text{mom}} = 0.717$ represents a $0.3$~s actuator lag at the
testbed's $10$~Hz control rate, but at one control cycle per ten
real-world minutes any physically reasonable vessel time constant is
negligible in comparison, so the cluster is given full command
authority every step rather than carrying spurious momentum across a
ten-minute gap. Second, the stiction floor is lowered from the
Decabot value ($0.025$ world units per step) to \textit{STICTION\_THRESHOLD}
$= 0.002$: the Decabot floor models motor static friction, which has no
vessel-scale analogue, and at the chosen operating point the original
value would incorrectly clip small commanded corrections. The
navigation gain is retuned from $k = 3.0$ to $k = 1.8$, with
$v_{\max} = 0.04$; empirically this lets the tracker follow the ridge
branch that descends along the middle island's north coast and makes
landfall there, a separate branch choice from the formation-scale one
above, rather than continuing east past it. With $168$ sample steps the trial spans
$168 \cdot 10$~min $= 28$~h of real field time, the full duration of
the archived record.

\begin{table}[!t]
\centering
\caption{Ocean HFR operating point versus the laboratory Decabot values of Table~\ref{tab:params}.}
\label{tab:ocean_params}
\footnotesize
\begin{tabular}{lrr}
\toprule
\textbf{Quantity} & \textbf{Decabot} & \textbf{Ocean HFR} \\
\midrule
$v_{\max}$ (world units/step) & 0.04 & 0.04 \\
$k$ & 3.0 & 1.8 \\
$\alpha_{\text{mom}}$ & 0.717 & 0 \\
stiction floor (world units/step) & 0.025 & 0.002 \\
seconds per step & 0.1 & 600 \\
formation scale & $1\times$ & $0.7\times$ \\
\bottomrule
\end{tabular}
\end{table}

\subsection{Noise Model}
\label{sec:noise_model}

Two independent noise channels perturb the samples available to the
estimator. Measurement noise is additive Gaussian on the $(u, v)$
components read by each robot:
\begin{equation}
    u_i^m = u(\mathbf{p}_i) + \eta_i^u, \quad
    v_i^m = v(\mathbf{p}_i) + \eta_i^v,
    \quad \eta \sim \mathcal{N}(0, \sigma_{uv}^2),
    \label{eq:meas_noise}
\end{equation}
with samples independent across robots, components, and time steps.
Position noise perturbs the location at which each robot samples the
field,
\begin{equation}
    (u_i^m, v_i^m) = \mathbf{v}(\mathbf{p}_i + \boldsymbol{\xi}_i), \quad
    \boldsymbol{\xi}_i \sim \mathcal{N}(\mathbf{0}, \sigma_p^2 \mathbf{I}),
    \label{eq:pos_noise}
\end{equation}
while the quadratic fit (\ref{eq:solve}) uses the nominal positions
$\mathbf{p}_i$: the robot reads the field at a slightly wrong place. To
first order this induces an effective measurement error
$\mathbf{J}(\mathbf{p}_i)\,\boldsymbol{\xi}_i$, which couples the two
channels through the local Jacobian magnitude and is tested directly by
the two-dimensional noise sweep of Section~\ref{sec:test_plan}.

\subsection{Testing Plan}
\label{sec:test_plan}

Both controllers are evaluated in Monte Carlo over a two-dimensional
noise grid, from one consistent initial condition per controller, so
that the success rate measures noise robustness alone rather than basin
geometry (Section~\ref{sec:limitations}). The separatrix tracker starts
from a formation straddling the separatrix at $(0, 0.35)$; the boundary
tracker starts adjacent to the left Okubo-Weiss boundary at
$(-0.5, 0.25)$. Every trial draws an independent initial cluster heading
uniformly from $[0, 2\pi)$; by Lemma~\ref{lem:gains} the estimator gains
are heading-isotropic, and the trials test this prediction.
\begin{itemize}
    \item $\sigma_{uv} \in \{0,\, 0.001,\, 0.005,\, 0.01,\, 0.02,\,
          0.05,\, 0.1\}$ m/s; the top level is doubled and re-swept
          until the worst-cell success rate falls below 10\%, so the
          failure cliff is always bracketed.
    \item $\sigma_p \in \{0,\, 0.005,\, 0.01,\, 0.02,\, 0.05\}$ m.
    \item 1000 trials per cell, at 10 Hz for at most 500 steps; a trial
          stops early when the separatrix tracker reaches the terminal
          saddle, or when the centroid leaves the domain (the boundary
          tracker rides its zero line out of the domain by design).
\end{itemize}

The following are recorded for each trial:
\begin{enumerate}
    \item \textbf{Task success.} For the separatrix tracker, reaching
          the terminal saddle (centroid within 0.06 of
          $\mathbf{p}^*_1$); for the boundary tracker, locking onto the
          $D = 0$ set ($|D| < 0.02$ sustained for 10 steps) with
          95th-percentile $|D|$ thereafter below 0.10. A trial whose
          formation collapses (below) fails regardless.
    \item \textbf{Time-to-band:} the first step at which the tracking
          quantity ($|x_c|$ for the separatrix, $|D(\mathbf{p}_c)|$ for
          the boundary) drops below its threshold and remains below it
          for 10 consecutive steps.
    \item \textbf{Straddle retention:} whether robots remain on both
          sides of the separatrix from the first straddling step until
          the stop, the failure count of [34].
    \item \textbf{Tracking error:} mean and 95th percentile of the
          tracking quantity over the on-structure phase (band entry to
          stop).
    \item \textbf{Formation shape error:} RMS error of the fifteen
          inter-robot distances against their nominal values; an
          excursion beyond twice the nominal pair length is a formation
          collapse.
    \item \textbf{Control effort:} $\int_0^T \|\mathbf{v}^{\text{des}}(t)\|
          \,dt$ over the trial.
\end{enumerate}

Success rates are reported as the fraction of successful trials in each
$(\sigma_{uv}, \sigma_p)$ cell.

\section{Results}
\label{sec:results}

\subsection{Estimator Accuracy versus Noise}

% %% TODO: SYNC WITH LATEST DRAFT -- figure needs generating/copying for
% thesis (source: experiments/estimator_accuracy_sweep.py and
% experiments/plot_estimator_accuracy.py).
\begin{figure}[!t]
    \centering
    \framebox[0.9\linewidth]{\rule{0pt}{3.0cm}}
    \caption{Relative estimation error of $\hat{D}$, $\nabla\hat{D}$, and
    $\hat{\mathbf{H}}_D$ (median over $2 \times 10^5$ noise draws per point;
    shaded interquartile range in (a)) for the steady double gyre of
    Section~\ref{sec:dg_example}, formation held static at the
    strain-region point $\mathbf{p}_c = (-0.35, -0.20)$. (a) Versus measurement noise $\sigma_{uv}$ at the nominal
    formation radius $\rho = 0.075$; the dashed line marks error equal to
    signal magnitude. (b) Versus formation radius $\rho$ at $\sigma_{uv} =
    0.01$; dotted curves are the noise-free truncation floors. $\hat{D}$
    and $\nabla\hat{D}$ exhibit the radius trade-off of (\ref{eq:rho_star});
    $\hat{\mathbf{H}}_D$ is limited by the structural model bias of
    Section~\ref{sec:sensitivity}.}
    \label{fig:est_accuracy}
\end{figure}

Estimator accuracy was measured open loop: the formation was held static
at the strain-region point $\mathbf{p}_c = (-0.35, -0.20)$ of the steady
double gyre, $2 \times 10^5$ independent noise draws were applied per
condition through the same fit used by the controllers, and the estimates
were compared with the analytic values of Appendix~\ref{app:double_gyre}.
The empirical standard deviation of every fitted coefficient matched the
closed forms of Lemma~\ref{lem:gains} within 0.2\%, and the per-robot
reading error induced by position noise matched the first-order
prediction of Section~\ref{sec:noise_model} within 1\%, confirming that
the two noise channels combine into a single effective measurement noise.

Figure~\ref{fig:est_accuracy} reports the accuracy of the three navigation
quantities as relative error, the error magnitude divided by the
magnitude of the true quantity; a value of one (dashed line) means the
estimate is as wrong as it is large and carries no directional
information. In Figure~\ref{fig:est_accuracy}(a) the three curves rise
linearly with $\sigma_{uv}$ and sit stacked a constant factor apart: each
additional derivative order multiplies the noise gain by roughly
$1/\rho$, about 13 at the nominal radius, the ladder of
Lemma~\ref{lem:gains}. Reading off the crossings of the dashed line,
$\hat{D}$ remains informative up to $\sigma_{uv} \approx 0.06$ m/s and
$\nabla\hat{D}$ up to $\sigma_{uv} \approx 0.01$ m/s, while
$\hat{\mathbf{H}}_D$ sits at the line even for the quietest sensors:
below $\sigma_{uv} \approx 0.005$ its error is set by the structural
model bias of Section~\ref{sec:sensitivity}, not by noise.
Figure~\ref{fig:est_accuracy}(b) fixes $\sigma_{uv} = 0.01$ and varies the
formation radius. Enlarging the ring averages noise over a longer
baseline, so the solid curves fall; the quadratic model simultaneously
fits the curved field more poorly, so the noise-free truncation floors
(dotted) rise. The compromise of (\ref{eq:rho_star}) appears as minima
at $\rho \approx 0.11$ for $\hat{D}$ and $\rho \approx 0.23$ for
$\nabla\hat{D}$. These optima are properties of the double gyre at this
location and noise level, entering through the third-derivative scale
$M_3$; the one-third-power scaling law is general, but a different
field, or a different region of the same field, shifts them.
$\hat{\mathbf{H}}_D$ exhibits no minimum: its floor is radius
independent, as predicted, and no formation size improves it, though its
eigen-directions remain accurate to $0.35°$ on and near the
separatrix and $3.8°$ at the generic test point
(Remark~\ref{rem:eigvec}).

The gradient threshold locates where closed-loop degradation should
begin: at $\sigma_{uv} \approx 0.01$ a single $\nabla\hat{D}$ estimate
is as wrong as it is right, so any one control decision approaches
chance. The closed loop does not extend this limit: the Monte Carlo
sweeps that follow place the 50\% success level of both controllers at
$\sigma_{uv} \approx 0.01$, coinciding with the single-estimate
threshold. Saturation and the per-cycle redraw of the noise bound the
damage of individual bad commands, but they buy no additional margin
for a traverse that requires a long sequence of correct decisions.

\subsection{Separatrix Controller Simulation Results}

% %% TODO: SYNC WITH LATEST DRAFT -- figure needs generating/copying for
% thesis (source: experiments/separatrix_clean_runs.py).
\begin{figure}[!t]
    \centering
    \framebox[0.9\linewidth]{\rule{0pt}{3.0cm}}
    \caption{Separatrix tracker (Logic C) centroid trajectories from six
    initial conditions, noise-free case, overlaid on the double-gyre
    streamlines and the two Okubo-Weiss diamonds. Start points are
    $\bullet$, saddles are $\times$. Every trajectory reaches $x = 0$ and
    continues along the wall trench through the terminal (bottom) saddle
    in both directions; a trial stops at domain exit or 150 steps after
    first saddle contact.}
    \label{fig:sep_trajectories}
\end{figure}

From the six starts of Figure~\ref{fig:sep_trajectories}, all six reach the
separatrix and pass within 0.02 of the bottom saddle
$\mathbf{p}^*_1$ (closest approach 0.0003--0.0194 across runs). Three of
the six begin already inside the FLOW band and band immediately;
the other three take 6, 18, and 35 steps. Mode occupancy over each run
is FLOW 6--37\%, SLIDE 34--48\%, and ATTRACT 10--24\%, with the
remaining time split between the flow-drift and attract-fallback
refinements of the same three modes; no run spends more than a few
percent of its steps outside these categories. Every run continues past
the terminal saddle rather than stopping there. This traces to the
structural bias of $\hat{\mathbf{H}}_D$ (Section~\ref{sec:sensitivity}):
near the well the estimated eigenvalues are indefinite and roughly an
order of magnitude too small, so the capture test
$\lambda_1\lambda_2 \geq 0$ of Section~\ref{sec:sep_controller} does not
fire even without noise. Section~\ref{sec:discussion} returns to this
behavior, continuous traversal of the separatrix network, with terminal
capture demonstrated as a tightened-threshold special case.

\begin{table}[!t]
\centering
\caption{Separatrix tracker: success rate (traverse) by noise cell,
1000 trials/cell, fixed straddling start $(0, 0.35)$.}
\label{tab:sep_success}
\footnotesize
\begin{tabular}{rrrrrr}
\toprule
$\sigma_{uv}$ \textbackslash\ $\sigma_p$ & 0.000 & 0.005 & 0.010 & 0.020 & 0.050 \\
\midrule
0.000 & 100.0 & 64.7 & 45.4 & 36.0 & 18.8 \\
0.001 &  99.3 & 61.9 & 43.7 & 33.8 & 19.9 \\
0.005 &  60.2 & 53.6 & 40.7 & 29.7 & 17.4 \\
0.010 &  43.3 & 40.6 & 39.6 & 28.1 & 18.4 \\
0.020 &  27.6 & 27.5 & 22.9 & 20.4 & 18.4 \\
0.050 &  15.1 & 16.4 & 14.7 & 15.2 & 15.0 \\
0.100 &  13.0 & 15.0 & 13.3 & 14.3 & 13.5 \\
\bottomrule
\end{tabular}
\end{table}

Table~\ref{tab:sep_success} reports traverse success (reaching the
terminal saddle without formation collapse) over the noise grid. At
$\sigma_p = 0$ the rate is a clean step: 100\% through
$\sigma_{uv} = 0.001$, then falling through the 50\% level between
$\sigma_{uv} = 0.005$ and $0.01$, landing at 43.3\% at
$\sigma_{uv} = 0.01$. This coincides with the open-loop
$\nabla\hat{D}$ reliability threshold reported above
($\sigma_{uv} \approx 0.01$, ratio 1.0): saturation and per-cycle
resampling do not extend the noise budget
of a long decision sequence beyond what a single gradient estimate can
support. Position noise degrades the rate along the same curve scaled
by $\|\mathbf{J}\|$, consistent with the effective-noise argument of
Section~\ref{sec:noise_model}: $\sigma_p = 0.01$ alone gives 45.4\%,
close to $\sigma_{uv} = 0.01$ alone at 43.3\%. Beyond the cliff, success
does not fall to zero but levels off near 13--15\%: at this noise level
the controller behaves as an unbiased random walk that occasionally
lands on the saddle by chance, so the straddle-retention metric of
Table~\ref{tab:sep_straddle} is the sharper failure indicator. Tracking
error over the on-trench phase (mean $|x_c|$) rises from 0.002 m at zero
noise to 0.086 m at the $\sigma_{uv} = 0.01$ cliff and saturates near
0.27 m at high noise (comparable to the domain half-width), reflecting
trials that never band at all. Success was independent of initial
heading to within the sampling noise of the sweep (quartile success
74.9--76.2\% at low noise, confirming the isotropy of
Lemma~\ref{lem:gains}).

\begin{table}[!t]
\centering
\caption{Separatrix tracker: straddle retention by noise cell (fraction
of 1000 trials, Michini-comparable metric [34]).}
\label{tab:sep_straddle}
\footnotesize
\begin{tabular}{rrrrrr}
\toprule
$\sigma_{uv}$ \textbackslash\ $\sigma_p$ & 0.000 & 0.005 & 0.010 & 0.020 & 0.050 \\
\midrule
0.000 & 1.000 & 0.625 & 0.417 & 0.338 & 0.164 \\
0.005 & 0.560 & 0.493 & 0.365 & 0.272 & 0.157 \\
0.010 & 0.379 & 0.366 & 0.361 & 0.258 & 0.165 \\
0.020 & 0.237 & 0.232 & 0.187 & 0.177 & 0.159 \\
0.050 & 0.126 & 0.137 & 0.112 & 0.128 & 0.122 \\
\bottomrule
\end{tabular}
\end{table}

\subsection{Okubo-Weiss Boundary Tracker Simulation Results}

% %% TODO: SYNC WITH LATEST DRAFT -- figure needs generating/copying for
% thesis (source: experiments/ow_clean_runs.py).
\begin{figure}[!t]
    \centering
    \framebox[0.9\linewidth]{\rule{0pt}{3.0cm}}
    \caption{Okubo-Weiss boundary tracker centroid trajectories from four
    initial conditions, noise-free case, overlaid on $D(\mathbf{p})$
    (red $D>0$ cores, blue $D<0$ exterior) and the diamond boundaries.
    Gyre centers are $+$, diamond corners are $\times$. After lock-on,
    each trajectory follows its diamond edge straight through the
    corners onto the neighboring edge, consistent with the line-network
    structure of $\mathcal{C}$ (Section~\ref{sec:dg_example}).}
    \label{fig:ow_trajectories}
\end{figure}

All four starts of Figure~\ref{fig:ow_trajectories} lock onto the
boundary in 13--16 steps and hold it: mean $|D|$ on-boundary is
0.0032--0.0075 (0.3--0.8\% of the well depth $\pi^4 A^2 = 0.974$), with
95th-percentile $|D|$ of 0.028--0.039. Two of the four runs each cross one
diamond corner and dwell there for 4--8 steps before continuing on the
new edge; the boundary tracker never turns to circulate a single
diamond, it travels along the connected line network, matching Problem
2's statement in Section~\ref{sec:problem_statement}.

\begin{table}[!t]
\centering
\caption{Okubo-Weiss tracker: success rate (track) by noise cell, 1000
trials/cell, fixed start $(-0.5, 0.25)$.}
\label{tab:ow_success}
\footnotesize
\begin{tabular}{rrrrrr}
\toprule
$\sigma_{uv}$ \textbackslash\ $\sigma_p$ & 0.000 & 0.005 & 0.010 & 0.020 & 0.050 \\
\midrule
0.000 & 100.0 & 100.0 & 100.0 &  98.6 &   9.7 \\
0.001 & 100.0 & 100.0 & 100.0 &  98.4 &   9.5 \\
0.005 &  93.8 &  93.7 &  91.6 &  72.8 &   8.1 \\
0.010 &  20.9 &  17.7 &  16.7 &  13.4 &   9.1 \\
0.020 &   8.8 &   8.0 &   9.6 &   7.2 &   7.9 \\
0.050 &   6.6 &   5.5 &   4.4 &   6.7 &   6.0 \\
\bottomrule
\end{tabular}
\end{table}

Table~\ref{tab:ow_success} shows a sharper cliff than the separatrix
tracker: full tolerance through $\sigma_{uv} = 0.001$ regardless of
$\sigma_p$, 93.8\% at $\sigma_{uv} = 0.005$, then a collapse to 20.9\%
at $\sigma_{uv} = 0.01$, the same coordinate as the separatrix
controller's cliff and the same open-loop gradient threshold. Position
noise is tolerated to $\sigma_p = 0.02$ at low $\sigma_{uv}$ (98.6\% at
$\sigma_{uv} = 0$) but the $\sigma_p = 0.05$ column is uniformly under
10\% regardless of $\sigma_{uv}$, since $\sigma_p = 0.05$ alone already
exceeds the reading-error budget the boundary law needs to keep
$|\hat{D}|$ below its lock-on threshold. On-boundary tracking error
($|D(\mathbf{p}_c)|$, mean and 95th percentile) is essentially unchanged
through $\sigma_{uv} = 0.001$ (mean 0.0008--0.0029) and rises steadily
to the well depth scale by $\sigma_{uv} = 0.05$ (mean 0.21, 95th
percentile 0.50). As with the separatrix tracker, success was
independent of initial heading (quartile success 77.6--79.4\% at low
noise).

\subsection{Comparison of the Two Controllers}

Both controllers place their 50\% success level at the same
$\sigma_{uv} \approx 0.01$, the point at which a single
$\nabla\hat{D}$ estimate is as wrong as it is right, reported above in
the estimator-accuracy results. Neither the flow-assisted
separatrix law nor the boundary law's tangential drift extends this
budget: the estimator, not the control law, sets the noise ceiling for
both structures. The boundary tracker's cliff is sharper (100\% to
20.9\% between $\sigma_{uv} = 0.005$ and $0.01$, versus the separatrix
tracker's more gradual 60.2\% to 43.3\% over the same interval) because
tracking $\mathcal{C}$ is a codimension-one regular-level-set problem
that fails outright once $\nabla\hat{D}$ stops resolving the crossing
direction, while the separatrix tracker degrades more gradually because
its FLOW mode can still make progress along the trench using the
ambient flow estimate $\mathbf{f}$ even when $\hat{\mathbf{H}}_D$ is
unreliable. This mirrors the asymmetry in the stability results of
Section~\ref{sec:stability}: the boundary is a single regular curve with
a global certificate; the separatrix is a trench whose capture depends
on curvature the estimator recovers only approximately.

\subsection{Ocean HFR Simulation}
\label{sec:results_ocean}

The separatrix tracker was run on the real, time-varying current field
of Section~\ref{sec:ocean_sim} from a single start north of the Santa
Barbara Channel entrance, $(34.40°\mathrm{N},
-120.39°\mathrm{W})$, using the operating point of
Table~\ref{tab:ocean_params}. Figure~\ref{fig:ocean_progression} shows the
path building up over the full 28-h trial against the instantaneous
current field at five points in time: the cluster descends through the
channel entrance, tracks the ridge through the bifurcation near
$(34.1°\mathrm{N}, -120.4°\mathrm{W})$, and continues to
landfall on the middle island's north coast. Figure~\ref{fig:ocean_ftle}
places the same trial against the 24-h forward-time FTLE field, computed
offline over the same window; the cluster's path
follows the dominant FTLE ridge running from the channel entrance to the
bifurcation closely, without the ridge being available to the controller
at any point, only the local $\hat{D}$ and $\hat{\mathbf{g}}$ estimates
that Section~\ref{sec:estimation} builds from the current samples at the
formation.

% %% TODO: SYNC WITH LATEST DRAFT -- figure needs generating/copying for
% thesis (source: experiments/ocean_hfr_2km_panel_progression.py).
\begin{figure}[!t]
    \centering
    \framebox[0.9\linewidth]{\rule{0pt}{3.0cm}}
    \caption{Centroid path so far (black) over the instantaneous current
    field (quiver), at $t = 0, 7, 14, 21, 28$~h into the 28-h trial. Green
    star is the start; red circle is the current position at each panel's
    time. The cluster descends through the channel entrance and turns
    toward the middle island as the current field evolves.}
    \label{fig:ocean_progression}
\end{figure}

% %% TODO: SYNC WITH LATEST DRAFT -- figure needs generating/copying for
% thesis (source: experiments/main_ocean_hfr_2km_ftle_overlay.py).
\begin{figure}[!t]
    \centering
    \framebox[0.9\linewidth]{\rule{0pt}{3.0cm}}
    \caption{Full 28-h centroid and robot paths overlaid on the 24-h
    forward-time FTLE field computed from the same 2~km data. The path
    tracks the dominant ridge from the channel entrance to the
    bifurcation and continues to landfall near
    $(34.04°\mathrm{N}, -120.26°\mathrm{W})$.}
    \label{fig:ocean_ftle}
\end{figure}

Figure~\ref{fig:ocean_ftle_snapshots} computes the same 24-h-style
forward FTLE anew at four anchor times spread through the record
(0, 5, 10, and 16~h, each integrated 12~h forward, the longest common
window the 29-frame archive supports at more than two anchors), showing
that the dominant ridge through the channel entrance persists across
the record while its finer structure downstream continues to evolve.
This is the sense in which the ocean experiment is genuinely
time-varying, unlike the steady double gyre of the rest of the chapter:
the structure the controller tracks is not fixed in advance and the
controller has no access to its future evolution, only the
instantaneous samples at the current formation position.

% %% TODO: SYNC WITH LATEST DRAFT -- figure needs generating/copying for
% thesis (source: experiments/ocean_hfr_2km_ftle_snapshots.py).
\begin{figure*}[!t]
    \centering
    \framebox[0.9\linewidth]{\rule{0pt}{3.0cm}}
    \caption{Forward-time FTLE (12-h window) anchored at four times
    across the 28-h record. The dominant ridge through the channel
    entrance is stable across all four anchors; downstream structure
    continues to evolve, confirming the field is genuinely time-varying
    over the trial.}
    \label{fig:ocean_ftle_snapshots}
\end{figure*}

Figure~\ref{fig:ocean_branch} tests the trial's sensitivity to the start
position with a $10 \times 10$ grid of starts jittered
$\pm 0.05°$ around the validated start, each run through the
full 28-h trial. Of the 100 starts, 84 land on the same
middle-island branch as the validated trial; the remaining 16, all at
the edges of the jitter box furthest from the ridge, are carried onto a
western branch that does not make landfall in the domain. The mean
distance from a start to the 95th-percentile FTLE ridge contour was
1.4~km, small relative to the $\sim 0.3°$ ($\sim 33$~km) domain
radius, so the branch boundary sits close to the jitter box: most of
the box is on one side of it. This is a basin-of-attraction result in
the sense of Section~\ref{sec:limitations}, not a robustness guarantee;
it shows that the tracked branch is the typical outcome for starts near
the validated one, not that every start converges to it.

% %% TODO: SYNC WITH LATEST DRAFT -- figure needs generating/copying for
% thesis (source: experiments/ocean_hfr_2km_branch_sensitivity.py).
\begin{figure}[!t]
    \centering
    \framebox[0.9\linewidth]{\rule{0pt}{3.0cm}}
    \caption{Landfall branch for a $10 \times 10$ grid of starts jittered
    around the validated start, over the 24-h forward FTLE background.
    84 of 100 starts (blue) reach the same middle-island branch as
    Figure~\ref{fig:ocean_ftle}; 16 (red), all near the jitter box edges,
    are carried onto a western branch instead.}
    \label{fig:ocean_branch}
\end{figure}

\section{Discussion}
\label{sec:discussion}

\subsection{Implications for Distributed Eulerian-Structure Following}

The two stability results of Section~\ref{sec:stability} differ in
strength for a structural reason. The Okubo-Weiss boundary is a
codimension-one regular level set: first-order information, a gradient
that does not vanish, fully determines the transversal direction and
yields a global certificate from a single quadratic Lyapunov function.
The separatrix trench lives where that first-order information vanishes
in the normal direction, so any controller and any proof must
synthesize the missing direction from curvature; the eigenstructure
dependence, mode switching, and tube assumptions of
Section~\ref{sec:stab_sep} are the price of that synthesis. The
separatrix result is the more difficult robotics contribution precisely
because the problem is harder; the boundary result is the stronger
theorem precisely because the problem is easier.

The results of Section~\ref{sec:results} show that a six-robot formation
sampling only its own local velocities recovers enough second-order
structure, not the field itself, to find and follow both kinds of
Eulerian feature, with no map, no forecast, and no communication beyond
the single synchronous sample each control cycle already requires. The
$\det(\mathbf{J})$ landscape is what makes this possible: both problems
reduce to following a scalar surrogate field that is computable from
the same twelve fitted coefficients, a trench for the separatrix and a
zero-crossing for the boundary, rather than to any procedure specific
to streamlines or Lagrangian integration. The two controllers are
consequently complementary rather than competing: one tracks a curve
defined by where curvature information vanishes to first order and must
be supplied by the Hessian, the other a curve defined by where the
Hessian itself is unnecessary because the gradient alone is regular
everywhere on it. A team that could estimate both could, in principle,
select which structure to follow from the local field alone, without an
operator choosing in advance.

\subsection{Park versus Traverse: One Threshold, Two Behaviors}
\label{sec:park_vs_traverse}

Theorem~\ref{thm:separatrix} and Remark~\ref{rem:capture_gap} identify
the terminal-capture test $\lambda_1\lambda_2 \geq 0$ as the reason the
deployed controller continues through every saddle in
Section~\ref{sec:results}: the biased $\hat{\mathbf{H}}_D$ makes that
test unreliable exactly where it would need to fire. An estimator-aware
alternative restores capture without touching the transverse Newton
snap, the flow step, or the slide step, using only the two quantities
Section~\ref{sec:sensitivity} shows are linear in the fitted
coefficients and free of the Hessian's structural bias, $\hat{D}$ and
$\hat{\mathbf{g}}$. The \emph{park step} adds a fourth mode, checked
before the FLOW test:
\begin{equation}
    \text{if } \hat{D} < D_{\text{capture}}:\quad
    \dot{\mathbf{p}}_c^{\,\text{des}} =
    -\begin{bmatrix}
        \mathrm{sat}(\hat{g}_x) \\ \mathrm{sat}(\hat{g}_y)
    \end{bmatrix},
    \label{eq:park_step}
\end{equation}
componentwise saturated gradient descent on $\hat{D}$, which converges
to the nearest local minimum of $D$, the flow saddle at the end of the
current trench segment, exactly as Theorem~\ref{thm:attract} predicts
for the exact-estimate case. $D_{\text{capture}}$ is the one-parameter
knob: $D_{\text{capture}} = -\infty$ disables the test and the
controller is trajectory-identical to Logic C (verified in
\textit{separatrix\_logic\_park\_step}, bit-exact off), while a finite
threshold inside the well (here $-0.5$, roughly half the double-gyre
well depth $\pi^4 A^2 \approx 0.974$) parks the team there.
Figure~\ref{fig:park_vs_traverse} runs both settings from the same start,
noise-free: traverse exits the domain at $(1.00, -0.51)$ after 228
steps, riding the crossing trench as Section~\ref{sec:results}
describes; park reaches the saddle and holds position with zero
tail-position standard deviation over its final 100 steps, settling
0.0124~m from $\mathbf{p}^*_1$. Both runs use the identical estimation
step and the identical selector below the park test; the only
difference is one threshold. This is the practical resolution of
Remark~\ref{rem:capture_gap}: continuation is not a controller limitation
to be designed around so much as the honest behavior of the theorem's
capture test under a real estimator, and a single threshold recovers
whichever behavior a mission calls for, terminal parking at a critical
point or open-ended traversal of the network.

% %% TODO: SYNC WITH LATEST DRAFT -- figure needs generating/copying for
% thesis (source: experiments/park_vs_traverse_demo.py).
\begin{figure}[!t]
    \centering
    \framebox[0.9\linewidth]{\rule{0pt}{3.0cm}}
    \caption{Park versus traverse from the same start $(0.10, -0.20)$,
    noise-free. Traverse (Logic C) rides the trench through the bottom
    saddle and exits the domain. Park (Logic Park,
    $D_{\text{capture}} = -0.5$) applies the park step
    (\ref{eq:park_step}) once $\hat{D}$ drops below threshold and holds
    the team 0.0124~m from the saddle.}
    \label{fig:park_vs_traverse}
\end{figure}

\subsection{Failure Modes}

Three behaviors fall outside the guarantees of
Section~\ref{sec:stability} and are worth stating plainly rather than
leaving implicit. First, the estimator is ill-conditioned whenever the
six sampling points drift toward a common conic
(Proposition~\ref{prop:conic}): $\det(\boldsymbol{\Phi})$ shrinks
continuously as the formation deforms toward one, and the fitted
coefficients degrade smoothly until, at the conic itself, the fit is
undefined. A ring formation without a center robot is the standing
example (Remark following Corollary~\ref{cor:pentagon}); the pentagon
plus center used throughout is one robot away from this failure by
construction, but a controller could in principle drive the formation
toward a degenerate shape through rotation or scaling commands not
considered here. Monitoring $\det(\boldsymbol{\Phi})$, or equivalently
the interior angles of the formation, and adding a shape-maintenance
term to the centroid command is the direct mitigation; the Monte Carlo
trials of Section~\ref{sec:results} flag the weaker symptom, an
inter-robot distance excursion beyond twice the nominal pair length, as
a formation collapse, and it does not occur in any reported trial.

Second, the separatrix controller can chatter across the FLOW-band
boundary (Equation~\eqref{eq:flow_band}) when the centroid sits near the threshold
and noise pushes $\hat{D}$ or the dimensionless ratio
$|\hat{D}|/\|\hat{\mathbf{H}}_D\|_F$ back and forth across
$\varepsilon_{\text{dim}}$. Because every mode shares the transverse
dynamics (Equation~\eqref{eq:n_dynamics}), chatter does not destabilize the
transverse channel; its cost is tangential, a lower effective
along-trench speed as the command alternates between the flow step and
the slide or attract step. $\varepsilon_{\text{dim}}$ trades this
against the crest crossing of Section~\ref{sec:dg_example}: too small a
threshold reintroduces the stall the flow-assisted mode was built to
remove, too large a threshold widens the band and the chatter with it.
The value used throughout, $\varepsilon_{\text{dim}} = 0.025$, was fixed
once and not retuned per experiment.

Third, the Okubo-Weiss boundary tracker's tangential drift can carry the
team through a diamond corner and onto the neighboring line of the
periodic network (Section~\ref{sec:dg_example}), rather than turning to
circulate the diamond it started on, which is what
Section~\ref{sec:results} reports rather than the circulation Problem 2
originally anticipated. The mechanism is the same corner behavior
covered by Remark~\ref{rem:ow_corners} following Theorem~\ref{thm:ow},
extended past the corner itself: once the flow fallback carries the team
onto the new line, the Newton action resumes there and the team departs
along it, and because the analytic double-gyre field tiles the plane
with no domain fence, this can carry a simulated team arbitrarily far
from the diamond it started on. On real, non-periodic fields this
would not recur past the first crossing encountered; in the simulation
it is worth flagging explicitly rather than truncating trials at the
domain edge without comment, which is what the Testing Plan does
(Section~\ref{sec:test_plan}).

\section{Limitations}
\label{sec:limitations}

This chapter restricts attention to the static double gyre ($\varepsilon = 0$).
In the time-varying case ($\varepsilon > 0$) the separatrix becomes a
material curve that evolves with the flow, and the $\det(\mathbf{J}) = 0$
boundary shifts in time; tracking either structure would require adapting
the control law to account for the flow unsteadiness. This extension is
left to future work.

Results are presented for a single field family (the double gyre). Whether
the separatrix tracker and OW boundary tracker generalize to other fields
with saddle-like or rotation-dominated regions is not addressed here,
though the estimation framework (Section~\ref{sec:estimation}) applies
to any smooth planar vector field.

All results are from simulation only. Hardware validation on the Decabot
testbed is not reported in this chapter; it is planned as a follow-on study.

The stability guarantees, and the reported success rates, are local.
Theorem~\ref{thm:separatrix} assumes initialization inside the tube
$T_w$, and all Monte Carlo trials start from a formation straddling the
separatrix. The attract mode converges to the critical set of $D$, which
includes the gyre-center maxima, so a team started deep inside a vortex
core parks at the core rather than on the separatrix: in supplementary
uniform-random-start trials over the domain, roughly half (47.9\%,
$n = 1000$) of starts reached the trench at zero noise. We do not
characterize the basin of attraction of the separatrix; prior
adaptive-navigation primitives [25], [4] likewise report local
convergence without basin estimates.

\section{Future Work}
\label{sec:sep_future}

Direct extensions of this work include:
\begin{itemize}
    \item \textbf{Time-varying double gyre.} The ocean HFR experiment of
          Section~\ref{sec:results_ocean} shows the controllers already
          operate on a genuinely time-varying field without modification;
          a time-derivative correction term to the navigation primitive,
          analyzed against the analytic $\varepsilon > 0$ double gyre where
          the exact structure motion is known, would let the stability
          results of Section~\ref{sec:stability} extend to that case with
          a quantified bound rather than empirical evidence alone.
    \item \textbf{Online formation reshaping.} Adapting the pair lengths
          $(L_k)$ and the SAS parameters to improve the conditioning of
          $\boldsymbol{\Phi}$ when the cluster is near a degenerate
          configuration.
    \item \textbf{Hardware deployment.} Validating both controllers on the
          Decabot testbed using the OptiTrack motion-capture system. The
          separatrix and Okubo-Weiss boundaries can be represented by a
          vector field emitted by an on-board computer and sampled by
          the robots' onboard sensors.
    \item \textbf{Extension to 3D.} Generalizing the quadratic estimator
          to three-dimensional fields, which would require at least 10 robots
          per component for full cubic determination, or a reduced-order
          estimator with different assumptions.
\end{itemize}

\section{Conclusion}
\label{sec:sep_conclusion}

This chapter presented a distributed framework for navigating a six-robot
formation along two Eulerian structural boundaries of the double gyre: the
separatrix and the Okubo-Weiss ($\det(\mathbf{J}) = 0$) boundary. A quadratic
estimator using simultaneous measurements from the six robots recovers the
Jacobian, its determinant, gradient, and Hessian at the cluster centroid with
no prior field knowledge, with closed-form noise gains verified empirically
to within $0.2\%$. Two adaptive controllers were derived, each with a
Lyapunov stability certificate; the separatrix result additionally extends
to a network-traversal theorem, since the trench a controller reaches
continues through every saddle it crosses rather than terminating there,
with single-saddle capture recovered as a tightened-threshold special case
demonstrated in simulation. Monte Carlo simulation over a two-dimensional
noise sweep placed the 50\% success level of both controllers at the same
noise level, $\sigma_{uv} \approx 0.01$, exactly the level at which a
single gradient estimate stops being informative: the estimator, not
either control law, sets the noise ceiling. On real, time-varying current
data from the Santa Barbara Channel the separatrix tracker followed the
dominant coherent-structure ridge for 28 hours of field time to a
consistent landfall, reproduced by 84 of 100 nearby starts, using only
the same local estimates validated in simulation, with no forecast, map,
or access to the field's future evolution.

This chapter's second-order estimation results rest on four fixed facts
that later work in this dissertation must not contradict: the
$\sigma_{uv} = 0.01$ noise level at which both the separatrix tracker
(43.3\%) and the Okubo-Weiss tracker (20.9\%) cross their respective
failure cliffs, the 47.9\% ($n = 1000$) basin-of-attraction figure for
uniform-random starts reaching the separatrix, the pentagon-plus-center
formation as the minimal shape one robot away from the degenerate-conic
failure of Proposition~\ref{prop:conic}, and the structural bias of
$\mathbf{H}_D$ that survives in its estimated eigenvectors (accurate to
fractions of a degree) but not in its estimated eigenvalues (indefinite
and order-of-magnitude too small near the wells). Read alongside
Chapter~\ref{ch:vector_field}, this chapter completes the dissertation's
two-step methodological arc: first-order Jacobian estimation from a
three-robot formation locates and characterizes isolated critical points,
while the second-order Hessian estimation developed here, from a
six-robot formation, extends that same distributed local-sampling
philosophy to continuous structures, trenches and boundaries, that a
single critical-point estimate cannot describe.

% ===================================================================
% CHAPTER 6: CONCLUSIONS
% ===================================================================
\chapter{Conclusions}
\label{ch:conclusions}

\section{Summary}

This dissertation developed a progression of tools for distributed
multirobot estimation and navigation in physical fields, moving from
the hardware substrate that makes such work testable, through the
geometric machinery every formation depends on, to two successive
orders of local field estimation.

Chapter~\ref{ch:testbed} established the Decabot hardware testbed: an
HSV-encoded printed vector field representation and a neural-network
sensor calibration procedure that removed the channel-interference
limitation of a prior testbed generation, achieving direction and
magnitude recovery accurate enough to validate closed-loop navigation
at full control rate. Every hardware result in the two chapters that
followed depends on this platform.

Chapter~\ref{ch:cluster_builder} generalized the cluster-space
kinematics underlying every formation used in this dissertation,
evolving the author's Master's thesis work on automated inverse-Jacobian
computation into a compositional grammar: a formation is a tree built
from a small library of closed-form charts, complete and square for any
tree by a direct dimension count, with its inverse Jacobian assembled in
a single depth-first pass rather than the $O(\text{depth}^2)$
recomputation the prior method required. A consensus orientation gauge
further removed a single-pointer-robot fragility that matters
specifically where the two following chapters drive a formation: the
low-signal neighborhood of a field critical point or boundary.

Chapter~\ref{ch:vector_field} showed that three robots, the SAS-3
triangle instantiation of Chapter~\ref{ch:cluster_builder}'s chart
library, are both necessary and sufficient to recover a vector field's
local Jacobian from simultaneous measurements alone, enabling critical
point localization, eigenvalue-based classification, and attraction and
orbital control laws, validated in simulation across eight canonical
fields and on hardware with a 100\% convergence success rate over 157
trials.

Chapter~\ref{ch:separatrix} extended this to second order: six robots,
the Pentagon-6 instantiation of the same chart library, recover a
flow's local Jacobian, Okubo-Weiss parameter, and Hessian, proven
minimal and almost-always sufficient, enabling two modified-Newton
controllers that track a separatrix trench and an Okubo-Weiss boundary
respectively, each with a Lyapunov stability certificate, validated on a
canonical double-gyre model and on 28 hours of real ocean current data.
This chapter is the direct answer to a question Chapter~\ref{ch:vector_field}
posed in its own Future Work (Section~\ref{sec:vf_future_work}) but
left open: that chapter's Jacobian-eigenvector framework, it suggested,
could in principle detect separatrices and heteroclinic orbits through
eigenvector alignment between neighboring critical points.
Chapter~\ref{ch:separatrix} is that suggestion delivered, generalized
past alignment between isolated points to a continuous second-order
estimate capable of tracking the boundary directly.

Appendix~\ref{app:four_robot} presented supplementary, unpublished
four-robot hardware trials generalizing Chapter~\ref{ch:vector_field}'s
three-robot result: simulation confirmed the theoretical noise-averaging
benefit of a four-robot overdetermined estimate, but the hardware
campaign did not realize it, a gap attributable to some combination of
motion-capture calibration bias, network load, and real-time execution
factors that this dissertation did not isolate.

\section{Future Work}

Several directions follow from the chapters above, some already stated
within their own Future Work sections and gathered here without
repetition.

\textbf{Time-varying fields.} Chapter~\ref{ch:vector_field}'s estimator
uses only instantaneous measurements and carries no state between
control cycles, so it does not assume field stationarity; validating
this on genuinely time-varying fields, beyond the reconstructed static
fields used there, is a natural next step.
Chapter~\ref{ch:separatrix}'s ocean HFR demonstration
(Section~\ref{sec:results_ocean}) already shows the separatrix tracker
operating on a real, time-varying field without modification, but a
time-derivative correction term, analyzed against the double gyre's
own $\varepsilon>0$ unsteady case where the exact structure motion is
known, would extend Chapter~\ref{ch:separatrix}'s stability results to
that case with a quantified bound rather than empirical evidence alone.

\textbf{Hardware validation of the second-order estimator.}
Chapter~\ref{ch:separatrix}'s results are simulation only; deploying
both the separatrix and Okubo-Weiss boundary controllers on the Decabot
testbed of Chapter~\ref{ch:testbed}, using a field emitted by an
onboard computer and sampled by the robots' own sensors, is the direct
extension.

\textbf{Extension to three dimensions.} Both estimation chapters are
restricted to planar fields. Chapter~\ref{ch:vector_field}'s first-order
estimator would need at least four non-coplanar robots in three
dimensions; Chapter~\ref{ch:separatrix}'s second-order estimator would
need at least ten robots per component for full cubic determination, or
a reduced-order estimator with different assumptions. Neither extension
is developed here.

\textbf{Online formation reshaping.} Chapter~\ref{ch:separatrix} notes
that adapting pair lengths and SAS parameters to improve the
conditioning of its formation matrix $\boldsymbol{\Phi}$, when the
cluster nears a degenerate configuration, remains future work; the same
idea applies to Chapter~\ref{ch:vector_field}'s triangle when it
approaches collinearity.

\textbf{Resolving the four-robot hardware gap.} Appendix~\ref{app:four_robot}'s
unexplained sim-to-hardware gap for the four-robot estimator would
benefit from a dedicated follow-up isolating calibration, network, and
real-time execution factors in turn, none of which the original
258-trial campaign was designed to separate.

\textbf{A typed field and controller interface.} Chapter~\ref{ch:cluster_builder}
notes that a typed interface between navigation primitives (such as
those of Chapters~\ref{ch:vector_field} and~\ref{ch:separatrix}) and the
kinematic generator would keep the reusable cluster-space core separate
from task-specific results; this remains unbuilt.

\section{Final Thoughts}

The four contributions of this dissertation share a single
methodological commitment: a multirobot team can recover meaningful
structure in an unknown field using nothing but its own simultaneous,
local measurements, no map, no forecast, and no prior model. What
changes from Chapter~\ref{ch:vector_field} to Chapter~\ref{ch:separatrix}
is only the order of the local model, linear to quadratic, and the
formation size that order requires, three robots to six, both
generalizations of the single kinematic substrate
Chapter~\ref{ch:cluster_builder} built from the author's earlier work.
Two robots short of what full generality would suggest are needed at
each order (three, not two; six, not five), and each of those extra
robots is doing exactly the work the minimality proofs
(Chapter~\ref{ch:vector_field}, Section~\ref{sec:vf_minimality};
Chapter~\ref{ch:separatrix}, Proposition~\ref{prop:minimality}) say it
must: no more, no less. The arc from a hardware testbed capable of
sensing a field at all, to a formation grammar capable of describing
any team that senses it, to two successive orders of what that team can
infer, is intended to read as one continuous argument rather than four
separate papers bound together, and the two open extensions this
chapter has just listed, time-varying fields and hardware validation of
the second-order estimator, are the two most direct ways to continue it.

% ===================================================================
% REFERENCES
% ===================================================================
\clearpage
\singlespacing
\addcontentsline{toc}{chapter}{References}

\begin{thebibliography}{99}

\bibitem{1} J. Cortés and M. Egerstedt, "Coordinated control of multi-robot systems: A survey," \textit{SICE Journal of Control, Measurement, and System Integration}, vol. 10, no. 6, pp. 495--503, 2017.

\bibitem{2} P. Ögren, E. Fiorelli, and N. E. Leonard, "Cooperative control of mobile sensor networks: Adaptive gradient climbing in a distributed environment," \textit{IEEE Transactions on Automatic Control}, vol. 49, no. 8, pp. 1292--1302, 2004.

\bibitem{3} L. Briñón-Arranz, A. Renzaglia, and L. Schenato, "Multirobot symmetric formations for gradient and hessian estimation with application to source seeking," \textit{IEEE Transactions on Robotics}, vol. 35, no. 3, pp. 782--789, 2019.

\bibitem{4} C. A. Kitts, R. T. McDonald, and M. A. Neumann, "Adaptive navigation control primitives for multirobot clusters: Extrema finding, contour following, ridge/trench following, and saddle point station keeping," \textit{IEEE Access}, vol. 6, pp. 17625--17642, 2018.

\bibitem{5} R. T. Mcdonald, C. A. Kitts, and M. A. Neumann, "Experimental implementation and verification of scalar field ridge, trench, and saddle point maneuvers using multirobot adaptive navigation," \textit{IEEE Access}, vol. 7, pp. 62950--62961, 2019.

\bibitem{6} R. K. Lee et al., "Multiple UAV adaptive navigation for three-dimensional scalar fields," \textit{IEEE Access}, vol. 9, pp. 122626--122654, 2021.

\bibitem{7} T. Adamek, C. A. Kitts, and I. Mas, "Gradient-based cluster space navigation for autonomous surface vessels," \textit{IEEE/ASME Transactions on Mechatronics}, vol. 20, no. 2, pp. 506--518, 2014.

\bibitem{8} S. McCammon et al., "Ocean front detection and tracking using a team of heterogeneous marine vehicles," \textit{Journal of Field Robotics}, vol. 38, no. 6, pp. 854--881, 2021.

\bibitem{9} S. R. Ramp et al., "Preparing to predict: the second autonomous ocean sampling network (AOSN-II) experiment in the Monterey Bay," \textit{Deep Sea Research Part II: Topical Studies in Oceanography}, vol. 56, no. 3-5, pp. 68--86, 2009.

\bibitem{10} J. Das et al., "Coordinated sampling of dynamic oceanographic features with underwater vehicles and drifters," \textit{The International Journal of Robotics Research}, vol. 31, no. 5, pp. 626--646, 2012.

\bibitem{11} J. Wang et al., "Dynamic plume tracking by cooperative robots," \textit{IEEE/ASME Transactions on Mechatronics}, vol. 24, no. 2, pp. 609--620, 2019.

\bibitem{12} D. Kularatne and M. A. Hsieh, "Tracking attracting manifolds in flows," \textit{Autonomous Robots}, vol. 41, no. 8, pp. 1575--1588, 2017.

\bibitem{13} M. Michini, H. Rastgoftar, M. A. Hsieh, and S. Jayasuriya, "Distributed formation control for collaborative tracking of manifolds in flows," in \textit{Proc. American Control Conference}, pp. 3874--3880, 2014.

\bibitem{14} X. Tricoche, Gerik Scheuermann, and H. Hagen, "Continuous topology simplification of planar vector fields," in \textit{Proc. IEEE Visualization}, pp. 159--166, 2001.

\bibitem{15} W. G. Pichel et al., "Marine debris collects within the North Pacific subtropical convergence zone," \textit{Marine Pollution Bulletin}, vol. 54, no. 8, pp. 1207--1211, 2007.

\bibitem{16} M. Mateus, R. Canelas, L. Pinto, and N. Vaz, "When tragedy strikes: potential contributions from ocean observation to search and rescue operations after drowning accidents," \textit{Frontiers in Marine Science}, vol. 7, p. 55, 2020.

\bibitem{17} E. Van Sebille et al., "The physical oceanography of the transport of floating marine debris," \textit{Environmental Research Letters}, vol. 15, no. 2, p. 023003, 2020.

\bibitem{18} K. Garg and D. Panagou, "Finite-time estimation and control for multi-aircraft systems under wind and dynamic obstacles," \textit{Journal of Guidance, Control, and Dynamics}, vol. 42, no. 7, pp. 1489--1505, 2019.

\bibitem{19} P. Keramea, K. Spanoudaki, G. Zodiatis, G. Gikas, and G. Sylaios, "Oil spill modeling: A critical review on current trends, perspectives, and challenges," \textit{Journal of Marine Science and Engineering}, vol. 9, no. 2, p. 181, 2021.

\bibitem{20} J. Wang, Y. Lin, R. Liu, and J. Fu, "Odor source localization of multi-robots with swarm intelligence algorithms: A review," \textit{Frontiers in Neurorobotics}, vol. 16, p. 949888, 2022.

\bibitem{21} H. Ji, M. Xu, W. Huang, and K. Yang, "The influence of oil leaking rate and ocean current velocity on the migration and diffusion of underwater oil spill," \textit{Scientific Reports}, vol. 10, no. 1, p. 9226, 2020.

\bibitem{22} C. Dong et al., "Oceanic mesoscale eddies," \textit{Ocean-Land-Atmosphere Research}, vol. 4, p. 0081, 2025.

\bibitem{23} Z. Li, X. Wang, J. Hu, F. P. Andutta, and Z. Liu, "A study on an anticyclonic-cyclonic eddy pair off Fraser Island, Australia," \textit{Frontiers in Marine Science}, vol. 7, p. 594358, 2020.

\bibitem{24} R. Deogharia, H. Gupta, S. Sil, A. Gangopadhyay, and A. Shee, "On the evidence of helico-spiralling recirculation within coherent cores of eddies using Lagrangian approach," \textit{Scientific Reports}, vol. 14, no. 1, p. 11014, 2024.

\bibitem{25} R. T. McDonald, C. A. Kitts, and M. A. Neumann, "Navigation of scalar fronts with multirobot clusters in simulation and experiment," \textit{IEEE Systems Journal}, vol. 14, no. 3, pp. 3755--3766, 2020.

\bibitem{26} R. K. Lee, C. A. Kitts, M. A. Neumann, and R. T. McDonald, "3-d adaptive navigation: Multirobot formation control for seeking and tracking of a moving source," in \textit{Proc. IEEE Aerospace Conference}, pp. 1--12, 2020.

\bibitem{27} F. Pagano et al., "Distributed Multi-Robot Active-Sensing of a Diffusive Source," \textit{IEEE Robotics and Automation Letters}, 2025.

\bibitem{28} Z. Jin, H. Li, Z. Qin, and Z. Wang, "Gradient-free cooperative source-seeking of quadrotor under disturbances and communication constraints," \textit{IEEE Transactions on Industrial Electronics}, vol. 72, no. 2, pp. 1969--1979, 2024.

\bibitem{29} S. Li, R. Kong, and Y. Guo, "Cooperative distributed source seeking by multiple robots: Algorithms and experiments," \textit{IEEE/ASME Transactions on Mechatronics}, vol. 19, no. 6, pp. 1810--1820, 2014.

\bibitem{30} Z. Deng and J. Luo, "Fully distributed algorithms for constrained nonsmooth optimization problems of general linear multiagent systems and their application," \textit{IEEE Transactions on Automatic Control}, vol. 69, no. 2, pp. 1377--1384, 2023.

\bibitem{31} W. Li, J. A. Farrell, S. Pang, and R. M. Arrieta, "Moth-inspired chemical plume tracing on an autonomous underwater vehicle," \textit{IEEE Transactions on Robotics}, vol. 22, no. 2, pp. 292--307, 2006.

\bibitem{32} Y. Deng, B. Zhu, and H. Duan, "Bioinspired bearing–based target enclosing control for unmanned aerial vehicle swarm," \textit{IEEE/ASME Transactions on Mechatronics}, vol. 30, no. 4, pp. 2699--2709, 2024.

\bibitem{33} C. Waight and C. Kitts, "A functional indoor testbed for multirobot adaptive navigation in vector field environments," in \textit{Proc. ASME Int. Design Engineering Technical Conf. Computers and Information in Engineering Conf.}, vol. 89251, 2025.

\bibitem{34} M. Michini, M. A. Hsieh, E. Forgoston, and I. B. Schwartz, "Robotic tracking of coherent structures in flows," \textit{IEEE Transactions on Robotics}, vol. 30, no. 3, pp. 593--603, 2014.

\bibitem{35} D. R. Nelson, D. B. Barber, T. W. McLain, and R. W. Beard, "Vector field path following for miniature air vehicles," \textit{IEEE Transactions on Robotics}, vol. 23, no. 3, pp. 519--529, 2007.

\bibitem{36} W. Yao, H. G. de Marina, B. Lin, and M. Cao, "Singularity-free guiding vector field for robot navigation," \textit{IEEE Transactions on Robotics}, vol. 37, no. 4, pp. 1206--1221, 2021.

\bibitem{37} D. Panagou, "Motion planning and collision avoidance using navigation vector fields," in \textit{Proc. IEEE International Conference on Robotics and Automation}, 2014, pp. 2513--2518.

\bibitem{38} T. Liddy, T.-F. Lu, P. Lozo, and D. Harvey, "Obstacle avoidance using complex vector fields," in \textit{Proc. Australasian Conference on Robotics and Automation}, 2008.

\bibitem{39} X. Wang et al., "Adaptive vector field guidance without a priori knowledge of course dynamics and wind," \textit{IEEE/ASME Transactions on Mechatronics}, vol. 27, no. 6, pp. 4597--4607, 2022.

\bibitem{40} T. Inanc, S. C. Shadden, and J. E. Marsden, "Optimal trajectory generation in ocean flows," in \textit{Proc. American Control Conference}, pp. 674--679, 2005.

\bibitem{41} R. N. Smith, M. Schwager, S. L. Smith, B. H. Jones, D. Rus, and G. S. Sukhatme, "Persistent ocean monitoring with underwater gliders: Adapting sampling resolution," \textit{Journal of Field Robotics}, vol. 28, no. 5, pp. 714--741, 2011.

\bibitem{42} G. Knizhnik, P. Li, X. Yu, and M. A. Hsieh, "Flow-based control of marine robots in gyre-like environments," in \textit{Proc. IEEE International Conference on Robotics and Automation (ICRA)}, pp. 3047--3053, 2022.

\bibitem{43} M. Michini, M. A. Hsieh, E. Forgoston, and I. B. Schwartz, "Experimental validation of robotic manifold tracking in gyre-like flows," in \textit{Proc. IEEE/RSJ International Conference on Intelligent Robots and Systems (IROS)}, pp. 2306--2311, 2014.

\bibitem{44} Y. Jiao, H. Hang, J. Merel, and E. Kanso, "Sensing flow gradients is necessary for learning autonomous underwater navigation," \textit{Nature Communications}, vol. 16, no. 1, p. 3044, 2025.

\bibitem{45} Z. Zhou, J. Liu, and J. Yu, "A survey of underwater multi-robot systems," \textit{IEEE/CAA Journal of Automatica Sinica}, vol. 9, no. 1, pp. 1--18, 2021.

\bibitem{46} C. A. Kitts and I. Mas, "Cluster space specification and control of mobile multirobot systems," \textit{IEEE/ASME Transactions on Mechatronics}, vol. 14, no. 2, pp. 207--218, 2009.

\bibitem{47} I. Mas and C. Kitts, "Obstacle avoidance policies for cluster space control of nonholonomic multirobot systems," \textit{IEEE/ASME Transactions on Mechatronics}, vol. 17, no. 6, pp. 1068--1079, 2012.

\bibitem{48} S. Tomer et al., "A low-cost indoor testbed for multirobot adaptive navigation research," in \textit{Proc. IEEE Aerospace Conference}, pp. 1--12, 2018.

\bibitem{49} S. T. Hart and C. A. Kitts, "Unifying control architecture for reactive particle swarms," \textit{IEEE/ASME Transactions on Mechatronics}, vol. 28, no. 2, pp. 873--883, 2022.

\bibitem{50} R. Cooper et al., "Initial Study of Multirobot Adaptive Navigation for Exploring Environmental Vector Fields," in \textit{Proc. IEEE Aerospace Conference}, 2019.

\bibitem{51} D. Rajabhandharaks, R. T. McDonald, M. A. Neumann, and C. A. Kitts, "Indoor testbed for vector field multirobot adaptive navigation: Feasibility study," in \textit{Proc. ASME Int. Design Engineering Technical Conf. Computers and Information in Engineering Conf.}, vol. 59292, p. V009T12A043, 2019.

\bibitem{52} G. Haller, "Lagrangian coherent structures," \textit{Annual
Review of Fluid Mechanics}, vol. 47, pp. 137--162, 2015.

\bibitem{53} Y. Sun, G. Ma, M. Liu, C. Li, and J. Liang, "Distributed
finite-time coordinated control for multi-robot systems with disturbances
and directed topologies," \textit{Transactions of the Institute of
Measurement and Control}, vol. 39, no. 12, pp. 2907--2919, 2017.

\bibitem{54} T. D. Harms, S. T. M. Dawson, and B. J. McKeon, "Lagrangian
gradient regression for the detection of coherent structures from sparse
trajectory data," \textit{Royal Society Open Science}, vol. 11, no. 9,
p. 240586, 2024.

\bibitem{55} G. Haller and G. Yuan, "Lagrangian coherent structures and mixing
in two-dimensional turbulence," \textit{Physica D: Nonlinear Phenomena},
vol. 147, no. 3-4, pp. 352--370, 2000.

\bibitem{56} T. Peacock and G. Haller, "Lagrangian coherent structures: The
hidden skeleton of fluid flows," \textit{Physics Today}, vol. 66, no. 2,
pp. 41--47, 2013.

\bibitem{57} A. Okubo, "Horizontal dispersion of floatable particles in the
vicinity of velocity singularities such as convergences," \textit{Deep-Sea
Research}, vol. 17, no. 3, pp. 445--454, 1970.

\bibitem{58} J. Weiss, "The dynamics of enstrophy transfer in two-dimensional
hydrodynamics," \textit{Physica D: Nonlinear Phenomena}, vol. 48, no. 2-3,
pp. 273--294, 1991.

\bibitem{59} B.-L. Hua and P. Klein, "An exact criterion for the stirring
properties of nearly two-dimensional turbulence," \textit{Physica D}, vol. 113,
no. 1, pp. 98--110, 1998.

\bibitem{60} M. Serra and G. Haller, "Objective Eulerian coherent
structures," \textit{Chaos: An Interdisciplinary Journal of Nonlinear
Science}, vol. 26, no. 5, p. 053110, 2016.

\bibitem{61} Y. N. Dauphin, R. Pascanu, C. Gulcehre, K. Cho, S. Ganguli,
and Y. Bengio, "Identifying and attacking the saddle point problem in
high-dimensional non-convex optimization," in \textit{Advances in Neural
Information Processing Systems}, vol. 27, pp. 2933--2941, 2014.

\bibitem{62} P. J. Nolan, M. Serra, and S. D. Ross, "Finite-time Lyapunov
exponents in the instantaneous limit and material transport,"
\textit{Nonlinear Dynamics}, vol. 100, no. 4, pp. 3825--3852, 2020.

\bibitem{63} J. C. R. Hunt, A. A. Wray, and P. Moin, "Eddies, streams, and
convergence zones in turbulent flows," in \textit{Proc. Summer Program,
Center for Turbulence Research}, Rep. CTR-S88, pp. 193--208, 1988.

\bibitem{64} S. C. Shadden, F. Lekien, and J. E. Marsden, "Definition and
properties of Lagrangian coherent structures from finite-time Lyapunov
exponents in two-dimensional aperiodic flows," \textit{Physica D}, vol. 212,
no. 3-4, pp. 271--304, 2005.

\bibitem{65} G. Haller, "Dynamic rotation and stretch tensors from a
dynamic polar decomposition," \textit{Journal of the Mechanics and Physics
of Solids}, vol. 86, pp. 70--93, 2016.

\bibitem{66} C. J. Waight, ``An Algorithm for Calculating the Inverse Jacobian
of Multirobot Systems in a Cluster Space Formulation,'' Master's thesis,
Santa Clara University, Santa Clara, CA, USA.
% %% TODO: SYNC WITH LATEST DRAFT -- confirm exact submission year with the
% user before final numbering; this is the Master's thesis Chapter 3 evolves
% (see Thesis_Execution_Plan.md addendum). It is very likely the same
% antecedent the source HTML/IDETC-II draft calls "the 2016 SCU toolbox
% thesis" ([FILL] bibitem 3 in cluster_kinematics_idetc.tex), but the exact
% year should be confirmed rather than assumed to be 2016.

% %% TODO: SYNC WITH LATEST DRAFT -- Chapter 3's own source material
% (cluster_kinematics_idetc.tex) has four [FILL] bibliography entries not
% yet resolved: a Mas and Kitts follow-on paper (cluster-space dynamic/
% partitioned control), the 2016 SCU toolbox thesis (very likely entry 66
% above -- confirm with the user rather than merging silently), a virtual-
% structure/formation-control citation, and a self-citation of the testbed
% paper (already covered by entry 33 above, so that one specifically can be
% dropped once confirmed). These are not assigned numbers here per
% Thesis_Execution_Plan.md Section 7, item 4: they cannot be merged into the
% consolidated list while their titles/venues/years remain unknown.

\end{thebibliography}

% ===================================================================
% APPENDICES
% ===================================================================
\appendix

\chapter{Vector Field Environments}
\label{app:vector_fields}

Six vector field environments centered at the origin were used to evaluate
the multirobot navigation algorithms of Chapter~\ref{ch:vector_field}. Let
$\rho=\sqrt{x^2+y^2}$ and $\theta=\operatorname{atan2}(y,x)$ (this polar
radius $\rho$ is renamed from the source material's bare $r$ to avoid
collision with the orbital radial magnitude of
Chapter~\ref{ch:vector_field}, Section~\ref{sec:control_primitives}).

\paragraph{Vortex.} Pure rotational flow:
\begin{equation}
\mathbf{v}_{\text{vortex}}(x,y) = \begin{bmatrix} -\rho\sin\theta \\ \rho\cos\theta \end{bmatrix}.
\end{equation}

\paragraph{Sinking Vortex.} Spiraling inward flow:
\begin{equation}
\mathbf{v}_{\text{sink-vortex}}(x,y) = 0.4\begin{bmatrix} -\rho\sin\theta \\ \rho\cos\theta \end{bmatrix} + 0.15\begin{bmatrix} -x \\ -y \end{bmatrix}.
\end{equation}

\paragraph{Sink.} Radial inward flow:
\begin{equation}
\mathbf{v}_{\text{sink}}(x,y) = \begin{bmatrix} -x \\ -y \end{bmatrix}.
\end{equation}

\paragraph{Source.} Radial outward flow:
\begin{equation}
\mathbf{v}_{\text{source}}(x,y) = \begin{bmatrix} x \\ y \end{bmatrix}.
\end{equation}

\paragraph{Spewing Vortex.} Spiraling outward flow:
\begin{equation}
\mathbf{v}_{\text{spew-vortex}}(x,y) = -0.4\begin{bmatrix} -\rho\sin\theta \\ \rho\cos\theta \end{bmatrix} - 0.15\begin{bmatrix} -x \\ -y \end{bmatrix}.
\end{equation}

\paragraph{Saddle.} Hyperbolic saddle point rotated $45°$; coordinates are
first rotated by $\pi/4$, the saddle field is applied, then rotated back by
$-\pi/4$:
\begin{equation}
\mathbf{v}_{\text{saddle}}(x,y) = \begin{bmatrix} -y \\ -x \end{bmatrix}.
\end{equation}

\chapter{Double-Gyre Analytic Forms}
\label{app:double_gyre}

For the steady double gyre with amplitude $A$, $x_f=x+1$, $y_f=y+0.5$,
define $s_x=\sin(\pi x_f)$, $c_x=\cos(\pi x_f)$, $s_y=\sin(\pi y_f)$,
$c_y=\cos(\pi y_f)$. The field is $u=-\pi A s_xc_y$, $v=\pi Ac_xs_y$.

The Jacobian entries are
\begin{equation}
\frac{\partial u}{\partial x} = -\pi^2Ac_xc_y, \qquad \frac{\partial u}{\partial y} = \pi^2As_xs_y,
\end{equation}
\begin{equation}
\frac{\partial v}{\partial x} = -\pi^2As_xs_y, \qquad \frac{\partial v}{\partial y} = \pi^2Ac_xc_y,
\end{equation}
so the flow is divergence free and
\begin{equation}
D = \det(\mathbf{J}) = -\pi^4A^2\bigl(c_x^2c_y^2 - s_x^2s_y^2\bigr) = -\frac{\pi^4A^2}{2}\bigl(\cos2\pi x_f + \cos2\pi y_f\bigr),
\label{eq:det_analytic}
\end{equation}
using $c^2=(1+\cos2\alpha)/2$ and $s^2=(1-\cos2\alpha)/2$. $D$ is positive
at the gyre centers ($D=\pi^4A^2$), negative at the corner saddles
($D=-\pi^4A^2$), and zero on the diamonds
$x_f\pm y_f\in\tfrac12+\mathbb{Z}$, by the factorization
$\cos2\pi x_f+\cos2\pi y_f = 2\cos\pi(x_f+y_f)\cos\pi(x_f-y_f)$.

Differentiating Equation~\eqref{eq:det_analytic},
\begin{equation}
\nabla D = \pi^5A^2\begin{bmatrix}\sin2\pi x_f\\\sin2\pi y_f\end{bmatrix},
\qquad
\mathbf{H}_D = 2\pi^6A^2\begin{bmatrix}\cos2\pi x_f & 0\\0 & \cos2\pi y_f\end{bmatrix}.
\label{eq:gradhess_analytic}
\end{equation}
The Hessian is diagonal everywhere, so $D$ is additively separable in
$x_f$ and $y_f$, which is what makes the trench normal form of
Assumption~\ref{ass:trench} exact for this field. The vorticity is
$\omega=-2\pi^2As_xs_y$, which vanishes on the separatrix $x_f=1$; there
$\mathbf{J}$ is symmetric and $\sqrt{-D}$ is the exact instantaneous
strain rate. The Okubo-Weiss parameter is $Q=-D$ by incompressibility, so
the $Q=0$ and $D=0$ contours coincide.

\chapter{Cluster Kinematics for the Formations Used in This Dissertation}
\label{app:cluster_kinematics}

This appendix collects the specific forward kinematics, inverse kinematics,
and inverse Jacobian for the two formations used in this dissertation, the
SAS-3 triangle of Chapter~\ref{ch:vector_field} and the pentagon-plus-center
of Chapter~\ref{ch:separatrix}, as concrete instances of the general chart
library and single-pass assembly algorithm of Chapter~\ref{ch:cluster_builder}.
Both were originally derived independently, in each chapter's own source
material, using notation that collides with the rest of this dissertation in
three places (Chapter~\ref{ch:intro}, Section~\ref{sec:notation}); the
renamings below apply the fixes stated there.

\section{Three-Robot SAS Triangle}
\label{app:sas_kinematics}

This is the SAS-3 chart of Chapter~\ref{ch:cluster_builder},
Section~\ref{sec:atom_sas}, restated here with the full worked forward and
inverse kinematics used by Chapter~\ref{ch:vector_field}.

\subsection{Forward Kinematics}

The forward kinematics transform robot positions $(x_1,y_1)$, $(x_2,y_2)$,
$(x_3,y_3)$ to SAS parameters and centroid coordinates. The centroid
position is
\begin{equation}
x_c = \frac{x_1+x_2+x_3}{3}, \qquad y_c = \frac{y_1+y_2+y_3}{3}.
\end{equation}
The side lengths are
\begin{equation}
p = \sqrt{(x_2-x_1)^2+(y_2-y_1)^2}, \qquad
q = \sqrt{(x_3-x_2)^2+(y_3-y_2)^2}, \qquad
s = \sqrt{(x_1-x_3)^2+(y_1-y_3)^2},
\end{equation}
where $s$ (renamed from the source material's bare $r$, per
Chapter~\ref{ch:intro}, Section~\ref{sec:notation}'s $r$-collision fix) is
the third side, the 1--3 diagonal, an intermediate quantity used only in the
law-of-cosines angle below. The interior angle at vertex 2 is
\begin{equation}
\beta = \arccos\left(\frac{p^2+q^2-s^2}{2pq}\right).
\end{equation}
The remaining two vertex angles, needed for neither this appendix's
kinematics nor any equation elsewhere in this dissertation, are
\begin{equation}
\phi_1 = \arccos\left(\frac{p^2+s^2-q^2}{2ps}\right), \qquad
\phi_3 = \arccos\left(\frac{q^2+s^2-p^2}{2qs}\right),
\end{equation}
renamed from the source material's $\alpha,\gamma$ to avoid the four-way
$\alpha$ collision of Chapter~\ref{ch:intro}, Section~\ref{sec:notation}.
The orientation angle is
\begin{equation}
\theta_c = \operatorname{atan2}(y_1-y_2,\,x_1-x_2),
\end{equation}
the direction from robot 2 to robot 1 (corrected from the source paper's
robot-1-to-robot-2 convention, which disagreed by $\pi$ with the inverse
kinematics below; the two conventions are mathematically equivalent up to a
constant reference shift, and the correction only affects the unused,
never-actively-commanded orientation variable, not any reported result).

\subsection{Inverse Kinematics}

Given the SAS parameters $(p,\beta,q)$ and centroid parameters
$(x_c,y_c,\theta_c)$, robot positions are recovered by constructing the
triangle in a local frame, with robot 2 at the origin, then applying
rotation and translation:
\begin{equation}
x_2^{\text{local}}=0,\ y_2^{\text{local}}=0, \qquad
x_1^{\text{local}}=p,\ y_1^{\text{local}}=0, \qquad
x_3^{\text{local}}=q\cos\beta,\ y_3^{\text{local}}=q\sin\beta.
\end{equation}
The local centroid is
\begin{equation}
x_c^{\text{local}} = \frac{x_1^{\text{local}}+x_2^{\text{local}}+x_3^{\text{local}}}{3}, \qquad
y_c^{\text{local}} = \frac{y_1^{\text{local}}+y_2^{\text{local}}+y_3^{\text{local}}}{3}.
\end{equation}
After centering the triangle at the local origin and applying rotation
$\theta_c$ and translation to $(x_c,y_c)$,
\begin{equation}
\begin{bmatrix} x_i \\ y_i \end{bmatrix} =
\begin{bmatrix} \cos\theta_c & -\sin\theta_c \\ \sin\theta_c & \cos\theta_c \end{bmatrix}
\begin{bmatrix} x_i^{\text{local}}-x_c^{\text{local}} \\ y_i^{\text{local}}-y_c^{\text{local}} \end{bmatrix}
+ \begin{bmatrix} x_c \\ y_c \end{bmatrix}
\end{equation}
for $i\in\{1,2,3\}$.

\subsection{Inverse Jacobian}

The inverse Jacobian maps shape-parameter velocities to robot velocities,
\begin{equation}
\begin{bmatrix} \delta x_1 \\ \delta y_1 \\ \delta x_2 \\ \delta y_2 \\ \delta x_3 \\ \delta y_3 \end{bmatrix}
= \mathbf{J}_{c,\text{SAS}}^{-1}
\begin{bmatrix} \delta x_c \\ \delta y_c \\ \delta\theta_c \\ \delta p \\ \delta\beta \\ \delta q \end{bmatrix},
\end{equation}
where $\mathbf{J}_{c,\text{SAS}}^{-1}$ (renamed from the source material's
bare $\mathbf{J}^{-1}$ to avoid collision with the field Jacobian
$\mathbf{J}$ of Chapters~\ref{ch:vector_field} and~\ref{ch:separatrix}) is a
$6\times6$ matrix of partial derivatives
$\partial x_i/\partial\xi_j$, $\partial y_i/\partial\xi_j$ for
$\xi_j\in\{x_c,y_c,\theta_c,p,\beta,q\}$, computed by differentiating the
inverse kinematics above with respect to each shape parameter. This is
precisely the size-3 leaf block $S_{\text{leaf},3}$ of
Chapter~\ref{ch:cluster_builder}, Section~\ref{sec:assembly}: its first
three columns are the sensitivity matrix used when this triangle appears as
a leaf inside a larger composition tree, and its remaining three columns
fill the shape-variable columns of the global inverse Jacobian directly.
% %% TODO: SYNC WITH LATEST DRAFT -- the source paper omits the explicit
% 6x6 entries "for brevity"; per Thesis_Execution_Plan.md Section 6, item 4
% and Section 10, question 5, these could now be derived mechanically from
% Chapter 3's general assembly algorithm rather than left omitted. Not done
% here; flagged for the user to decide.

\section{Six-Robot Pentagon Plus Center}
\label{app:pentagon_kinematics}

This is the Pentagon-6 chart of Chapter~\ref{ch:cluster_builder},
Section~\ref{sec:atom_pentagon}, under the consensus gauge of
Section~\ref{sec:gauge}, used by Chapter~\ref{ch:separatrix}. Two
concrete instantiations of this atom appear in this dissertation. The
minimal-symmetry version, five robots on a ring plus one at the center,
is given in full in Section~\ref{sec:atom_pentagon} and is not repeated
here. Chapter~\ref{ch:separatrix}'s own formation-control layer
(distinct from its estimator, which reads only the six measured
positions, per Section~\ref{sec:formation}) instead groups the six
robots as three rigid pairs and parameterizes their midpoint triangle,
described next; this is the pentagon-\emph{of-pairs} configuration
whose ring radius $\rho=0.075$ generates the pentagon-plus-center
formation only at the specific setpoint given in
Chapter~\ref{ch:separatrix}, Section~\ref{sec:formation} and its
simulation configuration (Chapter~\ref{ch:separatrix},
Section~\ref{sec:methods}).

\subsection{Forward Kinematics (Pentagon-of-Pairs)}

The pentagon-of-pairs cluster has six robots indexed $i=1,\ldots,6$,
arranged as three pairs. Pair $k$ ($k=1,2,3$) consists of robots
$2k-1$ and $2k$. The midpoint of pair $k$ is
$\mathbf{m}_k=(\mathbf{p}_{2k-1}+\mathbf{p}_{2k})/2$. The three
midpoints form a triangle; the centroid of that triangle is
$\mathbf{p}_c=(\mathbf{m}_1+\mathbf{m}_2+\mathbf{m}_3)/3$.

The SAS parameters of the midpoint triangle are defined as in the
three-robot case of Section~\ref{sec:atom_sas}: $p_1=\lVert\mathbf{m}_2-\mathbf{m}_1\rVert$,
$q_1=\lVert\mathbf{m}_3-\mathbf{m}_2\rVert$, and $\beta_1$ is the
interior angle at $\mathbf{m}_2$. The cluster heading is
$\theta_c=\operatorname{atan2}(m_{2,y}-m_{1,y},m_{2,x}-m_{1,x})$. Each
pair is described by its half-length $L_k/2$ and orientation angle
$\psi_k$, so the two robots in pair $k$ are at
$\mathbf{m}_k\pm(L_k/2)(\cos\psi_k,\sin\psi_k)^{\mathsf T}$.

The full twelve-state cluster configuration is
$(x_c,y_c,\theta_c,p_1,\beta_1,q_1,L_2,\psi_2,L_3,\psi_3,L_4,\psi_4)$,
where pair 1 uses $L_1$ fixed and $\psi_1=\theta_c$ by convention.

\subsection{Inverse Kinematics (Pentagon-of-Pairs)}

Given the twelve-state cluster configuration, the six robot positions
are recovered by: constructing the midpoint triangle from
$(p_1,\beta_1,q_1)$ in a local frame, following the same procedure as
the SAS-3 triangle of Section~\ref{sec:atom_sas}; rotating and
translating to the global frame using $\theta_c$, $x_c$, $y_c$; and
placing each pair of robots by offsetting from the midpoint along
$(\cos\psi_k,\sin\psi_k)$ by $L_k/2$ in each direction.

\subsection{Inverse Kinematic Jacobian (Pentagon-of-Pairs)}

The $12\times12$ inverse kinematic Jacobian maps desired cluster state
velocities $\dot{\boldsymbol{\xi}}$ to robot velocity commands
$(\dot x_1,\dot y_1,\ldots,\dot x_6,\dot y_6)^{\mathsf T}$:
\begin{equation}
\begin{bmatrix} \dot x_1\\\dot y_1\\\vdots\\\dot x_6\\\dot y_6 \end{bmatrix}
= \mathbf{J}_{c,\text{pentagon}}^{-1}
\begin{bmatrix} \dot x_c\\\dot y_c\\\dot\theta_c\\\dot p_1\\\dot\beta_1\\\dot q_1\\\dot L_2\\\dot\psi_2\\\dot L_3\\\dot\psi_3\\\dot L_4\\\dot\psi_4 \end{bmatrix},
\label{eq:inv_kin_jac_pentagon}
\end{equation}
where $\mathbf{J}_{c,\text{pentagon}}^{-1}$ (renamed from the source
material's bare $\mathbf{J}_c^{-1}$ for consistency with the
Chapter~\ref{ch:cluster_builder} notation convention) is a
$12\times12$ matrix of partial derivatives
$\partial(x_i,y_i)/\partial\xi_j$ computed analytically from the
inverse kinematics above. This pentagon-of-pairs parameterization is
precisely an instance of Chapter~\ref{ch:cluster_builder}'s
cluster-of-clusters construction (Section~\ref{sec:grammar}): a
three-robot SAS meta-node over three two-robot pair sub-clusters,
so Equation~\eqref{eq:inv_kin_jac_pentagon}'s $12\times12$ block is
exactly what Chapter~\ref{ch:cluster_builder}, Algorithm~\ref{alg:assemble}
assembles in a single depth-first pass from the SAS leaf rule
$S_{\text{leaf},3}$ at the root and the pair leaf rule
$S_{\text{leaf},2}$ (Equation in Section~\ref{sec:assembly}) at each of
the three children.
% %% TODO: SYNC WITH LATEST DRAFT -- the source separatrix paper states
% "the explicit entries are omitted here for brevity" for this 12x12
% Jacobian. Per Thesis_Execution_Plan.md Section 6, item 9 and Section 10,
% question 5, this could now be derived as a direct output of Chapter 3's
% general assembly algorithm rather than left omitted. Not done here;
% flagged for the user to decide.

\chapter{Four-Robot Generalization: Simulation and Hardware Trials (Unpublished)}
\label{app:four_robot}

This appendix generalizes Chapter~\ref{ch:vector_field}'s three-robot
critical-point estimator to a four-robot, overdetermined formation. It
predates the final, three-robot-only submitted version of that chapter
and was not peer-reviewed; it is included here as supplementary material
because the underlying hardware campaign was real and substantial (258
trials), not because its results were judged strong enough for
publication. Chapter~\ref{ch:vector_field}'s own minimality result
(Section~\ref{sec:vf_minimality}) proves three robots are
\emph{necessary}; it says nothing about whether a fourth robot can help,
and this appendix is the empirical answer to that separate question.

\section{Motivation}

Three robots are the minimum formation size for the linear (first-order)
field estimator of Chapter~\ref{ch:vector_field}; a fourth robot makes
the estimation problem overdetermined, which in principle allows
noise-averaging benefits unavailable to an exactly determined
three-robot fit. This appendix presents three estimator variants
explored for the four-robot case, a real 258-trial hardware campaign
(236 convergence trials plus 22 orbital trials), and an honest
discussion of why the hardware results did not realize the
noise-averaging benefit that simulation predicted.

\section{Three Estimator Variants}

All three variants use a four-robot square formation as their kinematic
layer, an instance of Chapter~\ref{ch:cluster_builder}'s Cross-Diagonal-4
chart (Section~\ref{sec:atom_xdiag}): the diagonal-based
$(d_1,d_2,r_1,r_2,\psi)$ parameterization used here for the Four-Robot
Cluster Space Controller below is exactly that chart's state vector,
giving this appendix a second purpose as a worked example of
Chapter~\ref{ch:cluster_builder}'s general grammar at a formation size
between the SAS-3 triangle and the Pentagon-6.

\subsection{Planar Center}

A single overdetermined least-squares plane fit through all four
robots' measurements, the direct four-robot generalization of
Chapter~\ref{ch:vector_field}, Section~\ref{sec:vf_estimation}'s
three-robot linear system: with four robots and three unknowns per
component, the system $\mathbf{A}\boldsymbol{\theta}_u=\mathbf{u}$,
$\mathbf{A}\boldsymbol{\theta}_v=\mathbf{v}$ is solved via the
Moore-Penrose pseudoinverse rather than direct inversion. This is the
variant carried through to the final 258-trial hardware campaign
described below (implementation:
\textit{estimate\_center\_planar}/\textit{four\_planar\_center\_finder}
in \textit{src/control/quad\_primitives.py}).

\subsection{Dual Jacobian}

Splits the four-robot square into two overlapping triangles, robots
$\{1,2,3\}$ and $\{2,3,4\}$, computes an independent three-robot
Jacobian estimate (Chapter~\ref{ch:vector_field},
Equation~\eqref{eq:vf_pstar}) from each triangle, and averages the two
resulting critical-point estimates (implementation:
\textit{calculate\_jacobian\_from\_triangle} /
\textit{calculate\_jacobian\_top} / \textit{calculate\_jacobian\_bottom} /
\textit{dual\_jacobian\_center\_finder} in
\textit{src/control/quad\_primitives.py}). This variant appeared only
in early paper drafts (Older Paper Drafts 4-6) and was dropped by the
draft that fixed the final 258-trial campaign; it was claimed superior
at tight orbital radii in simulation (0.111~m error at the smallest
commanded radius, a 58\% improvement over the Planar Center variant at
that radius), but the reason for its removal from the final hardware
campaign is not recorded in the surviving drafts, and is not asserted
here.

\subsection{Four-Robot Cluster Space Controller}

The kinematic layer underneath both estimators above: a diagonal-based
quadrilateral parameterization with eight shape parameters
$(d_1,d_2,r_1,r_2,\psi,x_c,y_c,\theta_c)$, an instance of
Chapter~\ref{ch:cluster_builder}'s Cross-Diagonal-4 chart
(Section~\ref{sec:atom_xdiag}) with forward and inverse kinematics
given there in full. This controller subsection appeared in every
draft of Chapter~\ref{ch:vector_field}'s source paper prior to the
final submission-track version, including a dedicated appendix giving
the forward centroid equations and the four-robot analog of the
inverse-Jacobian error-propagation derivation, but was cut entirely
from the final three-robot-only paper.

\section{Hardware Campaign}

The campaign ran 80 three-robot-triangle trials per field (matching
Chapter~\ref{ch:vector_field}, Section~\ref{sec:vf_hardware}) plus an
additional 80 four-robot-square vortex trials using the Planar Center
estimator, for a planned 258 hardware trials total. As with the
three-robot campaign, some trials were lost to file-management errors
(3 saddle plus 1 vortex trials from the three-robot side), giving 236
usable convergence trials; combined with 22 orbital control trials at
varying commanded radii, this is the full 258-trial figure.

The four-robot vortex estimation RMSE was 34.67~mm, compared to
approximately 25~mm for the three-robot hardware baseline on comparable
fields (Chapter~\ref{ch:vector_field}, Table~\ref{tab:estimation_results}).
This is confirmed directly from the surviving data file
\textit{trunk/robots\_4/4robot\_vortex\_stats.csv}: 80 runs, mean bias
16.14~mm, precision (standard deviation) 30.69~mm, combining to the
34.67~mm RMSE figure. By contrast, four-robot \emph{simulated} RMSE
(13.44--15.48~mm across field types) was at or below three-robot
simulated RMSE (12.65--18.45~mm): simulation predicted the
noise-averaging benefit of an overdetermined system, and the hardware
campaign did not realize it.

\section{Discussion: Why Hardware Underperformed Simulation}

The gap between the four-robot simulation prediction and the four-robot
hardware result is real and worth stating plainly rather than
explaining away. The surviving paper drafts attributed it to
motion-capture calibration bias, specifically a systematic offset of
approximately 10~mm and a directional bias reaching $-0.086$~m at a
0.400~m commanded orbital radius. That explanation is plausible but not
the only one consistent with the data, and this appendix does not
assert it more strongly than the evidence supports. A fourth robot on
the same 10~Hz WiFi control loop increases network traffic and could
introduce additional latency or packet loss not present in the
three-robot campaign; the four-robot Simulink model
(\textit{trunk/robots\_4/chris\_4R\_all\_primitives.slx}) is
substantially larger than its three-robot counterpart
(\textit{trunk/robots\_3/three\_robot\_cluster\_all\_primitives.slx}),
and a larger real-time model can itself introduce timing jitter. Either
of these, alongside or instead of calibration bias, could account for
the discrepancy. What the data support is narrower than a single root
cause: the estimation \emph{theory} is sound (simulation confirms the
predicted noise-averaging benefit), and the hardware gap reflects some
combination of calibration, network, and real-time execution factors
that this campaign was not designed to separate. Resolving this would
require a dedicated follow-up study, isolating each candidate factor in
turn, which was outside the scope of the original campaign and remains
outside the scope of this dissertation.

\section{Scope Note}

The scalar-field Hessian and Newton-saddle-finder material found in the
earliest draft of Chapter~\ref{ch:vector_field}'s source paper
(\textit{scalar\_newton\_saddle\_finder}, and related functions in
\textit{src/control/quad\_primitives.py}) belongs to a separate,
abandoned research thread on scalar-field navigation and is not part of
this appendix's scope, which is limited to the four-robot generalization
of the vector-field estimator.

\end{document}
